% PDF/A
\begin{filecontents*}[overwrite]{\jobname.xmpdata}
\Title{Rozšíření nástroje ORMorpher}
\Author{Miroslav Štochel}
\Language{cs-CZ}
\Subject{Report of research project, Department of Software Engineering, Faculty of Mathematics and Physics, Charles University}
\Keywords{objektově-relační mapování\sep překlad dotazů\sep mezireprezentace\sep ORMorpher}
\end{filecontents*}
\documentclass{article}

% Mezery v \path{...}, musí být před pdfx
\PassOptionsToPackage{obeyspaces,spaces}{url}
\usepackage[a-2u]{pdfx}
\usepackage[utf8]{inputenc}
\usepackage[T1]{fontenc}
\usepackage{lmodern}
\usepackage[czech]{babel}
\usepackage{array}
\usepackage{tabularx}
\usepackage{booktabs}
\usepackage{longtable}
\usepackage[ddmmyyyy]{datetime}
\usepackage{geometry}
\usepackage{enumitem}
\usepackage{url}
\usepackage[section]{placeins}
\usepackage{float}
\emergencystretch=1.5em
\usepackage{xcolor}
\usepackage{amssymb}
\usepackage{listings}
\usepackage{upquote}
\usepackage{graphicx}
\usepackage{tikz}
\usetikzlibrary{positioning,fit,arrows.meta}

% --- Schémata ---
\tikzset{
  sch/.style={font=\linespread{1}\fontsize{8}{9.6}\selectfont\sffamily, align=center},
  box/.style={sch, draw=black!55, rounded corners=2pt, fill=gray!10, inner sep=4pt, minimum height=9mm},
  hl/.style={box, fill=blue!12},
  ext/.style={box, fill=white, dashed},
  lbl/.style={font=\linespread{1}\fontsize{7}{8.4}\selectfont\sffamily, align=center, text=black!80, inner sep=2pt},
  arr/.style={-{Latex[length=1.8mm]}, draw=black!70, line width=0.5pt},
  grp/.style={draw=black!40, dashed, rounded corners=3pt, inner sep=7pt},
}

% --- Screenshot ---
% Soubory PNG leží ve složce obrazky/ vedle tohoto souboru. Dokud soubor chybí,
% vysází se na jeho místě rámeček se jménem souboru, takže se dokument přeloží i bez screenshotů.
% \snimek[podíl šířky textu]{jméno souboru bez přípony}{popisek}{label}
\newcommand{\snimek}[4][1]{%
  \begin{figure}[htbp]
  \centering
  \setlength{\fboxsep}{0pt}\setlength{\fboxrule}{0.4pt}%
  \IfFileExists{obrazky/#2.png}%
    {\fcolorbox{black!40}{white}{\includegraphics[width=\dimexpr#1\textwidth-2\fboxrule\relax]{obrazky/#2.png}}}%
    {\fcolorbox{black!40}{gray!8}{\parbox[c][4.5cm][c]{\dimexpr#1\textwidth-2\fboxrule\relax}{\centering\sffamily\small Zde bude screenshot\\[2pt]\texttt{obrazky/\detokenize{#2}.png}}}}%
  \caption{#3}\label{#4}
  \end{figure}}

% Screenshot s číslovanými body, které popisek vysvětluje.
% \bod{číslo}{x}{y}: x a y jsou podíly šířky a výšky obrázku od jeho levého horního rohu.
% \snimekbody[podíl šířky textu]{jméno souboru bez přípony}{popisek}{label}{body}
\newcommand{\bod}[3]{\node[circle, fill=red!75!black, text=white, inner sep=0pt, minimum size=4.4mm,
  font=\linespread{1}\fontsize{8}{8}\selectfont\sffamily\bfseries] at (#2,{1-(#3)}) {#1};}
\newcommand{\snimekbody}[5][1]{%
  \begin{figure}[htbp]
  \centering
  \begin{tikzpicture}
  \node[anchor=south west, inner sep=0pt, draw=black!40, line width=0.4pt] (img) at (0,0)
    {\includegraphics[width=\dimexpr#1\textwidth-0.8pt\relax]{obrazky/#2.png}};
  \begin{scope}[x={(img.south east)}, y={(img.north west)}]
  #5
  \end{scope}
  \end{tikzpicture}
  \caption{#3}\label{#4}
  \end{figure}}

% --- Výpisy kódu ---
\definecolor{lstbg}{gray}{0.96}
\lstset{
  basicstyle=\linespread{1}\fontsize{7.5}{9}\selectfont\ttfamily,
  backgroundcolor=\color{lstbg},
  frame=single, framerule=0.2pt, rulecolor=\color{black!30},
  breaklines=true, breakatwhitespace=false,
  columns=fullflexible, keepspaces=true,
  showstringspaces=false, tabsize=4, upquote=true,
  xleftmargin=2pt, xrightmargin=2pt,
  aboveskip=4pt, belowskip=6pt,
  captionpos=b,
  literate={á}{{\'a}}1 {č}{{\v{c}}}1 {ď}{{\v{d}}}1 {é}{{\'e}}1 {ě}{{\v{e}}}1 {í}{{\'i}}1 {ň}{{\v{n}}}1 {ó}{{\'o}}1 {ř}{{\v{r}}}1 {š}{{\v{s}}}1 {ť}{{\v{t}}}1 {ú}{{\'u}}1 {ů}{{\r{u}}}1 {ý}{{\'y}}1 {ž}{{\v{z}}}1 {Á}{{\'A}}1 {Č}{{\v{C}}}1 {É}{{\'E}}1 {Ě}{{\v{E}}}1 {Í}{{\'I}}1 {Ř}{{\v{R}}}1 {Š}{{\v{S}}}1 {Ú}{{\'U}}1 {Ů}{{\r{U}}}1 {Ý}{{\'Y}}1 {Ž}{{\v{Z}}}1 {→}{{$\rightarrow$}}1 {×}{{$\times$}}1 {…}{{\dots}}1 {–}{{--}}1 {—}{{---}}1 {≥}{{$\geq$}}1 {✓}{{$\surd$}}1,
}
\renewcommand{\lstlistingname}{Výpis}

\renewcommand{\dateseparator}{.}
\renewcommand{\baselinestretch}{1.5}
\geometry{margin=2cm}
\setlength{\parskip}{3pt}
% Osamocený první nebo poslední řádek odstavce na straně
\widowpenalty=10000
\clubpenalty=10000

\renewcommand{\arraystretch}{1.3}
\newcolumntype{L}[1]{>{\raggedright\arraybackslash}p{#1}}
\newcolumntype{Y}{>{\raggedright\arraybackslash}X}
% Sloupec jako podíl šířky textu včetně svých okrajů: podíly jedné tabulky dávají dohromady 1
% a tabulka je pak přesně \linewidth široká. Pro longtable, kde X není.
\newcolumntype{P}[1]{>{\raggedright\arraybackslash}p{\dimexpr#1\linewidth-2\tabcolsep\relax}}
% Popisek longtable stojí pod tabulkou (v \endlastfoot) a sází se jako popisek plovoucí tabulky:
% v základní velikosti písma, na šířku textu (\LTcapwidth se nastavuje za \begin{document})
% a se stejným odstupem od spodní linky. Je to definice z longtable.sty s \normalsize,
% podpěrou nad prvním řádkem a mezerou pod popiskem místo \vskip\baselineskip.
% Každá longtable má před \endfoot prázdný řádek vysoký jako \endlastfoot (linka a popisek): strana,
% na které tabulka pokračuje, tak popisku drží místo a popisek nikdy nevyjde sám na další stranu.
\makeatletter
\def\LT@makecaption#1#2#3{%
  \LT@mcol\LT@cols c{\hbox to\z@{\hss\parbox[t]\LTcapwidth{\normalsize
    \rule{\z@}{17.8pt}%
    \sbox\@tempboxa{#1{#2: }#3}%
    \ifdim\wd\@tempboxa>\hsize
      #1{#2: }#3%
    \else
      \hbox to\hsize{\hfil\box\@tempboxa\hfil}%
    \fi
    \endgraf\vskip 6pt}%
  \hss}}}
% Obsah: místo pro dvouciferné číslo podkapitoly (3.10)
\renewcommand*\l@subsection{\@dottedtocline{2}{1.5em}{2.8em}}
\makeatother

\def\ResearchProject{Research project}

% Odevzdávaný release: tag 2.4.2 z 23. 11. 2026
\def\Vydani{2.4.2}

% METADATA SECTION START
\def\StudentsFullName{Miroslav Štochel}
\def\StudyProgram{Informatika -- Softwarové a datové inženýrství}
\def\ProjectTitle{Rozšíření nástroje ORMorpher}
\def\ProjectType{\ResearchProject}
\def\SupervisorFullName{Ing. Pavel Koupil, Ph.D.}
\def\ConsultantsFullName{doc. RNDr. Irena Holubová, Ph.D., RNDr. Michal Kopecký, Ph.D.}
% METADATA SECTION END

\title{\ProjectTitle}
\author{\StudentsFullName}

\begin{document}
\shorthandoff{-}
\hyphenation{Wide-World-Importers Primary-Key}
\setlength{\LTcapwidth}{\textwidth}

\centerline{\Large \textbf{Report of research project}}
\centerline{ \textbf{Department of Software Engineering}}
\centerline{ \textbf{Faculty of Mathematics and Physics, Charles University}}
\bigskip
{\noindent\renewcommand{\arraystretch}{1}\begin{tabular}{ >{\raggedleft}m{4cm} l}
 \textbf{Solver:} & \StudentsFullName \\
 \textbf{Study program:} & \StudyProgram \\[1em]
 \textbf{Project title:} & \ProjectTitle \\
 \textbf{Project type:} & \ProjectType \\[1em]
 \textbf{Supervisor:} & \SupervisorFullName \\
 \textbf{Consultants:} & \ConsultantsFullName \\[1em]
 \textbf{Repository:} & \url{https://github.com/MiraStochel/orm-convertor}, release \texttt{\Vydani} \\
 \textbf{Date:} & 23.~listopadu~2026 \\
\end{tabular}}

\bigskip
\bigskip

\noindent\textbf{Předmluva} Zpráva uzavírá výzkumný projekt \emph{Rozšíření nástroje ORMorpher}. Výchozím stavem byl prototyp z~diplomové práce M.~Abraháma~\cite{abraham-repo}: překládal vybrané entity a mapování mezi frameworky Dapper, NHibernate a EF~Core a dotazy jen z~EF~Core do Dapperu. Cílem projektu bylo rozšířit prototyp o~složitější mapovací konstrukce, o~využití metadat z~databázového schématu a o~javovské frameworky Hibernate, EclipseLink a MyBatis~\cite{spec,zamer}. Release \texttt{\Vydani} překládá entity, mapování a dotazy mezi všemi šesti frameworky. Překlad ověřují dvě testovací sady, validační sada LDBC, katalogy dotazů FIBEN a SEDE a šest případových studií. Srovnání s~jazykovými modely proběhlo ve čtyřech bězích na vzorcích 13 a 43 buněk.

\tableofcontents
\clearpage

\section{Úvod}

\subsection{Motivace a smysl nástroje}

Pro migraci dat a evoluci databázového schématu existuje řada nástrojů. Podpora překladu aplikační datové vrstvy mezi objektově-relačními frameworky je omezená~\cite{ormorpher-jss}: existující řešení se zaměřují na migraci na úrovni databáze, přepis dotazů v~SQL nebo ladění výkonu v~rámci jednoho frameworku. Přechod na jiný framework kvůli výkonu, vyjadřovacím možnostem nebo platformě podle~\cite{ormorpher-jss} často vyžaduje rozsáhlý ruční přepis entit, mapování a dotazů a ztěžuje ověření, že nový dotaz vrací tytéž řádky jako původní. ORMorpher tento překlad provádí deterministicky přes mezireprezentaci (obr.~\ref{fig:pipeline}). Nástroj má dva účely:

\begin{itemize}[leftmargin=*,itemsep=1pt]
\item \textbf{pro vývojáře, který migruje aplikaci}: přeloží podporované konstrukce a u~každé nepřeložené vydá záznam s~důvodem,
\item \textbf{pro výzkum} \emph{deterministická baseline}: podmnožina překládaných konstrukcí je vymezená číslovanými rozhodnutími a artefakt nástroje i~jazykového modelu se hodnotí týmiž čtyřmi stupni ověření nad týmiž vstupy (kap.~\ref{sec:metodika}). Informace o~databázi určuje varianta protokolu: nástroj čte katalog, model dostane jeho DDL jen ve variantě \texttt{schema}. Čtyři zaznamenané běhy variantu \texttt{schema} neobsahují (kap.~\ref{sec:llm}).
\end{itemize}

\begin{figure}[H]
\centering
\begin{tikzpicture}[node distance=9mm and 6mm]
\node[box, text width=25mm, minimum height=13mm] (src) {artefakty zdroje\\třídy, mapování,\\dotazy};
\node[hl, text width=25mm, minimum height=13mm, right=of src] (parser) {parser zdrojového\\frameworku};
\node[hl, text width=25mm, minimum height=13mm, right=of parser] (ir) {mezireprezentace\\entity, mapování,\\dotazy};
\node[hl, text width=25mm, minimum height=13mm, right=of ir] (builder) {builder cílového\\frameworku};
\node[box, text width=25mm, minimum height=13mm, right=of builder] (dst) {artefakty cíle\\entity, mapování,\\dotazy};
\node[ext, text width=25mm, minimum height=13mm, below=of ir] (cat) {katalog databáze\\SQL Server 2022};
\node[box, text width=25mm, minimum height=13mm, below=of dst] (rec) {záznamy o~převodu,\\tabulka faktů,\\záznam běhu};
\draw[arr] (src) -- (parser);
\draw[arr] (parser) -- (ir);
\draw[arr] (ir) -- (builder);
\draw[arr] (builder) -- (dst);
\draw[arr] (cat) -- node[lbl, left]{doplnění faktů,\\která ve zdroji nejsou} (ir);
\draw[arr] (builder.south) |- (rec.west);
\end{tikzpicture}
\caption{Pipeline převodu.}
\label{fig:pipeline}
\end{figure}

\subsection{Vztah k~dalším projektům}

Projekt vznikl ve výzkumné skupině Adaptive Data Management (ADaM) na KSI MFF UK~\cite{adam-web}. Přímo navazuje na prototyp ORMorpher z~diplomové práce Milana Abraháma~\cite{abraham-repo} a na dva texty, které jeho návrh popisují~\cite{ormorpher-ase,ormorpher-jss}. Repozitář projektu je fork repozitáře prototypu: projekt z~něj převzal kód, mezireprezentaci a architekturu parser--builder a rozšířil je.

Ostatní práce skupiny jsou tematický kontext. Řeší příbuznou úlohu jinou metodou (jazykové modely) nebo na jiné úrovni (data a databázový systém) a projekt z~nich nepřebírá kód ani data. Návrh s~nimi souvisí nepřímo, přes prototyp: podle~\cite{ormorpher-ase} vychází mezireprezentace entit a mapování z~kategoriálního modelu~\cite{koupil-holubova-2022} a instrukční reprezentace dotazů z~přístupu k~dotazování nad více modely~\cite{koupil-crha-holubova-2025}. Obě skupiny prací uvádí tab.~\ref{tab:vztahy}.

\begin{table}[H]
\centering\footnotesize
\begin{tabularx}{\textwidth}{L{3.1cm} L{6.9cm} Y}
\toprule
\textbf{Práce} & \textbf{Co obsahuje} & \textbf{Vztah k~projektu} \\
\midrule
\multicolumn{3}{@{}l}{\textbf{Přímí předchůdci}} \\
diplomová práce M.~Abraháma a repozitář prototypu~\cite{abraham-repo} & překlad entit a mapování mezi frameworky Dapper, NHibernate a EF~Core, překlad dotazů z~EF~Core do Dapperu, Advisor nad ILP pro Dapper a EF~Core & výchozí kód: repozitář projektu je jeho fork \\
demonstrační příspěvek na ASE 2025~\cite{ormorpher-ase} & popis prototypu: překlad entit a dotazů a výběr frameworku v~interaktivním rozhraní nad EF~Core, NHibernatem a Dapperem & popisuje stav, ze kterého projekt vychází \\
rozšířená časopisecká verze, preprint~\cite{ormorpher-jss} & návrh mezireprezentace, překladová pravidla E1--E10 a Q1--Q15, architektura parser--builder, model ILP & návrh, který projekt rozšiřuje. Rozhodnutí projektu se zdůvodňují vůči němu \\
\midrule
\multicolumn{3}{@{}l}{\textbf{Tematický kontext}} \\
diplomové práce M.~Čorovčáka~\cite{corovcak-dp} a I.~Oboňové~\cite{obonova-dp} & migrace aplikačního kódu mezi relačními, grafovými a dokumentovými systémy přes univerzální objektový model (Čorovčák) a migrace dotazů mezi relačními a grafovými systémy (Oboňová), obojí velkými jazykovými modely & příbuzná úloha jinou metodou. ORMorpher překládá pravidly a mezireprezentací, deterministicky, a jeho výsledky jsou baseline pro přístupy založené na LLM (kap.~\ref{sec:exp}) \\
Database Mapping Optimization~\cite{mmcat-aimm}, mm-mapsearch~\cite{mm-mapsearch} & volba reprezentace dat podle pracovní zátěže: měření latencí v~PostgreSQL, MongoDB a Neo4j, naučené modely latence, prohledávání prostoru mapování metodou Monte Carlo Tree Search & jiná úroveň: volí uložení dat, ORMorpher převádí kód, kterým k~nim aplikace přistupuje. Společné s~Advisorem prototypu je rozhodování podle měření \\
MM-cat~\cite{mmcat-intro,koupil-holubova-2022}, MM-quecat, MM-evoque, MM-evocat~\cite{mm-quecat,mm-evoque,mm-evocat} & jednotná reprezentace multi-modelových dat založená na teorii kategorií, dotazování nad ní, migrace mezi modely a evoluce schématu & jiná úroveň: data a databázový systém. ORMorpher pracuje s~entitami, mapováním a dotazy aplikace nad relační databází. Z~kategoriálního modelu~\cite{koupil-holubova-2022} vychází mezireprezentace prototypu~\cite{ormorpher-ase} \\
\bottomrule
\end{tabularx}%
\caption{Přímí předchůdci projektu a tematický kontext ve skupině ADaM.}
\label{tab:vztahy}
\end{table}

\subsection{Návaznost na specifikaci}

Specifikace~\cite{spec} rozdělila práci do tří cílů:

\begin{itemize}[leftmargin=*,itemsep=1pt]
\item \textbf{.NET frameworky} (F1--F6, F11, F14, S1--S3, S5--S7). Splněno už ve verzi 1.2.0, kterou specifikace popisuje, s~mezemi u~F11 a F14 (tab.~\ref{tab:prehled-pozadavku}). Později přibyly výrazy, CTE, rekurze, okenní funkce a fallback na nativní SQL.
\item \textbf{Javovské frameworky} (F7--F10, F12, F13). Splněno od verze 2.0.0 s~mezí u~F13: Hibernate~\cite{hibernate-docs} a EclipseLink~\cite{eclipselink-docs} jako dva profily nad sdílenou vrstvou Jakarta Persistence, MyBatis~\cite{mybatis-docs} nad vrstvami jazyka Java a T-SQL, javovská testovací sada s~diferenčním ověřením a překlad mezi ekosystémy oběma směry.
\item \textbf{Advisor, benchmarking a experimenty} (T1--T7, F15, S4). Splněno částečně. Implementován je vstup případové studie, matice kategorií a metriky (T1--T3). Srovnání s~jazykovými modely (T4--T6) má protokol a čtyři běhy na vzorcích, všechny ve variantách T4 (\texttt{zero-shot} a \texttt{few-shot}). Plné experimenty patří navazující diplomové práci. Rozšíření Advisoru (F15), jeho porovnání (T7) a izolace spouštění měřeného kódu (S4) realizovány nebyly (rozh.~171).
\end{itemize}

Stav jednotlivých požadavků uvádí tab.~\ref{tab:prehled-pozadavku}. U~stavu \emph{splněno částečně} sloupec \emph{Rozsah a meze} uvádí, co je implementováno a co je mimo rozsah. Stav \emph{připraveno} znamená metodiku a běhy na vzorcích, ne provedený experiment. Specifikace u~T4--T6 žádá jen tuto přípravu a u~T1--T3 vstup, matici a měřicí infrastrukturu bez závazných hodnot. Číselná kritéria, která specifikace neobsahuje (mez 30~s u~S3, počty testů u~F12, rozsah případové studie u~T1 a silný model u~T4), pocházejí ze seznamu požadavků vedoucího (\path{docs/requirements.md}), jehož číslování specifikace přebírá.

{\footnotesize
\begin{longtable}{P{0.30} P{0.13} P{0.37} P{0.20}}
\toprule
\textbf{Požadavek} & \textbf{Stav} & \textbf{Rozsah a meze} & \textbf{Kde ve zprávě} \\
\midrule
\endfirsthead
\toprule
\textbf{Požadavek} & \textbf{Stav} & \textbf{Rozsah a meze} & \textbf{Kde ve zprávě} \\
\midrule
\endhead
\rule{0pt}{26pt}\\
\endfoot
\bottomrule
\caption{Požadavky specifikace v~release \Vydani.}\label{tab:prehled-pozadavku}\\
\endlastfoot
\textbf{F1--F3} Složené identifikátory, jejich překlad, cizí klíče na ně, vztahy 1:1, 1:N a N:M & splněno &-- & kap.~\ref{sec:co-jak}, příklady 2 a~3 \\
\textbf{F4--F6} Metadata ze schématu databáze, jejich sloučení se zdrojem, úplné mapování z~neúplného vstupu & splněno &-- & kap.~\ref{sec:co-jak}, příklad~1 \\
\textbf{F7--F9} Hibernate, MyBatis, EclipseLink & splněno & tři frameworky, ne ekosystém. Dynamicky parametrizované dotazy (F8) se překládají s~parametry a se značkou \texttt{<foreach>}. Dynamické SQL MyBatisu (\texttt{<if>}, \texttt{<choose>}, \texttt{<bind>}) se odmítá záznamem a lazy loading reference v~EclipseLinku není zaručen & kap.~\ref{sec:limity} \\
\textbf{F10} Překlad přes hranici ekosystémů & splněno & v~každé ze čtyř uspořádaných dvojic ekosystémů & příklady 2 a~3, tab.~\ref{tab:t2t3} \\
\textbf{F11} Validace cílových artefaktů & splněno částečně & implementováno a testováno: kontrola úplnosti mezireprezentace a omezení cíle před generováním, strukturované záznamy, odmítnutí dotazu, který by vrátil jiné řádky. Mimo rozsah: kontrola syntaxe generovaných souborů při převodu (\texttt{/convert}). Syntaxi ověřují testovací sady & kap.~\ref{sec:metodika} \\
\textbf{F12} Javovská testovací sada & splněno & build sady selže, má-li méně testů, než požadavek žádá & kap.~\ref{sec:testy} \\
\textbf{F13} Diferenční ověření výsledků dotazů & splněno částečně & implementováno a testováno v~obou sadách. Mimo rozsah: nastavitelná přesnost porovnání a nastavitelné zacházení s~hodnotou \texttt{NULL}. Počet desetinných míst a platných číslic je určen u~každého dotazu a \texttt{NULL} má jeden pevný zápis, aby se shodly sady ve dvou jazycích & tab.~\ref{tab:zajimavosti}, kap.~\ref{sec:metodika} \\
\textbf{F14} Celé třídy a dávkové vstupy v~rozhraní & splněno částečně & implementováno: více jednotek v~jednom převodu, archiv ZIP, výstup po souborech. Mimo rozsah: zobrazení mezireprezentace (rozh.~154), procházení repozitáře a jiné archivy než ZIP & kap.~\ref{sec:vlastni}, \ref{sec:odchylky} \\
\textbf{F15} Výběr cílového frameworku (Advisor) & nerealizováno & specifikace rozšíření podmínila časovými možnostmi. Advisor prototypu měří jen Dapper a EF~Core a je vyňatý ze záruk & tab.~\ref{tab:vynate}, \ref{tab:odchylky} \\
\midrule
\textbf{S1--S3} Modularita, determinismus, výkon překladu & splněno & překlad projektu trvá 0,2--0,7~s, mez je 30~s & tab.~\ref{tab:projekty}, \ref{tab:vykon}, kap.~\ref{sec:vstup-vystup} \\
\textbf{S4} Izolace a bezpečnost spouštění & splněno částečně & implementováno a testováno: přihlašovací údaje nejsou v~artefaktech, v~záznamech ani v~logu. Repozitář obsahuje jen údaje serverů, které si compose a CI vytvoří pro jeden běh. Mimo rozsah: izolace spouštění měřeného kódu, odložená s~F15 & kap.~\ref{sec:konfigurace} \\
\textbf{S5} Reprodukovatelné prostředí & splněno & experimentální běh běží v~týchž kontejnerech na vyžádání, build ho nespouští & kap.~\ref{sec:testy}, \ref{sec:reprodukce} \\
\textbf{S6} Pozorovatelnost a auditovatelnost & splněno & výsledek převodu obsahuje identifikátor běhu, zdroje metadat, záznamy a verze nástroje a frameworků & kap.~\ref{sec:vstup-vystup} \\
\textbf{S7} Uživatelská přívětivost & splněno & implementovány jsou funkce, které požadavek jmenuje: validace před spuštěním, stav zpracování, označení chyby v~souboru a na řádku, stažení výstupu a základní scénář ve čtyřech krocích (mez je pět). Test použitelnosti s~uživateli proveden nebyl & kap.~\ref{sec:tutorial}, \ref{sec:vlastni} \\
\midrule
\textbf{T1--T3} Případová studie, matice překladů, metriky korektnosti & splněno částečně & implementováno: vstup případové studie, matice kategorií a metriky experimentálního běhu. Mimo rozsah: výsledek jako závazek. Čísla jsou záznam jednoho běhu (rozh.~183, 184) & kap.~\ref{sec:vysledky} \\
\textbf{T4--T6} Příprava LLM baseline, RAG a agentní varianty, ablace & připraveno & protokol a čtyři kanály (T4--T6), čtyři běhy na vzorcích (jen varianty T4). Plné experimenty patří navazující diplomové práci & kap.~\ref{sec:llm} \\
\textbf{T7} Porovnání Advisoru & nerealizováno & odloženo s~F15 & kap.~\ref{sec:odchylky} \\
\end{longtable}}

Repozitář tytéž stavy označuje jako \emph{nárokovaný}, \emph{nárokovaný v~užším rozsahu} a \emph{vyňatý} (\path{docs/traceability.md}, anglicky kořenové \texttt{README.md}, \emph{Guarantees}). Stav \emph{vyňatý} tam mají řádky \emph{připraveno} a \emph{nerealizováno} a první věta požadavku S4. Požadavek S6 tam má užší rozsah, protože seznam požadavků vedoucího žádá v~záznamu převodu i~výsledky kompilace a testů. Ty jsou v~záznamech běhů sad.

Zpráva popisuje release \texttt{\Vydani} (tag z~23.~listopadu~2026): nasazení, příklady i~počty testů (tab.~\ref{tab:cisla}) odpovídají tomuto tagu. Výsledky běhů na vyžádání v~kap.~\ref{sec:exp} byly naměřeny nad staršími commity, které uvádí tab.~\ref{tab:zaznamy-behu}. Konec projektu podle schváleného záměru~\cite{zamer} je 1.~prosince~2026.

\subsection{Struktura dokumentu}

Kapitola~\ref{sec:nasazeni} popisuje nasazení a provoz. Kapitola~\ref{sec:pouziti} je uživatelská dokumentace s~příklady. Kapitola~\ref{sec:denik} popisuje vybraná rozhodnutí návrhu, vybraná selhání s~nápravou, odchylky od specifikace a využití AI. Kapitola~\ref{sec:exp} obsahuje metodiku experimentů, výsledky nástroje, srovnání s~jazykovými modely a postup reprodukce. Kapitola~\ref{sec:zaver} shrnuje plnění požadavků a navazující práci.

\section{Nasazení a běh}\label{sec:nasazeni}

Kapitola popisuje nasazení release \texttt{\Vydani}: obsah repozitáře, spuštění systému, konfiguraci, testovací sady, běhy na vyžádání a řešení problémů. Postupy pro spuštění systému a obou testovacích sad jsou zde úplné.

\subsection{Obsah repozitáře}\label{sec:repozitar}

Repozitář projektu je na adrese \url{https://github.com/MiraStochel/orm-convertor}~\cite{fork-repo}. Jeho složky ukazuje výpis~\ref{lst:strom}. Nástroj je ve složce \texttt{ORMConvertor/}.

\noindent\begin{minipage}{\linewidth}
\begin{lstlisting}[caption={Složky repozitáře.},label={lst:strom}]
orm-convertor/
|-- ORMConvertor/               nástroj
|   |-- ORMConvertor.sln        solution pro .NET 10
|   |-- Model/, Tests/, ...     projekty solution
|   |-- JavaTests/              javovská testovací sada (Maven, mimo .sln)
|   |-- FibenImport/            importní projekt (mimo .sln)
|   |-- SedeImport/             importní projekt (mimo .sln)
|   |-- database/               skripty databází
|   |-- database.Dockerfile     kontejnerová konfigurace: image databáze
|   \-- docker-compose.yml      kontejnerová konfigurace: služby
|-- docs/                       dokumentace
|-- benchmarks/                 zděděná složka
|-- diagrams/                   zděděná složka
\-- notes/                      zděděná složka
\end{lstlisting}
\end{minipage}

Převod prochází pěti moduly. Projekt \path{ORMConvertorAPI} přijme požadavek REST a předá ho orchestraci \path{OrmConvertor}. Orchestrace podle zdrojového a cílového frameworku vybere parser a builder z~wrapperů. Parser převede zdrojové artefakty do mezireprezentace (\texttt{Model}), projekt \path{DatabaseCatalog} do ní doplní fakta z~katalogu databáze a builder z~ní zapíše artefakty cíle. Orchestrace vrátí artefakty a diagnostické záznamy. K~wrapperům přistupuje jen přes rozhraní z~\path{AbstractWrappers}, takže nový framework je nový projekt wrapperu (S1). Všechny projekty solution uvádí tab.~\ref{tab:projekty}. Podrobnosti, které zpráva neopakuje, jsou v~dokumentaci repozitáře (tab.~\ref{tab:dokumentace}).

\begin{table}[H]
\centering\footnotesize
\begin{tabularx}{\textwidth}{L{7.3cm} Y}
\toprule
\textbf{Projekty} & \textbf{Odpovědnost} \\
\midrule
\texttt{Model}, \texttt{AbstractWrappers}, \texttt{Common} & mezireprezentace, rozhraní parserů a builderů s~deskriptory cílů a typy záznamů, sdílené převodníky \\
\texttt{DapperWrappers}, \texttt{EFCoreWrappers}, \texttt{NHibernateWrappers} & parser a builder každého .NET frameworku \\
\texttt{CSharpEntityParsing}, \texttt{LinqParsing}, \texttt{LinqBuilding} & sdílené čtení C\# entit, čtení a zápis LINQ \\
\texttt{JavaEntityParsing}, \texttt{JakartaPersistence} & lexer a parser podmnožiny Javy bez JVM, sdílená vrstva JPA (anotace, \texttt{orm.xml}, JPQL) \\
\texttt{HibernateWrappers}, \texttt{EclipseLinkWrappers}, \texttt{MyBatisWrappers} & dva profily nad vrstvou JPA a mapper nad vrstvami jazyka Java a T-SQL \\
\texttt{TransactSql}, \texttt{DatabaseCatalog} & čtení a zápis T-SQL nad ScriptDom, čtení metadat databáze \\
\texttt{OrmConvertor}, \texttt{ORMConvertorAPI} & orchestrace pipeline, REST API a statický frontend \\
\texttt{Advisor}, \texttt{AdvisorBenchmarking} & ILP model nad GLPK, měření generovaného kódu \\
\texttt{SampleData}, \texttt{Tests} & vzorky, příklady a katalogy LDBC, FIBEN a SEDE, sada xUnit~v3 s~experimentálními běhy \\
\bottomrule
\end{tabularx}%
\caption{Projekty solution (podrobně \texttt{docs/architecture.md}, §2).}
\label{tab:projekty}
\end{table}

\begin{table}[H]
\centering\footnotesize
\begin{tabular}{l l}
\toprule
\textbf{Co} & \textbf{Kde v~repozitáři} \\
\midrule
současné chování nástroje & \path{docs/architecture.md} \\
co se překládá a co ne, konstrukci po konstrukci & \texttt{docs/subset.md} \\
číslovaná rozhodnutí & \texttt{docs/decisions/} \\
záznamy běhů sad a měření & \texttt{docs/runs/} \\
hranice záruk & kořenové \texttt{README.md}, \emph{Guarantees} \\
provozní příručka (anglicky) & \path{ORMConvertor/README.md} \\
na čem lze pokračovat & \path{docs/open-items.md} \\
\bottomrule
\end{tabular}
\caption{Dokumentace v~repozitáři.}
\label{tab:dokumentace}
\end{table}

\subsection{Spuštění systému}\label{sec:spusteni}

\textbf{Předpoklady.} Hostitel architektury x86-64 (\texttt{amd64}), Docker s~Compose a BuildKit, virtuální stroj Dockeru s~alespoň čtyřmi procesory, síť při buildu a volné porty 5072 a 1444. .NET SDK, JDK ani databáze na hostiteli nejsou potřeba.

\noindent\begin{minipage}{\linewidth}
\begin{lstlisting}[language=bash,caption={Spuštění systému a ověření.}]
git clone --branch 2.4.2 https://github.com/MiraStochel/orm-convertor.git
cd orm-convertor/ORMConvertor
docker compose up --build
# ověření z druhého terminálu
docker compose logs mssql_db | grep holds          # LdbcSnb holds ..., Fiben holds ..., Sede holds ...
curl http://localhost:5072/orm/required-content    # odpověď 200
\end{lstlisting}
\end{minipage}

\textbf{Očekávaný stav.} Příkaz spustí dvě služby (obr.~\ref{fig:nasazeni}): aplikaci na \url{http://localhost:5072/orm/} a databázi SQL Server 2022 na portu 1444 hostitele (\texttt{sa}/\texttt{Testingorms123}, jen pro vývoj, data na volume \texttt{mssql\_data}). Úspěch má tři znaky:

\begin{itemize}[leftmargin=*,itemsep=1pt]
\item log služby \texttt{mssql\_db} obsahuje řádky \emph{LdbcSnb holds \ldots}, \emph{Fiben holds \ldots} a \emph{Sede holds \ldots} a žádný řádek \emph{\ldots loader stopped before it finished},
\item \path{GET /orm/required-content} vrátí odpověď 200 se seznamem jazyků, které každý zdrojový framework čte,
\item první převod z~kap.~\ref{sec:tutorial} má stav katalogu \emph{catalog read}.
\end{itemize}

Aplikace se spustí, jakmile server přijímá spojení, a na načtení datových sad nečeká. První start načte \texttt{LdbcSnb} asi za 5~minut, \texttt{Fiben} za 2 až 3~minuty a \texttt{Sede} za několik minut. Další starty už nenačítají nic. První build stahuje base image (asi 1~GB) a archivy datových sad (asi 0,9~GB) a na lince 50~Mb/s trvá asi deset minut. Velikosti a časy jsou orientační a závisejí na lince a na stroji. Když stažení datové sady selže, build skončí hlášením, které jmenuje přepínač pro build bez této sady (tab.~\ref{tab:sady}, tab.~\ref{tab:problemy}). \texttt{docker compose stop} služby zastaví a data ponechá, \texttt{docker compose down -v} odstraní i~volume s~databázemi.

\textbf{Bezpečnostní omezení.} Instance naslouchá bez TLS a žádný endpoint nevyžaduje přihlášení. Je určena pouze pro lokální a vývojové prostředí nebo pro důvěryhodnou síť. Veřejné nasazení vyžaduje reverzní proxy, která ukončuje HTTPS a řídí přístup, a zúžené nastavení \texttt{AllowedHosts}. Kroky uvádí \path{ORMConvertor/README.md}, \emph{Deploying a real instance}, předpoklady \path{docs/threat-model.md}. Konfiguraci konkrétní proxy repozitář neobsahuje.

\begin{figure}[H]
\centering
\begin{tikzpicture}[node distance=8mm and 12mm]
\node[hl, text width=27mm, minimum height=15mm] (app) {\texttt{ormconvertor}\\REST API a frontend\\port 5072};
\node[box, text width=27mm, minimum height=15mm, right=of app] (db) {\texttt{mssql\_db}\\SQL Server 2022\\port 1444};
\node[lbl, below=1mm of db] (dbs) {WideWorldImporters, LdbcSnb,\\Fiben, Sede};
\draw[arr] (app) -- node[lbl, above]{katalog} (db);
\node[grp, fit=(app)(db)(dbs), label={[lbl]above:bez profilu (systém)}] (sys) {};

\node[box, text width=27mm, minimum height=15mm, right=20mm of db] (tests) {\texttt{tests}\\.NET sada};
\node[box, text width=27mm, minimum height=15mm, right=of tests] (tdb) {\texttt{test\_db}\\SQL Server 2022\\LDBC scale factor 0,1};
\node[box, text width=27mm, minimum height=15mm, below=of tests] (jt) {\texttt{java\_tests}\\javovská sada};
\node[hl, text width=27mm, minimum height=15mm, below=of tdb] (tapp) {\texttt{test\_app}, \texttt{\_ldbc},\\\texttt{\_cases}, \texttt{\_dry}\\čtyři instance nástroje};
\draw[arr] (tests) -- (tdb);
\draw[arr] (jt) -- node[lbl, below]{HTTP} (tapp);
\draw[arr] (tapp) -- node[lbl, right]{katalog\\(ne \texttt{\_dry})} (tdb);
\draw[arr] (jt.north east) -- node[lbl, above, sloped]{JDBC} (tdb.south west);
\node[grp, fit=(tests)(tdb)(jt)(tapp), label={[lbl]above:profil \texttt{test}}] (tst) {};
\end{tikzpicture}
\caption{Služby souboru \texttt{docker-compose.yml}: systém a profil \texttt{test}.}
\label{fig:nasazeni}
\end{figure}

\subsection{Běh bez Dockeru}\label{sec:bez-dockeru}

Alternativně lze aplikaci nasadit bez Dockeru. Překlad běží všude, kde je .NET SDK verze 10.0.100 nebo novější verze 10.0 (\texttt{global.json}). Příkazy se spouštějí z~kořene klonu:

\noindent\begin{minipage}{\linewidth}
\begin{lstlisting}[language=bash,caption={Spuštění z~.NET CLI a ověření.}]
cd ORMConvertor
# volitelně: katalog databáze
dotnet user-secrets set "ConnectionStrings:CatalogDatabase" "<connection string>" --project ORMConvertorAPI
dotnet run --configuration Release --launch-profile http --project ORMConvertorAPI/ORMConvertorAPI.csproj
# ověření z druhého terminálu
curl http://localhost:5072/orm/required-content    # odpověď 200
\end{lstlisting}
\end{minipage}

Aplikace běží na \url{http://localhost:5072/orm/}. Bezpečnostní omezení z~kap.~\ref{sec:spusteni} platí i~zde. Proti běhu v~Dockeru se liší třemi věcmi:

\begin{itemize}[leftmargin=*,itemsep=1pt]
\item Bez klíče katalogu nástroj překládá jen s~konvencemi a převod má stav \emph{catalog not configured}. Požadavky F4 a F6 pak přes rozhraní neplatí.
\item Advisor nefunguje, protože \texttt{libadvisor.so} vzniká jen v~image. Překladu se to netýká.
\item Publikovaná instance (\texttt{dotnet publish}) se spouští z~adresáře publikace nebo s~\texttt{-{}-contentRoot}, jinak \texttt{/orm/} vrací prázdnou odpověď 404. Běží v~prostředí \texttt{Production} a konfiguraci čte z~proměnných prostředí.
\end{itemize}

\subsection{Konfigurace}\label{sec:konfigurace}

Systém spuštěný příkazem \texttt{docker compose up -{}-build} nastavení nepotřebuje: compose předá aplikaci oba connection stringy a profil \texttt{test} předá nastavení oběma testovacím sadám. Nastavuje se jen to, čím se běh od výchozího liší: katalog pro instanci mimo Docker, jiná databáze katalogu v~compose, databáze pro testovací sady spouštěné na hostiteli a limity vstupu (tab.~\ref{tab:konfigurace}). Výpis~\ref{lst:konfigurace} nastavuje katalog třemi způsoby.

\noindent\begin{minipage}{\linewidth}
\begin{lstlisting}[language=bash,caption={Nastavení katalogu databáze.},label={lst:konfigurace}]
CS="Server=localhost,1444;Database=WideWorldImporters;User ID=sa;Password=Testingorms123;TrustServerCertificate=true"
# instance z .NET CLI (prostředí Development): user secrets
dotnet user-secrets set "ConnectionStrings:CatalogDatabase" "$CS" --project ORMConvertorAPI
# publikovaná instance (prostředí Production): proměnná prostředí
export ConnectionStrings__CatalogDatabase="$CS"
# Docker: jiná databáze katalogu na témže serveru
ORMCONVERTOR_CATALOG_DATABASE=LdbcSnb docker compose up -d --build ormconvertor
\end{lstlisting}
\end{minipage}

Klíče se píší s~\texttt{:} v~user secrets a v~\texttt{appsettings.json} a s~\texttt{\_\_} v~proměnných prostředí. User secrets čte aplikace jen v~prostředí \texttt{Development}. Zda aplikace klíč katalogu načetla, ukáže stav katalogu u~převodu: rozhraní zobrazí \emph{catalog read} s~dobou čtení a odpověď \path{POST /orm/convert} má v~poli \texttt{catalogState} hodnotu 2. Hodnota 0 znamená, že klíč nastaven není, hodnota 1, že převod katalog nepotřeboval, a hodnota 3, že databáze neodpověděla. Žádný connection string nemá v~repozitáři hodnotu kromě přihlašovacích údajů serverů, které si compose a CI vytvoří pro jeden běh (S4). Úplný výčet klíčů je v~\path{ORMConvertor/README.md}, \emph{Configuration}.

\begin{table}[H]
\centering\footnotesize
\begin{tabularx}{\textwidth}{L{5cm} L{1.5cm} Y}
\toprule
\textbf{Klíč} & \textbf{Kdo} & \textbf{Význam a chování bez klíče} \\
\midrule
\path{ConnectionStrings:CatalogDatabase} & API & doplnění z~katalogu. Bez něj jen konvence, F4 a F6 přes rozhraní neplatí \\
\path{ConnectionStrings:AdvisorDatabase} & API & měření Advisoru. Bez něj hlášení, které klíč jmenuje \\
\path{ConnectionStrings:TestDatabase} & .NET sada & databázové testy. Bez něj se přeskočí s~důvodem \\
\path{ORMCONVERTOR_TEST_JDBC_URL}, \path{ORMCONVERTOR_API_URL} & Java sada & databáze a běžící instance nástroje, z~níž sada bere artefakty (rozh.~078) \\
\path{Parsing:MaxNestingDepth}, \path{Parsing:MaxQueryTextLength} & API & limit hloubky zanoření (128) a délky textu dotazu (1~048~576 znaků). Vstup nad limitem je \emph{Failure} \\
\path{ORMCONVERTOR_CATALOG_DATABASE} & compose & databáze, nad kterou instance čte katalog. Výchozí je \texttt{WideWorldImporters} \\
\bottomrule
\end{tabularx}%
\caption{Klíče konfigurace provozu.}
\label{tab:konfigurace}
\end{table}

\subsection{Build image a datové sady}\label{sec:build}

Tato část popisuje build infrastrukturu a k~běžnému spuštění není potřeba. Image vznikají ze dvou souborů \texttt{Dockerfile} a služby popisuje jeden \texttt{docker-compose.yml} se dvěma částmi: systém bez profilu a profil \texttt{test} (rozh.~039). \path{ORMConvertorAPI/Dockerfile} má pět stage: \texttt{advisor-native} (jediné místo, kde vzniká \texttt{libadvisor.so}), \texttt{dotnet-build}, \texttt{tests}, \texttt{java-tests} a \texttt{runtime}. \texttt{database.Dockerfile} staví image SQL Serveru 2022 se čtyřmi databázemi (tab.~\ref{tab:sady}). Base image jsou připnuté digestem a verze balíčků .NET commitnutým lock file (rozh.~186).

\begin{table}[H]
\centering\footnotesize
\begin{tabularx}{\textwidth}{L{2.9cm} Y L{5.0cm} L{3.4cm}}
\toprule
\textbf{Databáze} & \textbf{Sada} & \textbf{Odkud ji build bere} & \textbf{Přepínač, který ji vynechá} \\
\midrule
\texttt{WideWorldImporters} & vzorová databáze & release \path{Microsoft/sql-server-samples} & \texttt{WWI\_BACKUP=none} \\
\texttt{LdbcSnb} & LDBC SNB Interactive v1~\cite{ldbc}, scale factor 1, archiv dat 222~MB a validační sada 205~MB & \texttt{datasets.ldbcouncil.org} (rozh.~219) & \texttt{LDBC\_SCALE\_FACTOR=none} \\
\texttt{Fiben} & benchmark IBM Research~\cite{fiben}, archiv 80~MB & asset release \texttt{data\_sets} repozitáře, checksum v~\path{database/datasets.sha256} (rozh.~215) & \texttt{FIBEN\_DATASET=none} \\
\texttt{Sede} & dump dba.stackexchange.com pro dotazy SEDE~\cite{sede}, archiv 319~MB & archive.org, ověřeno MD5 & \texttt{SEDE\_DATASET=none} \\
\bottomrule
\end{tabularx}%
\caption{Databáze image \texttt{mssql\_db}. Vlastní kopii má projekt jen u~FIBEN, jehož původní repozitář už není dostupný. Ostatní sady zůstávají cizími zdroji buildu (rozh.~218). Překlad sám žádnou databázi nepotřebuje.}
\label{tab:sady}
\end{table}

\subsection{Testovací sady}\label{sec:testy}

Tato část slouží k~ověření vývoje a k~běžnému provozu není potřeba. Testovací sady jsou dvě, .NET a javovská. Každou reprodukuje jeden příkaz na stroji, kde je jen Docker (S5). Prostředí jim staví profil \texttt{test} (obr.~\ref{fig:nasazeni}, tab.~\ref{tab:testprofil}).

\textbf{Předpoklady.} Virtuální stroj Dockeru se čtyřmi procesory a 12~GB paměti bez swapu (rozh.~193), síť při buildu a zastavený systém (\path{docker compose stop ormconvertor mssql_db}). Systém a profil \texttt{test} se do 12~GB současně nevejdou.

\noindent\begin{minipage}{\linewidth}
\begin{lstlisting}[language=bash,caption={Obě sady v~kontejnerech.}]
docker compose --profile test build tests java_tests test_app test_app_ldbc test_app_cases test_app_dry test_db
docker compose --profile test run --rm tests         # .NET sada
docker compose --profile test run --rm java_tests    # javovská sada
\end{lstlisting}
\end{minipage}

\textbf{Kritérium úspěchu.} Žádný selhaný test v~obou sadách a počet testů shodný se záznamem velikosti sady pro release \texttt{\Vydani} (tab.~\ref{tab:cisla}). Přeskočený databázový test je v~profilu \texttt{test} selhání.

\textbf{Krok \texttt{build} není volitelný.} Příkaz \texttt{run} sám použije existující image a testuje stav kódu, z~něhož image vznikl. Sada přitom projde bez upozornění. Jediným signálem je počet testů jiný než v~záznamu.

\begin{table}[H]
\centering\footnotesize
\begin{tabularx}{\textwidth}{L{4cm} Y}
\toprule
\textbf{Služba} & \textbf{Role} \\
\midrule
\texttt{test\_db}, \texttt{test\_db\_init} & SQL Server jen s~daty LDBC scale factoru 0,1. \texttt{init} počká na načtení a vytvoří databáze sad \\
\texttt{tests} & .NET sada, mez paměti 6~GB. Schémata si vytváří a maže sama (rozh.~016) \\
\texttt{test\_app}, \texttt{\_ldbc}, \texttt{\_cases}, \texttt{\_dry} & čtyři instance nástroje (2~GB každá): katalog nad testovací databází, nad \texttt{LdbcSnb}, nad případovými studiemi a žádný \\
\texttt{java\_tests} & javovská sada, mez paměti 6~GB. Artefakty bere z~instancí přes HTTP \\
\bottomrule
\end{tabularx}%
\caption{Služby profilu \texttt{test}.}
\label{tab:testprofil}
\end{table}

Mimo kontejnery běží .NET sada příkazem \path{dotnet test Tests/Tests.csproj --configuration Release}. Javovská sada potřebuje JDK~25, Maven, databázi a běžící instanci s~katalogem nad testovací databází.

CI běží při změně \texttt{ORMConvertor/**} nebo workflow a má tři joby: \texttt{test} (.NET sada proti SQL Serveru 2022), \texttt{java-test} (Temurin 25, spuštěná aplikace, \texttt{mvn -B test}) a \texttt{dependencies} (známé zranitelnosti připnutých balíčků, běží také jednou týdně). Datové sady CI nenačítá, proto v~ní validátor LDBC a běhy katalogů přeskakují.

\subsection{Běhy na vyžádání}\label{sec:behy}

Běhy na vyžádání jsou případy obou testovacích sad zapnuté proměnnou prostředí (tab.~\ref{tab:behy}). Od testů se liší ve třech bodech:

\begin{itemize}[leftmargin=*,itemsep=1pt]
\item \textbf{Test} je případ testovací sady. Běží při buildu a v~CI a jeho výsledek je úspěch nebo selhání.
\item \textbf{Měření} (benchmark) je běh, který zaznamenává čas nebo paměť: běh FIBEN, běh SEDE, běh mezí a celý vstup. Vyžaduje vypnuté dynamické navýšení frekvence procesoru, běhy FIBEN a SEDE také měřicí režim serveru (tab.~\ref{tab:rezimy}). Krátký dotaz se opakuje a zaznamenává se medián.
\item \textbf{Dlouhý experimentální běh} trvá hodiny až dny: validátor LDBC, běh mezí a běh SEDE. Spouští se jako detached kontejner se jménem a bez \texttt{-{}-rm}.
\end{itemize}

Žádný běh na vyžádání není součástí buildu ani CI. Jeho výsledkem je záznam v~\texttt{docs/runs/} s~commitem, strojem a prostředím, ne závazek. Soubory běhu zůstávají mimo git. Doby v~tab.~\ref{tab:behy} byly naměřeny na měřicím stroji (Intel Core i7-1065G7, 4~jádra, 16~GB, Docker Desktop nad WSL2, virtuální stroj se čtyřmi procesory a 12~GB) nad commity z~tab.~\ref{tab:zaznamy-behu}. Příkazy jsou v~\path{ORMConvertor/README.md}, \emph{Tests}.

{\footnotesize
\begin{longtable}{P{0.14} P{0.36} P{0.31} P{0.19}}
\toprule
\textbf{Běh} & \textbf{Co dělá} & \textbf{Test a proměnné} & \textbf{Doba} \\
\midrule
\endfirsthead
\toprule
\textbf{Běh} & \textbf{Co dělá} & \textbf{Test a proměnné} & \textbf{Doba} \\
\midrule
\endhead
\rule{0pt}{26pt}\\
\endfoot
\bottomrule
\caption{Běhy na vyžádání.}\label{tab:behy}\\
\endlastfoot
validátor LDBC & přehraje validační sadu nad scale factorem 1 (138~747 řádků) a každé čtení každého frameworku srovná s~očekávanými řádky & \texttt{LdbcValidationTest}, \path{ORMCONVERTOR_LDBC_VALIDATION_ROWS} & 29,6~h (.NET) a 27,2~h (Java) \\
běh FIBEN, běh SEDE & každý dotaz katalogu do všech cílů, srovnání řádků se zdrojem a u~FIBEN s~validátorem, medián času & \texttt{FibenRunTest}, \texttt{SedeRunTest} & FIBEN 2,3~h a 3,1~h, SEDE 7,9~h a 14,4~h \\
experimentální běh & tabulka T2 × T3 (kap.~\ref{sec:reprodukce}) & \texttt{ExperimentRunTest}, \path{ORMCONVERTOR_EXPERIMENT_RESULTS} & asi hodina včetně buildu \\
běh mezí & syntetické řady v~jedné ose (joiny, zanoření, poddotazy, entity\ldots), čas a paměť překladu & \texttt{LimitSearchTest}, \path{ORMCONVERTOR_SERIES_AXES} & .NET 14~h, Java 32~h \\
celý vstup & všechny vstupy sloučené do jednoho projektu na framework, jen čas & \texttt{WholeInputTest} & minuty \\
cizí artefakty & překlad od modelu nebo od člověka projde stupni ověření & \texttt{foreign-run.sh} & podle počtu \\
srovnání s~LLM & protokol o~sedmi krocích (tab.~\ref{tab:protokol}) & \path{ORMCONVERTOR_LLM_RUN} & hodiny až den \\
\end{longtable}}

\begin{table}[H]
\centering\footnotesize
\begin{tabularx}{\textwidth}{L{4.6cm} Y L{5.2cm}}
\toprule
\textbf{Režim serveru} & \textbf{Nastavení} & \textbf{Účel} \\
\midrule
měřicí (výchozí, rozh.~145, 164) & databáze na procesorech 0--1, klienti na 2--3, paměť serveru 4096~MB, \texttt{MAXDOP}~1 & opakovatelná měření času \\
validační (\texttt{judge.env}, rozh.~164, 166) & všechny procesory, jednovláknové plány & souběžná čtení validátoru LDBC \\
\bottomrule
\end{tabularx}%
\caption{Režimy databázového serveru.}
\label{tab:rezimy}
\end{table}

Každý dlouhý běh vyžaduje nepřetržitý běh hostitele: pozastavené aktualizace Windows a vypnuté uspání. Ve validačním režimu musí každý příkaz compose, který pracuje s~databází, obsahovat \path{--env-file judge.env}. Bez něj compose databázi znovu vytvoří v~měřicím režimu a běžící validaci přeruší (kap.~\ref{sec:nepovedlo}). Z~Git Bash potřebují příkazy \texttt{docker compose run} předponu \path{MSYS_NO_PATHCONV=1}, jinak se cesty \texttt{/app/\ldots} v~argumentech \texttt{-e} přepíší na cesty Windows.

Běhy mají stejné ovládání. Běh zapne proměnná, u~validátoru LDBC a běhů katalogů connection string jejich databáze, a soubory se zapíší do adresáře běhu. Běhy FIBEN a SEDE řídí proměnné \path{ORMCONVERTOR_FIBEN_*} a \path{ORMCONVERTOR_SEDE_*}: \texttt{\_RESULTS} jmenuje adresář běhu, \texttt{\_QUERIES} vybere dotazy pro kontrolu a \texttt{\_MEASURE=0} vypne měření času. U~validátoru LDBC určuje \path{ORMCONVERTOR_LDBC_VALIDATION_ROWS} prefix validační sady (prázdná hodnota znamená celou sadu). U~běhu mezí vybírá \path{ORMCONVERTOR_SERIES_AXES} osy a \path{ORMCONVERTOR_SERIES_CATALOG=none} zapne dry run. Kontrola z~hostitele se nezaznamenává, záznam vzniká jen z~běhu v~zafixovaném prostředí. Jedinou výjimkou je kontrola, která změnila vstupy experimentu (tab.~\ref{tab:nepovedlo}).

\subsection{Problémy při nasazení}\label{sec:problemy}

Tab.~\ref{tab:problemy} uvádí možné problémy instalace a běžného použití podle projevu.

{\footnotesize
\begin{longtable}{P{0.24} P{0.24} P{0.22} P{0.30}}
\toprule
\textbf{Projev} & \textbf{Příčina} & \textbf{Diagnostika} & \textbf{Náprava} \\
\midrule
\endfirsthead
\toprule
\textbf{Projev} & \textbf{Příčina} & \textbf{Diagnostika} & \textbf{Náprava} \\
\midrule
\endhead
\rule{0pt}{26pt}\\
\endfoot
\bottomrule
\caption{Problémy při nasazení podle projevu.}\label{tab:problemy}\\
\endlastfoot
\multicolumn{4}{@{}l}{\textbf{Build image selže}} \\
hlášení \emph{No FIBEN data set of commit \ldots} (obdobně pro LDBC, SEDE a WideWorldImporters) & archiv datové sady není na své adrese dostupný & výstup \texttt{docker compose build mssql\_db} & přepínač, který hlášení jmenuje, postaví image bez té databáze (tab.~\ref{tab:sady}). \texttt{-{}-build-arg DATASETS\_URL=\ldots} určí jiné místo vlastních kopií \\
hlášení \emph{is not the archive datasets.sha256 names} & pod jménem archivu přišel jiný obsah & totéž hlášení jmenuje adresu & stáhnout archiv z~adresy v~tab.~\ref{tab:sady}, nebo sadu vynechat přepínačem \\
build tagu \texttt{2.4.0} skončí u~FIBEN & tag bere FIBEN z~repozitáře, který už není dostupný & výstup buildu & \texttt{FIBEN\_DATASET=none}, nebo release \texttt{\Vydani} \\
\midrule
\multicolumn{4}{@{}l}{\textbf{Služba se nespustí nebo aplikace neodpovídá}} \\
start služby skončí chybou \texttt{cpuset} & virtuální stroj Dockeru má méně než čtyři procesory a compose připíná služby na procesory 0--1 a 2--3 & \texttt{docker info}, řádek \emph{CPUs} & \texttt{.wslconfig} se čtyřmi procesory a 12~GB, pak \texttt{wsl -{}-shutdown} a start Docker Desktopu \\
start služby skončí na obsazeném portu 1444 nebo 5072 & port publikuje jiný kontejner nebo proces. Port 1444 publikují také zděděné benchmarky & \texttt{docker ps} & zastavit druhou službu \\
kontejner nebo proces hostitele je ukončen pro nedostatek paměti & systém a profil \texttt{test} najednou přesáhnou 12~GB virtuálního stroje & \texttt{docker stats}, \texttt{docker compose ps} & před profilem \texttt{test} zastavit systém a odstranit zbylé kontejnery profilu (rozh.~193) \\
\texttt{/orm/} vrací prázdnou odpověď 404 & publikovaná instance je spuštěná mimo content root & log aplikace: \emph{Request reached the end of the middleware pipeline without being handled} & \texttt{cd} do adresáře publikace, nebo \texttt{-{}-contentRoot} \\
\texttt{/orm/advisor-test} vrací 400 s~chybou načtení \texttt{libadvisor.so} & Advisor běží mimo image a knihovna pro Windows neexistuje & odpověď endpointu & spustit systém v~kontejneru. Překlad funguje i~bez knihovny \\
\midrule
\multicolumn{4}{@{}l}{\textbf{Katalog není dostupný nebo nedoplňuje}} \\
stav \emph{catalog not configured}. Entita ze zdroje Dapper nebo MyBatis vyjde bez klíče a NHibernate, Hibernate a EclipseLink ji odmítnou záznamem \emph{Failure} & klíč \path{ConnectionStrings:CatalogDatabase} není nastaven & stav katalogu u~převodu, pole \texttt{catalogState} & nastavit klíč (kap.~\ref{sec:konfigurace}). Není to vada: F4 a F6 bez databáze neplatí \\
záznam \emph{Incompleteness} \uv{No table matching the entity was found in the catalog} & katalog čte jinou databázi, než jakou vstup popisuje & jméno databáze v~connection stringu & \path{ORMCONVERTOR_CATALOG_DATABASE} u~compose, \texttt{ConnectionStrings\_\_\allowbreak CatalogDatabase} u~instance na hostiteli \\
převod čeká desítky sekund a skončí stavem \emph{catalog unreachable} & connection string jen s~hostname: SqlClient nejdřív hledá instanci a mez 5~s při tom neplatí (rozh.~198) & pole \texttt{catalogReadMilliseconds} v~odpovědi & port nebo protokol ve zdroji dat (\texttt{localhost,1444}, \texttt{tcp:host}) \\
databáze datové sady na serveru chybí nebo je neúplná & loader skončil před dokončením & \texttt{docker compose logs mssql\_db}: řádek \emph{\ldots loader stopped before it finished} a nad ním hlášení loaderu & podle hlášení loaderu (místo na disku, stažený archiv). Příští start načítá znovu \\
\midrule
\multicolumn{4}{@{}l}{\textbf{Testy neodpovídají zdrojovému kódu nebo selhávají}} \\
sada projde s~jiným počtem testů než v~záznamu & \texttt{run} bez \texttt{build}: sada běžela nad image ze staršího stavu kódu a její výsledek neověřuje aktuální zdrojový kód & počet testů v~souhrnu proti tab.~\ref{tab:cisla} & vždy dvojice \texttt{build} a \texttt{run} \\
diferenční testy javovské sady na hostiteli selhávají & instance pro sadu čte katalog nad vzorovou databází místo testovací & záznam \emph{Incompleteness} v~odpovědi převodu & \texttt{ConnectionStrings\_\_\allowbreak CatalogDatabase} nad testovací databází \\
\end{longtable}}

\section{Použití}\label{sec:pouziti}

ORMorpher převádí zdrojové soubory datové vrstvy aplikace z~jednoho ORM frameworku do jiného. Vstupem jsou soubory s~entitami (třídy, které popisují tabulky databáze), s~mapováním a s~dotazy. Výstupem jsou tytéž entity, mapování a dotazy zapsané pro cílový framework a seznam záznamů o~tom, co převod doplnil, změnil nebo nepřevedl. Výstup není hotový sestavitelný projekt: jeho zařazení do vlastní aplikace je na uživateli.

Kapitoly~\ref{sec:start} až~\ref{sec:vlastni} vedou krok za krokem: spuštění aplikace, první převod na dodaném příkladu a převod vlastních souborů. Kapitoly~\ref{sec:vstup-vystup} až~\ref{sec:stranky} popisují čtení výsledku, řešení problémů, další příklady a ostatní stránky nástroje. Kapitoly~\ref{sec:co-jak} až~\ref{sec:api} jsou reference: podporované konstrukce, meze překladu a API. Uživatelské rozhraní je anglické a nástroj se v~něm jmenuje ORMConvertor.

\subsection{Spuštění aplikace}\label{sec:start}

\textbf{Cíl.} Aplikace běží na počítači uživatele a její úvodní stránka je otevřená v~prohlížeči.

\textbf{Předpoklady.}

\begin{itemize}[leftmargin=*,itemsep=1pt]
\item počítač s~procesorem x86-64 (kap.~\ref{sec:spusteni}), připojení k~internetu a orientačně 30~GB volného místa na disku,
\item nainstalovaný a spuštěný Docker s~Compose a s~alespoň čtyřmi procesory: ve Windows a macOS Docker Desktop, v~Linuxu Docker Engine,
\item Git, nebo archiv ZIP release stažený z~adresy \url{https://github.com/MiraStochel/orm-convertor/archive/refs/tags/2.4.2.zip},
\item volné porty 5072 a 1444 a webový prohlížeč.
\end{itemize}

Jiný software potřeba není. Příkazy (výpis~\ref{lst:start}) se zadávají do terminálu: ve Windows do aplikace PowerShell, v~macOS a Linuxu do aplikace Terminál.

\noindent\begin{minipage}{\linewidth}
\begin{lstlisting}[language=bash,caption={Spuštění aplikace.},label={lst:start}]
docker version                                                               # krok 1
git clone --branch 2.4.2 https://github.com/MiraStochel/orm-convertor.git    # krok 2
cd orm-convertor/ORMConvertor                                                # krok 3
docker compose up --build                                                    # krok 3
\end{lstlisting}
\end{minipage}

\begin{enumerate}[leftmargin=*,itemsep=2pt]
\item \textbf{Ověření Dockeru.} Spusťte Docker Desktop a zadejte \texttt{docker version}. \emph{Správný stav:} výpis obsahuje oddíly \emph{Client} a \emph{Server}. Bez oddílu \emph{Server} Docker neběží.
\item \textbf{Stažení repozitáře.} Zadejte příkaz \texttt{git clone}, nebo rozbalte stažený archiv ZIP. \emph{Správný stav:} vznikla složka \texttt{orm-convertor} (z~archivu \texttt{orm-convertor-2.4.2}) s~podsložkou \texttt{ORMConvertor}.
\item \textbf{Spuštění.} Přejděte do podsložky \texttt{ORMConvertor} a zadejte \texttt{docker compose up -{}-build}. Docker stáhne base image a datové sady, sestaví image aplikace a databáze a spustí obě služby (obr.~\ref{fig:nasazeni}). Terminál vypisuje průběh. První spuštění trvá orientačně desítky minut podle linky a stroje, další spuštění už nic nestahuje. Okno terminálu zůstává otevřené po celou dobu práce s~aplikací.
\item \textbf{Otevření aplikace.} V~prohlížeči otevřete adresu \url{http://localhost:5072/orm/}. \emph{Správný stav:} zobrazí se stránka s~nadpisem \emph{ORMConvertor} a sedmi kartami (obr.~\ref{fig:rozcestnik}). Když se prohlížeč nepřipojí, build ještě neskončil a terminál dál vypisuje jeho kroky. Po jeho dokončení načtěte stránku znovu.
\end{enumerate}

\textbf{Ukončení a další spuštění.} Aplikaci ukončí klávesy Ctrl+C v~okně terminálu nebo příkaz \texttt{docker compose stop} zadaný ve složce \texttt{ORMConvertor}. Data databáze zůstanou zachována. Znovu ji spustí příkaz \texttt{docker compose up}.

Aplikace je určena pro lokální použití: nevyžaduje přihlášení a nešifruje spojení (kap.~\ref{sec:spusteni}). Velikosti, časy a konfiguraci uvádí kap.~\ref{sec:nasazeni}, řešení chyb při spuštění tab.~\ref{tab:problemy}.

\snimek[0.8]{01-rozcestnik}{Úvodní stránka na adrese \texttt{/orm/}. Karta \emph{Translation} vlevo nahoře vede tlačítkem \emph{Open the translator} na překladovou obrazovku, karta \emph{Tutorial} na výukovou stránku. Tytéž stránky nabízí menu vpravo nahoře.}{fig:rozcestnik}

\FloatBarrier
\subsection{První převod}\label{sec:tutorial}

\textbf{Cíl.} Převést příklad dodaný s~nástrojem z~frameworku Dapper do frameworku EF~Core, přečíst výsledek a stáhnout ho. Příklad má dvě třídy jazyka C\# a jeden dotaz v~SQL. Postup trvá asi pět minut.

\textbf{Předpoklady.} Běžící aplikace (kap.~\ref{sec:start}). Znalost obou frameworků potřeba není.

\textbf{Pojmy.} \emph{Jednotka} je jeden vstupní soubor. \emph{Artefakt} je jeden výstupní soubor. \emph{Katalog} je popis tabulek, sloupců a klíčů, který nástroj čte z~připojené databáze. \emph{Záznam} je jedna zpráva nástroje o~tom, co při převodu doplnil nebo nepřevedl.

\begin{enumerate}[leftmargin=*,itemsep=2pt]
\item \textbf{Otevření překladové obrazovky.} Na úvodní stránce klikněte na \emph{Open the translator}. \emph{Výsledek:} stránka \emph{Translation} s~oddíly \emph{Start from an example}, \emph{Direction}, \emph{Input files} a \emph{Convert} pod sebou. Stránka je předvyplněná jiným příkladem, který krok~2 nahradí.
\item \textbf{Načtení příkladu.} V~seznamu oddílu \emph{Start from an example} zvolte \emph{A source that states little (Dapper → EF~Core, 3 files)} a klikněte na \emph{Load example}. \emph{Výsledek} (obr.~\ref{fig:priklad-nacteny}): oddíl \emph{Direction} ukazuje \emph{Dapper} vlevo a \emph{EF~Core} vpravo. Oddíl \emph{Input files} obsahuje tři jednotky (\texttt{Customer.cs}, \texttt{CustomerTransaction.cs} a \texttt{CustomerQuery.sql}) a nad nimi řádek \emph{Dapper input is read as: C\# source, SQL query.} Pod tlačítkem \emph{Convert} je řádek \emph{Loaded \uv{A source that states little}: 3 files.} Pokud se stránka zeptá, zda má nahradit rozpracovaný vstup, potvrďte.
\item \textbf{Převod.} Klikněte na \emph{Convert}. \emph{Výsledek:} převod trvá méně než sekundu. Stránka se posune k~oddílu \emph{Result} a pod tlačítkem je řádek \emph{Converted: 3 artifacts, 33 records.}
\end{enumerate}

\snimekbody[0.9]{02-priklad-nacteny}{Překladová obrazovka po kroku~2. Bod~1: seznam příkladů a tlačítko \emph{Load example}. Bod~2: směr převodu (\emph{Direction}), zdroj vlevo a cíl vpravo. Bod~3: jazyky, ve kterých se vstup zdroje čte. Bod~4: přidání vlastních souborů, archivu ZIP nebo prázdné jednotky. Bod~5: jednotka s~názvem, jazykem a editorem, do kterého se vstup vkládá. Bod~6: tlačítko \emph{Convert}, které spustí převod (krok~3).}{fig:priklad-nacteny}{%
  \bod{1}{0.521}{0.099}\bod{2}{0.316}{0.176}\bod{3}{0.243}{0.285}%
  \bod{4}{0.279}{0.310}\bod{5}{0.530}{0.376}\bod{6}{0.094}{0.981}}

\begin{enumerate}[resume,leftmargin=*,itemsep=2pt]
\item \textbf{Kontrola hlavičky výsledku.} První řádek oddílu \emph{Result} (obr.~\ref{fig:tutorial}) uvádí směr s~verzemi obou frameworků (\emph{Dapper 2.1.79 → EF~Core 10.0.10}), údaj \emph{source SQL dialect not declared}, stav katalogu, verzi nástroje a identifikátor běhu. \emph{Správný stav:} štítek \emph{catalog read} s~dobou čtení. Jiný stav katalogu vysvětluje tab.~\ref{tab:katalog-stavy}.
\item \textbf{Porovnání vstupu a výstupu.} Pod hlavičkou stojí vlevo sloupec \emph{What you sent} se vstupem a vpravo sloupec \emph{What came back} s~výstupem. Každý panel výstupu uvádí na řádku \emph{from} jednotku, ze které vznikl. Výpisy~\ref{lst:ex1-in} a~\ref{lst:ex1-out} ukazují vstup a výstup. Vstupní třída neuvádí tabulku, typy sloupců ani klíč. Výstupní třída je obsahuje jako atributy EF~Core, které nástroj doplnil z~katalogu (tab.~\ref{tab:ex1-zmeny}). Dotaz v~SQL vyšel jako metoda v~LINQ.
\end{enumerate}

\snimek{03-vysledek}{Výsledek převodu (kroky~4 a~5). Nahoře hlavička se stavem katalogu a tlačítko \emph{Download all as ZIP}, vlevo vstup, vpravo výstup. Podtržená místa výstupu jsou údaje z~tabulky faktů.}{fig:tutorial}

\noindent\begin{minipage}{\linewidth}
\begin{lstlisting}[language={[Sharp]C},caption={První převod, vstup: \texttt{Customer.cs}, \texttt{CustomerTransaction.cs} (zkráceno) a \texttt{CustomerQuery.sql}.},label={lst:ex1-in}]
namespace DapperEntities;
public class Customer
{
    public int CustomerID { get; set; }
    public required string CustomerName { get; set; }
    public DateTime AccountOpenedDate { get; set; }
    public decimal? CreditLimit { get; set; }
    public List<CustomerTransaction> Transactions { get; set; } = [];
}
public class CustomerTransaction
{
    public int CustomerTransactionID { get; set; }
    public int CustomerID { get; set; }
    public DateTime TransactionDate { get; set; }
    public decimal TransactionAmount { get; set; }
}
-- CustomerQuery.sql
SELECT c.CustomerName, c.CreditLimit
FROM Sales.Customers AS c
WHERE c.CreditLimit > 2000
ORDER BY c.AccountOpenedDate DESC
\end{lstlisting}
\end{minipage}

\noindent\begin{minipage}{\linewidth}
\begin{lstlisting}[language={[Sharp]C},caption={První převod, výstup: \texttt{Customer.cs} a metoda dotazu.},label={lst:ex1-out}]
[Table("Customers", Schema = "Sales")]
[Index(nameof(CustomerName), IsUnique = true, Name = "UQ_Sales_Customers_CustomerName")]
public class Customer
{
    [Key]
    [Column(TypeName="int")]
    public required int CustomerID { get; set; }

    [Column(TypeName="nvarchar")]
    [MaxLength(100)]
    public required string CustomerName { get; set; }

    [Column(TypeName="date")]
    public required DateTime AccountOpenedDate { get; set; }

    [Column(TypeName="decimal")]
    [Precision(18, 2)]
    public decimal? CreditLimit { get; set; }

    public List<CustomerTransaction> Transactions { get; set; } = [];
}

public static IQueryable Query(DbContext ctx)
{
    return ctx.Set<Customer>()
        .Where(c => c.CreditLimit > 2000)
        .OrderByDescending(c => c.AccountOpenedDate)
        .Select(c => new { CustomerName = c.CustomerName, CreditLimit = c.CreditLimit });
}
\end{lstlisting}
\end{minipage}

\begin{table}[H]
\centering\footnotesize
\begin{tabularx}{\textwidth}{L{6.4cm} Y}
\toprule
\textbf{Řádek výstupu} & \textbf{Význam} \\
\midrule
\texttt{[Table("Customers", Schema = "Sales")]} & třída patří k~tabulce \texttt{Sales.Customers} \\
\texttt{[Key]} & sloupec \texttt{CustomerID} je primární klíč \\
\texttt{[Column(TypeName="nvarchar")]}, \texttt{[MaxLength(100)]} & typ a délka sloupce \\
\texttt{[Precision(18, 2)]} & přesnost a měřítko desetinného čísla \\
\texttt{required}, \texttt{decimal?} & sloupec hodnotu vyžaduje, nebo ji mít nemusí \\
\texttt{[Index(\ldots, IsUnique = true, \ldots)]} & unikátní omezení nad sloupcem \texttt{CustomerName} \\
\midrule
\texttt{ctx.Set<Customer>()} & \texttt{FROM Sales.Customers} \\
\texttt{.Where}, \texttt{.OrderByDescending}, \texttt{.Select} & \texttt{WHERE}, \texttt{ORDER BY \ldots DESC}, seznam sloupců za \texttt{SELECT} \\
\bottomrule
\end{tabularx}%
\caption{Co výstup obsahuje navíc proti vstupu. Atributy pocházejí z~katalogu databáze.}
\label{tab:ex1-zmeny}
\end{table}

\begin{enumerate}[resume,leftmargin=*,itemsep=2pt]
\item \textbf{Tabulka faktů.} Pás \emph{Facts} pod oběma sloupci (obr.~\ref{fig:fakta}) má 27 řádků, jeden za každý údaj mapování. Sloupec \emph{Origin} uvádí původ údaje a sloupec \emph{Outcome}, zda je údaj ve výstupu. \emph{Správný stav:} všech 27 řádků má původ \emph{supplied by the catalog} a stav \emph{carried}. Všechny údaje tedy doplnil katalog a všechny jsou ve výstupu. Ukázání na řádek zvýrazní odpovídající místo ve výstupu.
\item \textbf{Záznamy.} Pás \emph{Records} (obr.~\ref{fig:zaznamy}) má 33 záznamů tří druhů (tab.~\ref{tab:ex1}). Kliknutí na štítek s~počtem zobrazí jen záznamy toho druhu. \emph{Správný stav:} žádný záznam druhu \emph{Failure}, tedy všechny tři artefakty vznikly.
\end{enumerate}

\snimek[0.9]{04-fakta}{Pás \emph{Facts} (krok~6): jeden řádek za údaj mapování se sloupci entita, vlastnost, druh údaje, hodnota, původ (\emph{Origin}) a stav (\emph{Outcome}).}{fig:fakta}

\snimek{05-zaznamy}{Pás \emph{Records} (krok~7). Nahoře štítky s~počty po druzích a jejich vysvětlení, pod nimi tabulka záznamů, zde filtrovaná na druh \emph{Incompleteness}.}{fig:zaznamy}

\begin{table}[H]
\centering\footnotesize
\begin{tabularx}{\textwidth}{L{2.6cm} Y L{4.6cm}}
\toprule
\textbf{Druh a počet} & \textbf{Příklad záznamu} & \textbf{Co uživatel udělá} \\
\midrule
\emph{Supplied} (27) & \emph{Length 100 supplied by the database catalog from \textquotesingle{}Sales.Customers\textquotesingle{}.} Údaj není ve vstupu a nástroj ho převzal z~databáze & nic, pokud je připojená databáze databází aplikace \\
\emph{Incompleteness} (4) & \emph{The catalog states the foreign key FK\_Sales\_CustomerTransactions\_CustomerID\_Sales\_Customers towards \textquotesingle{}Customer\textquotesingle{}, but the entity has no navigation property of that type; the relation is not generated.} Dva záznamy uvádějí cizí klíč bez vlastnosti v~entitě, dva strategii generování klíče, kterou katalog neobsahuje (rozh.~064) & doplní do entity vlastnost vztahu, pokud vztah potřebuje \\
\emph{Convention} (2) & \emph{No generation strategy was stated; EF Core generates the value of a single-property integer or Guid key on its own, which is a convention of the target, not a fact of the source.} & ověří, že konvence EF~Core odpovídá tomu, jak databáze klíč plní \\
\bottomrule
\end{tabularx}%
\caption{Záznamy prvního převodu. Bez katalogu vyjdou dvě entity s~atributem \texttt{[Keyless]} a tři záznamy.}
\label{tab:ex1}
\end{table}

\begin{enumerate}[resume,leftmargin=*,itemsep=2pt]
\item \textbf{Stažení.} Klikněte na \emph{Download all as ZIP} pod hlavičkou výsledku. \emph{Výsledek:} prohlížeč uloží archiv se třemi soubory (výpis~\ref{lst:zip}). Archiv se jmenuje \texttt{conversion.zip}.
\end{enumerate}

\noindent\begin{minipage}{\linewidth}
\begin{lstlisting}[caption={Obsah souboru \texttt{conversion.zip} po prvním převodu.},label={lst:zip}]
conversion.zip
|-- Customer.cs               entita EF Core
|-- CustomerTransaction.cs    entita EF Core
\-- query.cs                  metoda dotazu v LINQ, bez třídy
\end{lstlisting}
\end{minipage}

\textbf{Použití výstupu} je mimo nástroj. Kdo chce výstup spustit ve vlastní aplikaci, vloží soubory entit a metodu ze souboru \texttt{query.cs} do svého projektu s~EF~Core~10, který už má třídu kontextu a připojení k~databázi. Výstup je správný, když se projekt sestaví a dotaz vrátí tytéž řádky jako původní dotaz v~SQL. Tytéž kontroly dělají testy nástroje ve čtyřech stupních (kap.~\ref{sec:metodika}), převod přes rozhraní je nedělá.

\textbf{Shrnutí.} Uživatel provedl čtyři akce: otevření překladové obrazovky, načtení příkladu, převod a stažení. Výsledek se čte v~pořadí stav katalogu, porovnání vstupu s~výstupem, tabulka faktů a záznamy. Týž postup nabízí stránka \emph{Tutorial} v~nástroji (obr.~\ref{fig:vyukova-stranka}): její příklad je živý a její krok~4 ho převede na téže instanci.

\snimek[0.8]{06-vyukova-stranka}{Výuková stránka \emph{Tutorial}: tentýž převod v~pěti krocích.}{fig:vyukova-stranka}

\FloatBarrier
\subsection{Převod vlastních souborů}\label{sec:vlastni}

\textbf{Cíl.} Převést soubory vlastního projektu.

\textbf{Předpoklady.} Běžící aplikace a soubory datové vrstvy: entity, mapování a kód s~dotazy. Třídy, které entitami nejsou (služby, DTO), do vstupu nepatří.

\begin{enumerate}[leftmargin=*,itemsep=2pt]
\item \textbf{Směr.} V~oddílu \emph{Direction} zvolte zdrojový framework vlevo a cílový vpravo. Směr se volí jako první, protože určuje, které soubory vstup přijme. \emph{Výsledek:} řádek pod nadpisem \emph{Input files} vyjmenuje jazyky, které zdroj čte, například \emph{EF~Core input is read as: C\# source.} Co každý framework čte, uvádí tab.~\ref{tab:frameworky}.
\item \textbf{Dialekt SQL.} Seznam pod směrem ponechte na hodnotě \emph{Source SQL dialect: not declared}. Hodnota \emph{Source SQL is another database system} je určena pro zdroj, jehož SQL je psané pro jinou databázi než SQL Server: nástroj pak toto SQL nečte a dotaz odmítne záznamem (tab.~\ref{tab:chyby-vstup}).
\item \textbf{Soubory.} Tlačítko \emph{Add files or a ZIP\ldots} přidá soubory z~disku, tlačítko \emph{Add empty unit} prázdnou jednotku, do které lze text vložit. Každý soubor je jedna jednotka s~názvem, jazykem a editorem. Jazyk určuje přípona: \texttt{.cs}, \texttt{.java}, \texttt{.xml}, \texttt{.sql}, \texttt{.hql} a \texttt{.jpql}. U~jiné přípony se jazyk zvolí v~seznamu u~jednotky. \emph{Remove} odebere jednotku, \emph{Clear} všechny. Vstup zůstává uložen v~prohlížeči i~po znovunačtení stránky.
\item \textbf{Archiv ZIP.} Archiv projektu se rozbalí v~prohlížeči a může mít libovolnou strukturu složek. Každý soubor v~jazyce, který zdroj čte, se stane jednotkou pojmenovanou cestou v~archivu. Ostatní soubory vyjmenuje seznam \emph{\ldots files from \ldots not used} s~důvodem (obr.~\ref{fig:archiv-zip}), například: přípona neurčuje jazyk, zdroj tento jazyk nečte, soubor je prázdný nebo šifrovaný. Nesprávné složení vstupu se opraví odebráním jednotek tlačítkem \emph{Remove}, nebo tlačítkem \emph{Clear}, změnou směru a novým přidáním archivu. Přidaný archiv se k~jednotkám na stránce připojuje, nenahrazuje je.
\item \textbf{Převod.} Klikněte na \emph{Convert}. Před odesláním stránka zkontroluje, že vstup má aspoň jednu jednotku, že žádná není prázdná, že každá má jazyk, který zdroj čte, a že XML je správně utvořené. Nálezy vypíše pod nadpisem \emph{The input is not ready to send} a převod neodešle. Syntaktické chyby v~C\#, Javě a SQL hlásí server záznamem \emph{Failure} s~řádkem a sloupcem. Chyba v~jedné jednotce převod ostatních nezastaví.
\item \textbf{Čtení výsledku.} Části výsledku jsou tytéž jako v~kap.~\ref{sec:tutorial}. U~vlastních souborů se čtou v~pořadí stav katalogu, záznamy \emph{Failure}, porovnání vstupu s~výstupem, fakta se stavem jiným než \emph{carried} a ostatní záznamy (kap.~\ref{sec:vstup-vystup}).
\end{enumerate}

Katalog instance spuštěné podle kap.~\ref{sec:start} čte vzorovou databázi WideWorldImporters. Pro vlastní projekt doplní fakta jen tehdy, když instance čte databázi tohoto projektu (kap.~\ref{sec:konfigurace}). Bez ní převod použije jen to, co uvádí zdroj, a konvence cíle.

\snimek{09-archiv-zip}{Archiv ZIP po přidání (krok~4): rozbalený seznam nepoužitých souborů s~důvodem u~každého a pod ním jednotky pojmenované cestou v~archivu.}{fig:archiv-zip}

\FloatBarrier
\subsection{Výsledek převodu a diagnostické záznamy}\label{sec:vstup-vystup}

Výsledek má čtyři části: hlavičku se záznamem běhu, artefakty vedle vstupu, tabulku faktů a záznamy. Převod je deterministický (S2): týž vstup nad týmž katalogem a touž verzí nástroje dává bajtově shodné artefakty i~záznamy. Test \texttt{RunRecordTest} to ověřuje ve všech 36 směrech na převodu bez katalogu.

\textbf{Stav katalogu} v~hlavičce uvádí, zda měl převod k~dispozici databázi (tab.~\ref{tab:katalog-stavy}).

\begin{table}[H]
\centering\footnotesize
\begin{tabularx}{\textwidth}{L{3.3cm} Y L{5.6cm}}
\toprule
\textbf{Štítek} & \textbf{Význam} & \textbf{Co uživatel udělá} \\
\midrule
\emph{catalog read} & fakta, která zdroj neuvádí, byla doplněna z~katalogu & nic \\
\emph{catalog unused} & katalog je připojen, převod se ho neměl na co ptát & nic \\
\emph{catalog not configured} & instance nemá připojení k~databázi, převod použil jen konvence & nastaví katalog (kap.~\ref{sec:konfigurace}). Instance v~Dockeru ho má nastavený \\
\emph{catalog unreachable} & připojení je nastaveno, čtení selhalo a převod pokračoval s~konvencemi & ověří, že databáze běží, a převod zopakuje (tab.~\ref{tab:problemy}) \\
\bottomrule
\end{tabularx}%
\caption{Stavy katalogu v~hlavičce výsledku.}
\label{tab:katalog-stavy}
\end{table}

\textbf{Tabulka faktů} (pás \emph{Facts}) má řádek za každý mapovací fakt: tabulku, sloupec, typ, klíč, vztah. Původ je jméno jednotky, která fakt uvádí (bez známé jednotky \emph{stated by the input}), \emph{supplied by the catalog} nebo \emph{derived by the tool}. Stav je \emph{carried} (fakt je ve výstupu), \emph{dropped} (cíl ho nevyjádří) nebo \emph{not generated} (artefakt entity nevznikl). Ukázání na řádek zvýrazní fakt ve vstupu i~ve výstupu (obr.~\ref{fig:kniha-tvrzeni}).

\textbf{Záznamy} (pás \emph{Records}, požadavek F11) mají osm druhů. Tab.~\ref{tab:zaznamy} je řadí podle závažnosti. Fakt, který cíl nevyjádří, použitá konvence, rozpor mezi zdroji a neúplný vstup mají vždy záznam.

\begin{table}[H]
\centering\footnotesize
\begin{tabularx}{\textwidth}{L{2.1cm} Y L{5.2cm}}
\toprule
\textbf{Druh} & \textbf{Význam} & \textbf{Co uživatel udělá} \\
\midrule
\multicolumn{3}{@{}l}{\textbf{Vyžaduje zásah}} \\
\emph{Failure} & artefakt nevznikl: jednotku nelze přečíst, cíl vyžaduje fakt, který chybí, nebo by dotaz vrátil jiné řádky & opraví vstup, nebo artefakt napíše ručně \\
\midrule
\multicolumn{3}{@{}l}{\textbf{Vyžaduje kontrolu}} \\
\emph{Loss} & zdroj fakt uvádí a cíl nebo mezireprezentace ho nevyjádří. Artefakt vznikl bez něj & doplní konstrukci v~cílovém projektu ručně, pokud ji potřebuje \\
\emph{Incompleteness} & chybí fakt, který cíl potřebuje: vztah bez sloupců, neznámý typ, tabulka nenalezená v~katalogu & doplní fakt do zdroje nebo do výstupu \\
\emph{Conflict} & dva zdroje uvádějí pro týž fakt různé hodnoty. Platí priorita zdrojů (rozh.~017) & ověří, která hodnota je správná \\
\emph{Convention} & výstup obsahuje hodnotu, kterou zdroj neuvádí a kterou určila konvence cíle nebo nástroje & ověří, že hodnota odpovídá databázi \\
\midrule
\multicolumn{3}{@{}l}{\textbf{Informativní}} \\
\emph{Supplied} & fakt dodal katalog databáze, ne zdroj (F4--F6) & nic \\
\emph{Fallback} & dotaz vyšel celý v~nativním SQL, protože jazyk cíle konstrukci nevyjádří (rozh.~113). Vrací tytéž řádky & počítá s~vazbou artefaktu na SQL Server \\
\emph{Omitted} & tvar LINQ, který provider NHibernatu nevyjádří. Dotaz vyšel jen jako HQL & nic \\
\bottomrule
\end{tabularx}%
\caption{Druhy diagnostických záznamů podle závažnosti.}
\label{tab:zaznamy}
\end{table}

Čtyři z~pěti druhů, které vyžadují zásah nebo kontrolu, na příkladech:

\begin{itemize}[leftmargin=*,itemsep=1pt]
\item \emph{Failure.} Třída Dapperu bez klíče převáděná do Hibernatu na instanci bez katalogu: \emph{The target requires PrimaryKey and neither the source nor a catalog supplied it; the entity\textquotesingle{}s artifacts are not generated.} Náprava: instance s~katalogem nad databází aplikace.
\item \emph{Loss.} Vlastnost \texttt{@Embedded private Address address;} v~entitě JPA: \emph{The property is mapped by @Embedded, which the intermediate representation does not carry (excluded area 2); it is left out of the entity with everything stated about it (decision 177).} Vlastnost ve výstupu chybí a uživatel ji doplní ručně.
\item \emph{Incompleteness.} Vlastnost \texttt{private java.time.Instant createdAt;} převáděná do EF~Core: typ není ve slovníku typů a výstup opakuje jméno \texttt{java.time.Instant}, které C\# nepřijme. Uživatel ho ve výstupu nahradí typem cíle.
\item \emph{Convention.} Klíč EF~Core bez uvedené strategie převáděný do Hibernatu: artefakt obsahuje strategii \texttt{IDENTITY} (rozh.~187). Uživatel ověří, že sloupec v~databázi je \texttt{IDENTITY}.
\end{itemize}

Tab.~\ref{tab:chyby-vypocet} uvádí jevy, které uživatel ve výstupu najde, a odlišuje doplněnou informaci, konvenci, ztrátu a selhání.

{\footnotesize
\begin{longtable}{P{0.25} P{0.14} P{0.41} P{0.20}}
\toprule
\textbf{Jev} & \textbf{Stav} & \textbf{Příčina} & \textbf{Záznam} \\
\midrule
\endfirsthead
\toprule
\textbf{Jev} & \textbf{Stav} & \textbf{Příčina} & \textbf{Záznam} \\
\midrule
\endhead
\rule{0pt}{26pt}\\
\endfoot
\bottomrule
\caption{Diagnostické stavy převodu: jev ve výstupu, jeho příčina a záznam, který ho hlásí.}\label{tab:chyby-vypocet}\\
\endlastfoot
výstup obsahuje atributy, které ve vstupu nejsou & doplněná informace & doplnění metadat ze schématu databáze. Zdroj každého atributu uvádí záznam, například \emph{Length 100 supplied by the database catalog from \textquotesingle{}Sales.Customers\textquotesingle{}} pro \texttt{[MaxLength(100)]} & \emph{Supplied} \\
klíč bez strategie generování & konvence & katalog neobsahuje údaj, jakým mechanismem se sloupec plní (rozh.~064). Cíl použije svou konvenci & \emph{Incompleteness} a \emph{Convention} \\
vztah, který obsahuje jen katalog & neúplnost & cizí klíč bez navigační vlastnosti vztah negeneruje & \emph{Incompleteness} \\
dotaz vyšel v~nativním SQL (obr.~\ref{fig:fallback}) & fallback & jazyk cíle konstrukci nevyjádří. Artefakt vrací tytéž řádky, ale je vázán na SQL Server & \emph{Fallback} \\
chybí inicializátor vlastnosti & ztráta & inicializátor, který není literál shodný v~obou jazycích, se vypouští (například \texttt{= new ArrayList<>()} ze zdroje v~Javě) & \emph{Loss} \\
verze je jen ke čtení & ztráta & \texttt{@Version int} vyjde v~EF~Core jako \texttt{[ConcurrencyCheck]}. Zvyšovat ji musí aplikace & \emph{Loss} \\
třída kontextu EF~Core ve výstupu chybí & ztráta & třída \texttt{DbContext} se čte jako konfigurace, ne jako entita (rozh.~132) & \emph{Loss} \\
dotaz v~LINQ je delší než zdroj & bez ztráty & \texttt{SUM} jazyka SQL nad prázdnou skupinou je \texttt{NULL}, \texttt{Sum} EF~Core nula. Builder píše tvar se stejným výsledkem (rozh.~155) & -- \\
artefakt chybí bez chyby ve vstupu & selhání & výjimka při překladu, tedy vada nástroje. Zbytek převodu doběhne (rozh.~197) & \emph{Failure} \\
\end{longtable}}

\snimek{07-podtrzeny-fakt}{Fakt \texttt{Precision(18, 2)} z~tabulky faktů zvýrazněný ve vstupu vlevo a na dvou místech výstupu vpravo (EF~Core → Hibernate).}{fig:kniha-tvrzeni}

\snimek{12-fallback}{Dotaz zapsaný fallbackem: dva záznamy \emph{Fallback} (Hibernate → EF~Core).}{fig:fallback}

\FloatBarrier
\subsection{Řešení problémů se vstupem}\label{sec:chyby-vstup}

Tab.~\ref{tab:chyby-vstup} je řazena podle projevu. První tři řádky hlásí stránka před odesláním, ostatní jsou záznamy ve výsledku. Vyřešený případ: po odstranění středníku na řádku~5 jednotky \texttt{Customer.cs} převod vrátí záznam \emph{Failure} s~pozicí, stránka řádek označí a odkaz \emph{go to line~5} na něj přejde (obr.~\ref{fig:chyba-radek}). Po opravě řádku a novém kliknutí na \emph{Convert} výsledek záznam neobsahuje. Zbylé dvě jednotky se převedly i~při chybě.

Syntaktická chyba a záměrné odmítnutí jsou dva různé případy. Dotaz psaný pro jinou databázi skončí bez deklarace dialektu syntaktickou chybou jen tam, kde text není syntakticky T-SQL (například \texttt{LIMIT}). Text, který gramatika T-SQL přijme, se přeloží jako T-SQL, i~když volá funkci jiného systému. S~deklarací \emph{Source SQL is another database system} (v~API \texttt{"declaredSourceDialect": 1000}) nástroj text nečte vůbec a odmítne ho záznamem \emph{Failure}. Entity se převedou v~obou případech.

{\footnotesize
\begin{longtable}{P{0.31} P{0.22} P{0.21} P{0.26}}
\toprule
\textbf{Projev (hlášení)} & \textbf{Vstup, který ho vyvolá} & \textbf{Příčina} & \textbf{Oprava} \\
\midrule
\endfirsthead
\toprule
\textbf{Projev (hlášení)} & \textbf{Vstup, který ho vyvolá} & \textbf{Příčina} & \textbf{Oprava} \\
\midrule
\endhead
\rule{0pt}{26pt}\\
\endfoot
\bottomrule
\caption{Problémy se vstupem podle projevu.}\label{tab:chyby-vstup}\\
\endlastfoot
\emph{Customer.cs: the content is empty.} & jednotka bez textu & stránka prázdnou jednotku neodešle & vložit obsah, nebo \emph{Remove} \\
\emph{choose the language of the unit - its file extension names none} & soubor \texttt{notes.txt} & přípona neurčuje jazyk & zvolit jazyk v~seznamu u~jednotky \\
\emph{Dapper does not read HQL query} & soubor \texttt{.hql} při zdroji Dapper & zdroj tento jazyk nečte & změnit zdroj v~oddílu \emph{Direction}, nebo jednotku odebrat \\
\midrule
\emph{The C\# source could not be read at line 5, column 38: \{ or ; or => expected.} & \texttt{public int CustomerID \{ get; set \}} & syntaktická chyba ve zdroji & \texttt{\{ get; set; \}} \\
\emph{The C\# source could not be read at line 1, column 1: Invalid expression term \textquotesingle{}<\textquotesingle{}.} & mapování XML v~jednotce s~jazykem \emph{C\# source} & jazyk jednotky neodpovídá obsahu & zvolit u~jednotky jazyk \emph{XML document} \\
\emph{The SQL could not be parsed at line 1, column 73: Expected OFFSET but encountered LIMIT instead.} & \texttt{SELECT \ldots{} ORDER BY c.CustomerName LIMIT 10} & dotaz není T-SQL & \texttt{SELECT TOP 10 \ldots{} ORDER BY c.CustomerName} \\
\emph{The source declares that its literal SQL is written for another database system, which this version does not read} & jakýkoli dotaz v~SQL s~deklarací jiného systému & záměrné odmítnutí: jediný dialekt nástroje je SQL Server 2022 (rozh.~088) & přepsat dotaz do T-SQL a deklaraci zrušit \\
\emph{The target requires PrimaryKey and neither the source nor a catalog supplied it; the entity\textquotesingle{}s artifacts are not generated.} & třída Dapperu nebo MyBatisu, cíl NHibernate, Hibernate nebo EclipseLink & zdroj klíč neuvádí a katalog není připojen nebo tabulku neobsahuje & instance s~katalogem nad databází aplikace, nebo cíl EF~Core, kde entita vyjde s~\texttt{[Keyless]} \\
\emph{No table matching the entity was found in the catalog; its mapping facts cannot be completed.} & třída \texttt{CustomerDto}, nebo entita z~jiné databáze & třída není entita, nebo katalog čte jinou databázi & třídu vynechat, nebo nastavit katalog (kap.~\ref{sec:konfigurace}) \\
\emph{An earlier source states length 50, a later one 80. \ldots{} the first value is kept.} & dvě jednotky deklarují entitu \texttt{Customer} & fakta se sloučí a rozpor rozhodne pořadí jednotek (rozh.~017) & druhou definici odebrat \\
\emph{No source unit yielded an entity or a query to generate from} & vstup bez entity i~bez dotazu & žádná jednotka neobsahuje to, co její jazyk deklaruje & zkontrolovat jazyk a obsah jednotek \\
\end{longtable}}

\snimek{08-chyba-radek}{Syntaktická chyba ve vstupu: označený řádek~5, hlášení s~řádkem a sloupcem a odkaz \emph{go to line~5}.}{fig:chyba-radek}

\FloatBarrier
\subsection{Praktické příklady převodu}\label{sec:priklady}

Všechny příklady jsou v~nástroji k~načtení volbou \emph{Start from an example}. Výstupy jsou skutečné artefakty release \texttt{\Vydani} s~katalogem WideWorldImporters, zkrácené o~gettery a settery. Příklad~1 je první převod z~kap.~\ref{sec:tutorial}. Tab.~\ref{tab:priklady} uvádí, co každý příklad ukazuje.

\begin{table}[H]
\centering\footnotesize
\begin{tabularx}{\textwidth}{L{3.2cm} Y Y L{1.5cm}}
\toprule
\textbf{Příklad} & \textbf{Vstup} & \textbf{Co příklad ukazuje} & \textbf{Záznamů} \\
\midrule
1: Dapper → EF~Core & dvě třídy bez atributů a jeden dotaz v~SQL & doplnění faktů z~katalogu, čtení výsledku & 33 \\
2: EF~Core → Hibernate & čtyři entity s~atributy a šest dotazů v~LINQ & převod mezi ekosystémy: složený klíč, typy, dotazy & 9 \\
3: EclipseLink → EF~Core (pokročilý) & deset entit JPA a jeden dotaz v~JPQL se sedmi joiny a čtrnácti poddotazy & chování na rozsáhlém dotazu & 29 \\
\bottomrule
\end{tabularx}%
\caption{Příklady převodu.}
\label{tab:priklady}
\end{table}

\subsubsection*{Příklad 2: EF~Core → Hibernate}

\textbf{Cíl.} Převést datovou vrstvu z~.NET do Javy a zjistit, jak se změní složený klíč, cizí klíč, typy a dotaz. \textbf{Vstup.} Příklad \emph{Order book: four entities, six LINQ queries}: zákazník, objednávka, řádek objednávky se složeným klíčem, produkt a šest dotazů. Zdroj uvádí tabulky, klíče i~vztahy sám, takže převod katalog nepotřebuje. \textbf{Postup.} Kroky 2 a~3 z~kap.~\ref{sec:tutorial} s~tímto příkladem.

Výpisy~\ref{lst:ex2-in} a~\ref{lst:ex2-out} ukazují jednu entitu a jeden dotaz. Čísla v~komentářích doplňuje zpráva a označují konstrukce, jejichž překlad vysvětluje tab.~\ref{tab:ex2}.

\begin{lstlisting}[language={[Sharp]C},caption={Příklad 2, vstup (výběr): \texttt{OrderLine.cs} a \texttt{OpenOrders.cs}.},label={lst:ex2-in}]
[Table("OrderLines", Schema = "Ordering")]
[PrimaryKey(nameof(SalesOrderId), nameof(LineNumber))]                              // (1)
public class OrderLine
{
    public int SalesOrderId { get; set; }
    public short LineNumber { get; set; }
    [ForeignKey(nameof(SalesOrderId))] public required SalesOrder Order { get; set; }   // (2)
    public int ProductId { get; set; }
    [ForeignKey(nameof(ProductId))] public required Product Product { get; set; }
    public int Quantity { get; set; }
    [Precision(18, 2)] public decimal UnitPrice { get; set; }                       // (3)
    [Precision(5, 2)] public decimal DiscountPercent { get; set; }
}

public List<SalesOrder> OpenOrders(int customerId, DateTime placedAfter, int skip, int take)
{
    return ctx.SalesOrders
        .Where(o => o.CustomerId == customerId && o.OrderedAt >= placedAfter
                 && new[] { "New", "Paid", "Packed" }.Contains(o.Status))           // (4)
        .OrderByDescending(o => o.OrderedAt).ThenBy(o => o.OrderNumber)
        .Skip(skip).Take(take).ToList();                                            // (5)
}
\end{lstlisting}

\begin{lstlisting}[language=Java,caption={Příklad 2, výstup (výběr): \texttt{OrderLine.java} bez getterů a setterů a metoda dotazu.},label={lst:ex2-out}]
@Entity
@Table(name = "OrderLines", schema = "Ordering")
@IdClass(OrderLine.OrderLineId.class)                                               // (1)
public class OrderLine {
    @Id @Column(name = "SalesOrderId") private Integer SalesOrderId;                // (1)
    @Id @Column(name = "LineNumber") private Short LineNumber;                      // (1)
    @Column(name = "ProductId", nullable = false, insertable = false, updatable = false) private int ProductId;
    @Column(name = "Quantity", nullable = false) private int Quantity;
    @Column(name = "UnitPrice", precision = 18, scale = 2, nullable = false) private BigDecimal UnitPrice;   // (3)
    @Column(name = "DiscountPercent", precision = 5, scale = 2, nullable = false) private BigDecimal DiscountPercent;
    @ManyToOne(optional = false)
    @JoinColumn(name = "SalesOrderId", referencedColumnName = "SalesOrderId", nullable = false,
                insertable = false, updatable = false)                              // (2)
    private SalesOrder Order;
    @ManyToOne(optional = false)
    @JoinColumn(name = "ProductId", referencedColumnName = "ProductId", nullable = false)
    private Product Product;

    public static class OrderLineId implements Serializable {                       // (1)
        private Integer SalesOrderId;
        private Short LineNumber;
        /* konstruktory, equals, hashCode */
    }
}

public static TypedQuery<SalesOrder> query(EntityManager em, int customerId, LocalDateTime placedAfter,
                                           int skip, int take) {
    return em.createQuery("""
        select o
        from SalesOrder o
        where o.CustomerId = :customerId and o.OrderedAt >= :placedAfter and o.Status in ('New', 'Paid', 'Packed')
        order by o.OrderedAt desc, o.OrderNumber asc
        """, SalesOrder.class)                                                      // (4)
        .setParameter("customerId", customerId)
        .setParameter("placedAfter", placedAfter)
        .setFirstResult(skip)                                                       // (5)
        .setMaxResults(take);
}
\end{lstlisting}

\begin{table}[H]
\centering\footnotesize
\begin{tabularx}{\textwidth}{L{0.6cm} L{5.4cm} Y}
\toprule
& \textbf{Zdroj (EF~Core)} & \textbf{Cíl (Hibernate)} \\
\midrule
(1) & složený klíč jedním atributem třídy \texttt{[PrimaryKey(a, b)]} & \texttt{@IdClass}, dvě vlastnosti s~\texttt{@Id} a vnořená třída klíče \texttt{OrderLineId} \\
(2) & cizí klíč, který je zároveň částí klíče & \texttt{@JoinColumn} s~\texttt{insertable = false, updatable = false}, aby sloupec zapisovala jen vlastnost klíče \\
(3) & \texttt{[Precision(18, 2)] decimal} & \texttt{@Column(precision = 18, scale = 2)} a typ \texttt{BigDecimal}. Obdobně \texttt{DateTime} → \texttt{LocalDateTime} \\
(4) & pole hodnot a \texttt{Contains} & \texttt{in (\ldots)} v~textu JPQL. Parametry metody jsou pojmenované parametry dotazu \\
(5) & \texttt{Skip} a \texttt{Take} & \texttt{setFirstResult} a \texttt{setMaxResults} na objektu dotazu \\
\midrule
& \textbf{Záznam} & \textbf{Za co} \\
\midrule
& \emph{Loss} (2) & neunikátní index, který mezireprezentace nenese (rozh.~130), a inicializátor \texttt{[]} \\
& \emph{Incompleteness} (4) & katalog instance (WideWorldImporters) tabulky schématu \texttt{Ordering} neobsahuje. Zdroj fakta uvádí sám, není to chyba \\
& \emph{Convention} (3) & dvakrát \texttt{IDENTITY} místo \texttt{Auto} (rozh.~187) a jméno třídy klíče \\
\bottomrule
\end{tabularx}%
\caption{Příklad 2: překlad označených konstrukcí a záznamy.}
\label{tab:ex2}
\end{table}

\textbf{Výsledek.} Šestnáct artefaktů: čtyři entity a ke každému dotazu metoda a samotný text JPQL. Všech 72 faktů má stav \emph{carried}.

\subsubsection*{Příklad 3 (pokročilý): EclipseLink → EF~Core, rozsáhlý dotaz}

\textbf{Cíl.} Ověřit chování nástroje na rozsáhlém dotazu. \textbf{Předpoklady.} Znalost JPQL a LINQ. \textbf{Vstup.} Příklad \emph{Marketplace: one long query with seven joins and fourteen subqueries} není výukový. Je to zátěžový a regresní případ: jeden dotaz kombinuje sedm joinů, poddotazy zanořené do pěti úrovní, kvantifikované porovnání, \texttt{like} s~escape znakem, seskupení s~\texttt{having} a stránkování. Nástroj ho přeloží do jazyka cíle bez záznamu \emph{Fallback} a \emph{Failure}. Že převod příkladu nemá záznam \emph{Failure}, ověřuje test \texttt{ExampleCatalogTest}. Experimentální běh ověřuje artefakty příkladu do 3.~stupně, shodu řádků u~něj neměří (tab.~\ref{tab:vstupy}).

\textbf{Co dotaz počítá.} Pro každou dvojici regionu a zákazníka vrací ukazatele prodeje: počet objednávek a řádků, počet kusů, hrubou tržbu, poskytnutou slevu, vrácenou částku, průměrnou cenu a další. Počítá jen odeslané a doručené objednávky ze zadaného období, vynechává zákazníky s~testovacím e-mailem a uplatňuje další podmínky s~poddotazy. Skupiny filtruje podle tržby a počtu objednávek, řadí podle tržby a stránkuje.

\textbf{Kroky překladu.} Deset entit JPA vyjde jako jedenáct entit EF~Core: entitu \texttt{ProductTag} nástroj syntetizoval z~\texttt{@ManyToMany} (rozh.~005). Třída \texttt{SalesDashboard} s~dotazem (výpis~\ref{lst:ex3-in}) vyjde jako jedna metoda v~LINQ (výpis~\ref{lst:ex3-out}). Tab.~\ref{tab:ex3} uvádí překlad jednotlivých konstrukcí a místa, kde je zápis v~cíli delší než ve zdroji.

\begin{lstlisting}[language=Java,caption={Příklad 3, vstup (výňatek): \texttt{SalesDashboard.java}.},label={lst:ex3-in}]
return em.createQuery("""
        select r.Name as Region, c.CustomerId as CustomerId, c.DisplayName as Customer,
               count(distinct o.OrderId) as Orders, count(ol) as Lines, sum(ol.Quantity) as Units,
               sum(ol.Quantity * ol.UnitPrice) as GrossRevenue,
               sum(case when ol.Discount is null then 0 else ol.Quantity * ol.UnitPrice * ol.Discount end) as DiscountGiven,
               coalesce(sum(rf.Amount), 0) as Refunded, avg(ol.UnitPrice) as AverageUnitPrice,
               ... ,
               (select count(rv3) from Review rv3 where rv3.CustomerId = c.CustomerId) as ReviewsWritten
        from OrderLine ol
        join SalesOrder o on o.OrderId = ol.OrderId
        join Customer c on c.CustomerId = o.CustomerId
        join Region r on r.RegionId = c.RegionId
        join Product p on p.ProductId = ol.ProductId
        join Category cat on cat.CategoryId = p.CategoryId
        join Supplier s on s.SupplierId = p.SupplierId
        left join Refund rf on rf.OrderId = ol.OrderId and rf.LineNumber = ol.LineNumber
        where o.Status in ('SHIPPED', 'DELIVERED')
          and o.OrderedAt between :fromDate and :toDate
          and ...
          and lower(c.Email) not like '%!_test@%' escape '!'
          and cat.CategoryId in (select leaf.CategoryId from Category leaf
                where leaf.ParentCategoryId in (select branch.CategoryId from Category branch
                      where branch.ParentCategoryId is null and exists (select p2.ProductId from Product p2
                            where p2.CategoryId = leaf.CategoryId and p2.ListPrice > (select avg(p3.ListPrice)
                                  from Product p3
                                  where p3.SupplierId in (select s2.SupplierId from Supplier s2
                                        where s2.RegionId = c.RegionId and s2.Rating >= 4)))))
          and ...
          and not exists (select rf2.RefundId from Refund rf2 join SalesOrder o2 on o2.OrderId = rf2.OrderId
                where o2.CustomerId = c.CustomerId and rf2.Amount > (select max(ol2.UnitPrice) from OrderLine ol2
                      where ol2.OrderId = o2.OrderId))
          and ol.UnitPrice > any (select p4.ListPrice from Product p4
                where p4.CategoryId = p.CategoryId and p4.Discontinued = true)
        group by r.Name, c.CustomerId, c.DisplayName
        having sum(ol.Quantity * ol.UnitPrice) > :minRevenue and count(distinct o.OrderId) >= :minOrders
           and avg(ol.UnitPrice) >= (select avg(p5.ListPrice) from Product p5 where ...)
        order by GrossRevenue desc, r.Name asc, c.CustomerId asc
        """, Tuple.class)
    .setParameter("fromDate", fromDate) ... .setFirstResult(skip).setMaxResults(take).getResultList();
\end{lstlisting}

\begin{lstlisting}[language={[Sharp]C},caption={Příklad 3, výstup (výňatek): metoda dotazu v~LINQ pro EF~Core.},label={lst:ex3-out}]
public static IQueryable Query(DbContext ctx, DateTime fromDate, DateTime toDate, short minRating,
                               decimal minRevenue, long minOrders, int skip, int take)
{
    return ctx.Set<OrderLine>()
        .Join(ctx.Set<SalesOrder>(), ol => ol.OrderId, o => o.OrderId, (ol, o) => new { ol, o })
        .Join(ctx.Set<Customer>(), t => t.o.CustomerId, c => c.CustomerId, (t, c) => new { t.ol, t.o, c })
        ...
        .LeftJoin(ctx.Set<Refund>(), t => new { t.ol.OrderId, t.ol.LineNumber },
                  rf => new { rf.OrderId, rf.LineNumber }, (t, rf) => new { t.ol, t.o, t.c, t.r, t.p, t.cat, t.s, rf })
        .Where(t => new[] { "SHIPPED", "DELIVERED" }.Contains(t.o.Status)
            && t.o.OrderedAt >= fromDate && t.o.OrderedAt <= toDate && ...
            && !(t.c.Email.ToLower().Contains("_test@"))
            && ctx.Set<Category>().Where(leaf => ctx.Set<Category>().Where(branch => branch.ParentCategoryId == null
                 && ctx.Set<Product>().Where(p2 => p2.CategoryId == leaf.CategoryId && p2.ListPrice > ctx.Set<Product>()
                      .Where(p3 => ctx.Set<Supplier>().Where(s2 => s2.RegionId == t.c.RegionId && s2.Rating >= 4)
                           .Select(s2 => s2.SupplierId).Contains(p3.SupplierId)).Average(p3 => p3.ListPrice)).Any())
                 .Select(branch => branch.CategoryId).Contains(leaf.ParentCategoryId.Value))
                 .Select(leaf => leaf.CategoryId).Contains(t.cat.CategoryId)
            && ...
            && ctx.Set<Product>().Where(p4 => p4.CategoryId == t.p.CategoryId && p4.Discontinued == true)
                 .Select(p4 => p4.ListPrice).Any(v => t.ol.UnitPrice > v))
        .GroupBy(t => new { t.r.Name, t.c.CustomerId, t.c.DisplayName })
        .Where(g => g.Sum(t => t.ol.Quantity * t.ol.UnitPrice) > minRevenue
            && g.Select(t => t.o.OrderId).Distinct().Count() >= minOrders && ...)
        .Select(g => new { Region = g.Key.Name, CustomerId = g.Key.CustomerId, Customer = g.Key.DisplayName,
            Orders = g.Select(t => t.o.OrderId).Distinct().Count(), Lines = g.Count(), ...,
            GrossRevenue = g.Sum(t => t.ol.Quantity * t.ol.UnitPrice),
            DiscountGiven = (g.Count(t => (t.ol.Discount == null ? 0
                    : t.ol.Quantity * t.ol.UnitPrice * t.ol.Discount) != null) > 0
                ? (decimal?)g.Sum(t => (t.ol.Discount == null ? 0
                    : t.ol.Quantity * t.ol.UnitPrice * t.ol.Discount)) : null),
            Refunded = g.Sum(t => t.rf.Amount), ...,
            ReviewsWritten = ctx.Set<Review>().Where(rv3 => rv3.CustomerId == g.Key.CustomerId).Count() })
        .OrderByDescending(p => p.GrossRevenue).ThenBy(p => p.Region).ThenBy(p => p.CustomerId)
        .Skip(skip).Take(take);
}
\end{lstlisting}

\begin{table}[H]
\centering\footnotesize
\begin{tabularx}{\textwidth}{L{4.6cm} Y}
\toprule
\textbf{JPQL} & \textbf{LINQ} \\
\midrule
joiny & řetěz \texttt{Join} s~anonymními typy, \texttt{LeftJoin} nad dvojicí sloupců \\
\texttt{between} & dvojice porovnání (pravidlo Q14 článku~\cite{ormorpher-jss}) \\
\texttt{like \ldots escape} & \texttt{Contains} (rozh.~051, 102) \\
\texttt{> any} & \texttt{Any(v => \ldots)} (rozh.~119) \\
\texttt{count(distinct)} & \texttt{Select(\ldots).Distinct().Count()} \\
\texttt{sum} nad skupinou bez hodnoty & podmíněný výraz, který vrací \texttt{null}, protože \texttt{Sum} EF~Core by vrátil nulu (rozh.~155). Zápis je delší, výsledek je shodný \\
\midrule
\textbf{Záznam} & \textbf{Za co} \\
\midrule
\emph{Incompleteness} (10) & entity, které katalog instance neobsahuje \\
\emph{Loss} (13) & inicializátory kolekcí, sloupec \texttt{CHAR(2)}, verze, strategie klíčů, které EF~Core anotacemi nevyjádří \\
\emph{Convention} (6) & klíče, které EF~Core generuje sám, třída \texttt{SalesDashboard}, která do výstupu nejde, spojovací entita a \texttt{between} \\
\bottomrule
\end{tabularx}%
\caption{Příklad 3: překlad konstrukcí dotazu a záznamy. Tabulka faktů má 186 řádků, 178 se stavem \emph{carried} a 8 se stavem \emph{dropped}.}
\label{tab:ex3}
\end{table}

\subsection{Další stránky nástroje}\label{sec:stranky}

Běžný převod potřebuje jen stránku \emph{Translation}. Ostatní stránky z~menu mají tento účel:

\begin{itemize}[leftmargin=*,itemsep=1pt]
\item \textbf{Tutorial} je výuková stránka s~prvním převodem v~pěti krocích (kap.~\ref{sec:tutorial}).
\item \textbf{Examples} je výkladová stránka (obr.~\ref{fig:examples}). Obsahuje schéma pipeline převodu, základní vstup (týž zákazník zapsaný v~každém ze šesti frameworků) a 22 popsaných příkladů ve skupinách. U~každého příkladu ukazuje skutečný vstup, výstup a záznamy, které vrátila tato instance. Každý příklad je také v~seznamu \emph{Start from an example} překladové obrazovky: načte se a převede kroky 2 a~3 z~kap.~\ref{sec:tutorial}. Doporučený první příklad je Dapper~→~EF~Core. Po načtení lze cíl přepnout na kterýkoli ze šesti frameworků.
\item \textbf{LDBC, FIBEN a SEDE} jsou výzkumné katalogy cizích dotazů, na kterých se nástroj ověřuje (kap.~\ref{sec:exp}). K~běžnému převodu nejsou potřeba. Stránka katalogu ukazuje řádek za dotaz a sloupec za cíl (obr.~\ref{fig:katalog}). Jméno jazyka (LINQ, HQL, JPQL, SQL) znamená, že dotaz vyšel v~jazyce cíle. Štítek \emph{native SQL} znamená fallback, štítek \emph{refused} odmítnutí s~důvodem. Stránka slouží ke zjištění, zda nástroj daný tvar dotazu do daného cíle přeloží. Každý dotaz lze na stránce převést.
\item \textbf{Advisor} je experimentální a je vyňatý ze záruk nástroje (tab.~\ref{tab:vynate}). Pro zadané entity a zátěž dotazů měří kandidáty a doporučí framework. Měří jen Dapper a EF~Core, funguje jen v~instanci v~Dockeru a jeho funkci žádný test neověřuje.
\end{itemize}

\snimek[0.8]{10-examples}{Výkladová stránka \emph{Examples}: schéma pipeline, základní vstup a příklady ve skupinách.}{fig:examples}

\snimek{11-katalog}{Stránka katalogu LDBC: řádek za dotaz, sloupec za cíl. Text bez štítku je jazyk cíle, štítek \emph{native SQL} je fallback.}{fig:katalog}

\FloatBarrier
\subsection{Podporované vstupní a výstupní konstrukce}\label{sec:co-jak}

Převod má tři fáze. Parser zdrojového frameworku přečte jednotky do mezireprezentace nezávislé na frameworku. Volitelná fáze doplní z~katalogu databáze fakta, která zdroj neuvádí. Builder cílového frameworku z~mezireprezentace zapíše artefakty. Konstrukce, kterou mezireprezentace nenese nebo cíl nevyjádří, vede k~záznamu (kap.~\ref{sec:limity}). Tab.~\ref{tab:frameworky} uvádí, co každý framework čte a píše, tab.~\ref{tab:slovnik}, co mezireprezentace nese.

\begin{table}[H]
\centering\footnotesize
\begin{tabularx}{\textwidth}{L{1.75cm} Y L{5.0cm} L{4.6cm}}
\toprule
\textbf{Framework} & \textbf{Entity a mapování} & \textbf{Dotazy čte z} & \textbf{Dotazy píše jako} \\
\midrule
Dapper & třída C\# & SQL, SQL ve voláních Dapperu v~kódu & SQL s~metodou nad \texttt{IDbConnection} \\
EF~Core & třída C\# s~anotacemi. Fluent konfigurace v~\texttt{OnModelCreating} a \texttt{IEntityTypeConfiguration<T>} jako vymezená podmnožina & LINQ, nativní SQL v~kódu & LINQ \\
NHibernate & třída C\# a \texttt{hbm.xml} & LINQ, HQL jako samostatná jednotka, nativní SQL a HQL v~kódu & HQL a vedle něj LINQ nad \texttt{session.Query<T>()} \\
Hibernate & anotace \texttt{jakarta.persistence} (i~\texttt{orm.xml}), profil Hibernatu & JPQL a HQL Hibernatu, nativní SQL v~kódu & JPQL, HQL tam, kde JPQL tvar nemá (rekurze) \\
EclipseLink & totéž, profil EclipseLinku & JPQL, nativní SQL v~kódu & JPQL \\
MyBatis & doménová třída a \texttt{<resultMap>} & dokument mapperu (SQL), \texttt{@Select} na rozhraní & dokument mapperu a metoda rozhraní mapperu \\
\bottomrule
\end{tabularx}%
\caption{Co každý framework čte a píše.}
\label{tab:frameworky}
\end{table}

\begin{table}[H]
\centering\footnotesize
\begin{tabularx}{\textwidth}{L{3cm} Y}
\toprule
\textbf{Oblast} & \textbf{Konstrukce} \\
\midrule
projekce & sloupce, \texttt{DISTINCT}, konstanta pod aliasem, okenní funkce, agregace do seznamu (\texttt{STRING\_AGG}) \\
filtrace & podmínkový strom s~AND, OR a NOT, poddotazy \texttt{IN}, \texttt{EXISTS}, skalární a kvantifikované porovnání \\
joiny & vnitřní, levý, pravý, plný a křížový, join s~podmínkou nad rámec rovností \\
seskupení a řazení & agregace, seskupení i~podle výrazu, \texttt{HAVING}, řazení i~podle agregátu, stránkování \\
výrazy & aritmetika, konkatenace, uzavřený slovník funkcí, \texttt{CASE} \\
parametry & pojmenované, poziční, zachycené hodnoty C\#, kolekce za \texttt{IN} \\
skládání dotazů & množinové operace, CTE (\texttt{WITH}), odvozená tabulka, rekurzivní CTE \\
\midrule
mapování & složené primární klíče se strategií generování pro každou část (rozh.~006), vícesloupcové cizí klíče, vztahy 1:1, 1:N a N:M (ten jako explicitní spojovací entita, rozh.~005), unikátní omezení, verze \\
\bottomrule
\end{tabularx}%
\caption{Co mezireprezentace nese: slovník dotazů a mapování.}
\label{tab:slovnik}
\end{table}

\textbf{Doplnění z~katalogu} (F4--F6). Když zdroj neuvádí tabulku, typ, délku, klíč nebo cizí klíč, převod načte fakt z~katalogu připojené databáze a vznikne záznam \emph{Supplied}. V~případě rozporu platí hodnota zdroje (rozh.~015, 017) a vznikne záznam \emph{Conflict}.

\textbf{Fallback.} Co jazyk cíle nevyjádří, napíší EF~Core, NHibernate, Hibernate a EclipseLink celým dotazem v~nativním SQL přes API frameworku, vždy se záznamem \emph{Fallback} (rozh.~113, tab.~\ref{tab:fallback}). Dapper a MyBatis fallback nepotřebují.

\begin{table}[H]
\centering\footnotesize
\begin{tabular}{l l}
\toprule
\textbf{Konstrukce} & \textbf{Cíle, které ji píší nativním SQL} \\
\midrule
množinové operace & NHibernate \\
plný vnější join, CTE, agregace do seznamu & NHibernate, EclipseLink \\
rekurze, okenní funkce & EF~Core, NHibernate, EclipseLink \\
agregát nad celým výsledkem & EF~Core \\
\bottomrule
\end{tabular}
\caption{Příklady konstrukcí zapsaných fallbackem.}
\label{tab:fallback}
\end{table}

\textbf{Hibernate a EclipseLink} jsou dva profily nad jednou vrstvou Jakarta Persistence. Čtou a píší tytéž anotace a liší se výchozími hodnotami a chováním. Výpis~\ref{lst:profily} ukazuje rozdíl na jedné entitě: anotaci \texttt{@Nationalized} EclipseLink nemá, a proto do něj převod zapíše typ sloupce. Generátor klíče bez strategie znamená v~Hibernatu sekvenci \texttt{Tags\_SEQ}, a převod ji do EclipseLinku zapíše výslovně, protože tatáž anotace by tam znamenala generátor \texttt{TABLE}. Hlavní rozdíly uvádí tab.~\ref{tab:profily}, úplný výčet \texttt{docs/subset.md}.

\noindent\begin{minipage}{\linewidth}
\begin{lstlisting}[language=Java,caption={Rozdíl profilů: vstup pro Hibernate a výstup pro EclipseLink (výběr).},label={lst:profily}]
// vstup (Hibernate)
@Id @GeneratedValue
private Integer id;
@Nationalized @Column(name = "Label", length = 40)
private String label;

// výstup (EclipseLink)
@Id
@GeneratedValue(strategy = GenerationType.SEQUENCE, generator = "Tag_id_gen")
@SequenceGenerator(name = "Tag_id_gen", sequenceName = "Tags_SEQ", allocationSize = 50)
@Column(name = "id")
private Integer id;
@Column(name = "Label", length = 40, columnDefinition = "nvarchar(40)")
private String label;
\end{lstlisting}
\end{minipage}

\begin{table}[H]
\centering\footnotesize
\begin{tabularx}{\textwidth}{L{2.8cm} L{3.9cm} L{2.2cm} Y}
\toprule
\textbf{Rozdíl} & \textbf{Hibernate} & \textbf{EclipseLink} & \textbf{Co s~ním převod dělá} \\
\midrule
\texttt{@GeneratedValue} bez strategie (\texttt{AUTO}) & sekvence \texttt{<tabulka>\_SEQ} & generátor \texttt{TABLE}: tabulka \texttt{SEQUENCE} s~řádkem \texttt{SEQ\_GEN} & mechanismus zdroje přenáší do cíle. \emph{Auto} z~EF~Core a~NHibernatu píše do JPA jako \texttt{IDENTITY} (rozh.~187) \\
nationalizace textu & anotace \texttt{@Nationalized} & anotaci nemá & do EclipseLinku píše \texttt{columnDefinition} (rozh.~080) \\
hodnota \texttt{CHAR(n)} & předá ji doplněnou mezerami, jak ji vrací SQL Server & koncové mezery odřízne & v~každém cíli zachová hodnotu, jakou vrací SQL (rozh.~169) \\
\texttt{@Version} nad \texttt{byte[]} & přijímá & odmítá & do EclipseLinku píše sloupec jen pro čtení (rozh.~179) \\
plný vnější join & vyjádří & nevyjádří & do EclipseLinku píše dotaz nativním SQL (rozh.~180) \\
\bottomrule
\end{tabularx}%
\caption{Kde se profily Hibernate 7.4.5 a EclipseLink 5.0.0 liší.}
\label{tab:profily}
\end{table}

\subsection{Meze překladu}\label{sec:limity}

O~konstrukci, kterou nástroj nepřevede celou, platí jedno ze tří tvrzení. Uživatel je rozliší podle záznamu:

\begin{itemize}[leftmargin=*,itemsep=1pt]
\item \textbf{Nepodporováno.} Artefakt nevznikne a záznam \emph{Failure} konstrukci jmenuje. Patří sem každý dotaz, který by v~cíli vrátil jiné řádky (rozh.~053, 065, 070).
\item \textbf{Přeloženo se ztrátou.} Artefakt vznikne bez konstrukce a záznam \emph{Loss} ji jmenuje. Přeložený dotaz vrací tytéž řádky.
\item \textbf{Podporováno bez záruky.} Funkce existuje, záruky nástroje se na ni nevztahují a žádný test ji neověřuje. Platí pro Advisor.
\end{itemize}

Oblasti vyňaté ze záruk uvádí tab.~\ref{tab:vynate}. Tab.~\ref{tab:limity-priklady} ukazuje hlavní meze na krátkém vstupu s~výsledkem, který uživatel dostane, tab.~\ref{tab:limity} ostatní meze s~technickým důvodem. Úplný výčet konstrukci po konstrukci je v~\texttt{docs/subset.md}.

\begin{table}[H]
\centering\footnotesize
\begin{tabularx}{\textwidth}{L{5.2cm} L{3.2cm} Y}
\toprule
\textbf{Vyňatá oblast} & \textbf{Tvrzení} & \textbf{Co to znamená} \\
\midrule
Advisor a benchmarking & podporováno bez záruky & bez testu funkce, měří dva frameworky ze šesti, spouští cizí kód bez izolace \\
dědičnost, komponenty a spojené tabulky (\texttt{<subclass>}, \texttt{@Embedded}, \texttt{[Owned]}\ldots) & přeloženo se ztrátou & mezireprezentace nese jen plochou třídu. K~prvku vznikne záznam \emph{Loss} a jeho vlastnost ve výstupu chybí \\
jiný dialekt než SQL Server 2022 & nepodporováno & jediný dialekt nástroje \\
experimenty nad rámec T1--T3 & mimo rozsah & experimenty T4--T6 s~jazykovými modely a T7 s~Advisorem \\
\bottomrule
\end{tabularx}%
\caption{Vyňaté oblasti.}
\label{tab:vynate}
\end{table}

{\footnotesize
\begin{longtable}{P{0.25} P{0.32} P{0.43}}
\toprule
\textbf{Mez} & \textbf{Vstup} & \textbf{Výsledek} \\
\midrule
\endfirsthead
\toprule
\textbf{Mez} & \textbf{Vstup} & \textbf{Výsledek} \\
\midrule
\endhead
\rule{0pt}{26pt}\\
\endfoot
\bottomrule
\caption{Hlavní meze na krátkém vstupu a očekávaný diagnostický výsledek.}\label{tab:limity-priklady}\\
\endlastfoot
dědičnost & \texttt{public class Customer : Party}, třída \texttt{Party} ve vstupu není & \emph{Loss}: \emph{The class Customer lists Party first among its base types \ldots{} Customer is translated with the members it declares itself}. Entita vyjde jen se svými vlastnostmi \\
komponenta & \texttt{@Embedded private Address address;} & \emph{Loss}: vlastnost ve výstupu chybí \\
dynamické SQL MyBatisu & \texttt{<if test="name != null"{}>c.CustomerName = \#\{name\}</if>} & \emph{Failure}: \emph{The statement carries <if>, whose expansion an OGNL expression decides, so the statement is a family of statements rather than one \ldots{} no artifact was generated.} \\
funkce mimo slovník v~podmínce & \texttt{WHERE LEFT(c.CustomerName, 3) = \textquotesingle{}Tai\textquotesingle{}} & \emph{Failure}: \emph{The WHERE clause uses \textquotesingle{}LEFT(c.CustomerName, 3)\textquotesingle{}, which is a construct outside the vocabulary of expressions the query representation carries, and a query emitted without its filter would return different rows; no artifact was generated.} \\
funkce mimo slovník v~projekci & \texttt{SELECT LEFT(c.CustomerName, 3) AS Prefix, c.CreditLimit} & \emph{Loss}: výraz se z~projekce vypustí. Když v~projekci nezbude nic, je to \emph{Failure} \\
jiný dialekt SQL & deklarace \emph{Source SQL is another database system} & \emph{Failure} každého dotazu v~SQL, \emph{Loss} typu sloupce zapsaného v~SQL. Katalog fakta doplní a jeden záznam \emph{Conflict} uvede, že zdroj deklaruje jiný systém, než jaký katalog popisuje (rozh.~091) \\
zdroj bez klíče a bez katalogu & třída Dapperu nebo MyBatisu na instanci bez katalogu & do NHibernatu, Hibernatu a EclipseLinku \emph{Failure}, do EF~Core entita s~\texttt{[Keyless]} \\
neznámý jazykový typ & \texttt{private java.time.Instant createdAt;} & \emph{Incompleteness}: výstup opakuje jméno typu ze zdroje \\
\end{longtable}}

{\footnotesize
\begin{longtable}{P{0.36} P{0.36} P{0.28}}
\toprule
\textbf{Konstrukce} & \textbf{Důvod} & \textbf{Projev} \\
\midrule
\endfirsthead
\toprule
\textbf{Konstrukce} & \textbf{Důvod} & \textbf{Projev} \\
\midrule
\endhead
\rule{0pt}{26pt}\\
\endfoot
\bottomrule
\caption{Ostatní meze, jejich technický důvod a projev.}\label{tab:limity}\\
\endlastfoot
\texttt{fetch = LAZY} na referenci v~EclipseLinku & bez weavingu načte EclipseLink referenci hned & \emph{Loss} \\
hodnota \texttt{CHAR(n)} v~EclipseLinku & EclipseLink koncové mezery odřízne, ostatní frameworky ne (rozh.~169) & \emph{Convention} (artefakt mezery doplní nebo ořízne), jinak \emph{Loss} \\
tvar LINQ, který provider NHibernatu nevyjádří & HQL je závazný tvar, LINQ doplněk & \emph{Omitted} \\
okenní funkce mimo slovník (rámec okna, \texttt{LAG}, \texttt{NTILE}) & slovník obsahuje tři pořadové funkce a pět agregátů nad oknem (rozh.~161) & \emph{Loss} v~projekci (ze zdroje HQL \emph{Failure}), jinde \emph{Failure} \\
dotaz skládaný za běhu (\texttt{QueryOver}, \texttt{CreateCriteria}) & nemá text, který lze přečíst & \emph{Failure} \\
doplňkový balíček (Dapper.Contrib, MyBatis-Plus) & není součástí frameworku (rozh.~133) & \emph{Loss}. Tabulku a klíč doplní katalog \\
znak \texttt{?} (u~Hibernatu a NHibernatu i~\texttt{:} a \texttt{\textquotesingle}) ve jméně sloupce nebo tabulky v~hranatých závorkách, když dotaz jde fallbackem & API frameworku pro nativní SQL hledá v~textu značky parametrů a literály. Escape znak má jen Hibernate (měřeno, rozh.~217) & Hibernate píše znak za zpětným lomítkem. EclipseLink a NHibernate \emph{Failure} \\
kolekční parametr v~nativním SQL EF~Core. LINQ skládané nad nativním SQL & EF~Core~10 naváže interpolovanou kolekci jako jedinou hodnotu. LINQ nad nativním SQL spouští jako vlastní poddotaz a kroky nad ním by měnily množinu řádků (rozh.~113) & \emph{Failure} \\
dotaz volaný jménem (\texttt{createNamedQuery}, \texttt{GetNamedQuery}) & volání obsahuje jen jméno. Text dotazu je v~definici (\texttt{@NamedQuery}, \texttt{<query>}) a nástroj ho čte tam (rozh.~143, 081) & \emph{Incompleteness} u~volání. Dotaz vznikne z~jednotky s~definicí \\
příliš hluboký nebo dlouhý dotaz & limit zanoření (128) a délky textu dotazu (tab.~\ref{tab:konfigurace}) & \emph{Failure}, který limit jmenuje \\
\end{longtable}}

\subsection{Rozhraní API}\label{sec:api}

Převod je jedno volání \texttt{POST /orm/convert}. Výpis~\ref{lst:api-req} obsahuje úplné tělo požadavku pro převod jedné třídy z~Dapperu do EF~Core a příkaz, který ho odešle. Pole \texttt{sources} je seznam jednotek, každá s~názvem, jazykem a obsahem. Nepovinné pole \texttt{declaredSourceDialect} má hodnotu 10 pro SQL Server 2022 a 1000 pro jiný systém. Číselné hodnoty uvádí tab.~\ref{tab:api-vycty}.

\noindent\begin{minipage}{\linewidth}
\begin{lstlisting}[caption={Požadavek: soubor \texttt{request.json} a jeho odeslání.},label={lst:api-req}]
{
  "sourceOrm": 10,
  "targetOrm": 30,
  "sources": [
    {
      "name": "Customer.cs",
      "contentType": 90,
      "content": "public class Customer\n{\n    public int CustomerID { get; set; }\n    public string CustomerName { get; set; }\n}\n"
    }
  ]
}

curl -X POST http://localhost:5072/orm/convert -H "Content-Type: application/json" -d @request.json
\end{lstlisting}
\end{minipage}

Odpověď 200 (výpis~\ref{lst:api-res}) obsahuje záznam běhu, artefakty v~poli \texttt{sources}, záznamy v~poli \texttt{records} a tabulku faktů v~poli \texttt{facts}. Každý artefakt uvádí entitu a jednotky, ze kterých vznikl. Na instanci s~katalogem WideWorldImporters vrátí požadavek jeden artefakt, 12 záznamů a 9 faktů.

\begin{lstlisting}[caption={Odpověď na požadavek z~výpisu~\ref{lst:api-req} (zkráceno, vynechaná místa označují tři tečky).},label={lst:api-res}]
{
  "runId": "cf1e131e-c897-4f65-abac-efa50e8c7161",
  "toolVersion": "2.4.2",
  "sourceFramework": 10, "sourceFrameworkVersion": "2.1.79",
  "targetFramework": 30, "targetFrameworkVersion": "10.0.10",
  "targetDatabaseDialect": 10,
  "sources": [
    { "contentType": 10, "name": null, "entity": "Customer", "units": ["Customer.cs"],
      "content": "using Microsoft.EntityFrameworkCore;\r\n...[Table(\"Customers\", Schema = \"Sales\")]\r\n..." }
  ],
  "records": [
    { "kind": 5, "framework": 30, "artifact": null, "entity": "Customer", "property": null, "category": 1,
      "feature": null, "unit": null, "query": null,
      "reason": "Table name 'Customers' supplied by the database catalog." },
    ...
  ],
  "facts": [
    { "entity": "Customer", "property": null, "category": 1, "value": "Customers", "origin": 2, "unit": null,
      "outcome": 1, "source": null, "targets": [ { "artifact": 10, "start": 146, "length": 11 } ] },
    ...
  ],
  "catalogState": 2, "catalogReadMilliseconds": ..., "declaredSourceDialect": null,
  "maxNestingDepth": 128, "maxQueryTextLength": 1048576
}
\end{lstlisting}

\begin{table}[H]
\centering\footnotesize
\begin{tabularx}{\textwidth}{L{3.6cm} Y}
\toprule
\textbf{Pole} & \textbf{Hodnoty} \\
\midrule
\texttt{sourceOrm}, \texttt{targetOrm} & Dapper 10, NHibernate 20, EF~Core 30, Hibernate 40, EclipseLink 50, MyBatis 60 \\
\texttt{contentType} jednotky & XML 30, SQL 40, HQL 50, JPQL 80, soubor C\# 90, soubor Javy 100 \\
\texttt{contentType} artefaktu & entita C\# 10, metoda dotazu C\# 20, XML 30, SQL 40, HQL 50, entita Javy 60, metoda dotazu Javy 70, JPQL 80, dotaz LINQ 110 \\
\texttt{kind} záznamu & \emph{Failure} 1, \emph{Loss} 2, \emph{Convention} 3, \emph{Incompleteness} 4, \emph{Supplied} 5, \emph{Conflict} 6, \emph{Fallback} 7, \emph{Omitted} 8 \\
\texttt{catalogState} & \emph{not configured} 0, \emph{unused} 1, \emph{read} 2, \emph{unreachable} 3 \\
\texttt{origin} faktu & zdroj 1, katalog 2, odvozeno nástrojem 3 \\
\texttt{outcome} faktu & \emph{carried} 1, \emph{dropped} 2, \emph{not generated} 3 \\
\bottomrule
\end{tabularx}%
\caption{Číselné hodnoty v~požadavku a odpovědi \texttt{POST /orm/convert}.}
\label{tab:api-vycty}
\end{table}

\textbf{Chybové stavy} převodu jsou dva. Požadavek, který nelze zpracovat (neznámý framework, chybějící pole \texttt{sources}), vrátí odpověď 400 ve tvaru \texttt{ProblemDetails} (výpis~\ref{lst:api-err}). Tělo, které není platný JSON, vrátí prázdnou odpověď 400 a jiný typ obsahu odpověď 415. Vstup, který nelze přeložit, i~vada nástroje při překladu vrátí odpověď 200 se záznamem \emph{Failure} a s~tím, co se přeložilo (rozh.~197).

\noindent\begin{minipage}{\linewidth}
\begin{lstlisting}[caption={Odpověď 400 na požadavek s~\texttt{"sourceOrm": 99} a záznam \emph{Failure} v~odpovědi 200 na požadavek s~prázdným polem \texttt{sources}.},label={lst:api-err}]
{ "type": "https://tools.ietf.org/html/rfc9110#section-15.5.1", "title": "Bad Request", "status": 400,
  "detail": "Source ORM not supported" }

{ "kind": 1, "framework": 30, "artifact": null, "entity": null, "property": null, "category": null,
  "feature": null, "unit": null, "query": null,
  "reason": "The request carried no source unit with content, so nothing was generated." }
\end{lstlisting}
\end{minipage}

Úplný popis rozhraní vrací každá instance na adrese \path{/orm/openapi/v1.json}. Swagger UI na \texttt{/orm/swagger} běží jen ve vývojovém prostředí, tedy v~kontejneru systému a při \texttt{dotnet run}, ne v~publikované instanci.

\section{Návrh a implementace}\label{sec:denik}

Kapitola uvádí rozsah práce, vybraná rozhodnutí návrhu, vybraná selhání s~nápravou, odchylky od specifikace a využití AI. Rozhodnutí a selhání jsou výběr, ne úplný seznam. Úplné podklady jsou v~repozitáři: číslovaná rozhodnutí (\texttt{docs/decisions/}), záznamy běhů (\texttt{docs/runs/}), audity (\texttt{docs/audits/}) a anotace tagů.









% Návrh oddílu 4.1 „Architektura" pro docs/report/report.tex (stav tagu 2.4.2).
%
% Kam: jako první oddíl kapitoly 4, před \subsection{Rozsah a postup práce}.
%      Dosavadní oddíly 4.1 až 4.5 se tím posunou na 4.2 až 4.6.
% Co používá: makra a styly preambule zprávy (box, hl, arr, lbl, sloupce L a Y, \path),
%      popisky tab:projekty, tab:zajimavosti, sec:metodika a sec:limity.
% Po vložení upravit ve zprávě:
%   - odstavec „Struktura dokumentu" (1.4) a úvodní větu kapitoly 4 doplnit o architekturu,
%   - v kap. 2.1 u odstavce „Převod prochází pěti moduly" odkázat na kap. \ref{sec:architektura},
%   - připnuté verze šesti frameworků jsou i v oddíle „Druh joinu u navigace přes referenci",
%     jedno z obou míst nahradit odkazem.
% Zdroje faktů: docs/architecture.md (§1 až §4, §7), rozh. 009, 017, 023, 032, 039, 076, 082, 113, 153.

\subsection{Architektura}\label{sec:architektura}

Převod má tři fáze. Parser zdrojového frameworku převede vstup do mezireprezentace, fáze doplnění do ní přidá fakta z~katalogu databáze a builder cílového frameworku z~ní zapíše artefakty. Všech 36 směrů mezi šesti frameworky tak pokrývají parsery a buildery šesti wrapperů, ne překladač pro každou dvojici. Odpovědnosti projektů solution uvádí tab.~\ref{tab:projekty}, jejich závislosti obr.~\ref{fig:vrstvy}.

\begin{figure}[htbp]
\centering
\begin{tikzpicture}
\node[box, text width=52mm] (api) at (0,0) {\texttt{ORMConvertorAPI}\\REST API a statický frontend};
\node[box, text width=52mm] (orch) at (0,-1.5) {\texttt{OrmConvertor}\\orchestrace a factory};
\node[box, text width=40mm] (cat) at (5.6,-1.5) {\texttt{DatabaseCatalog}\\čtení katalogu, doplnění faktů};
\node[hl, text width=62mm] (net) at (-3.7,-3.2) {wrappery .NET\\Dapper, EF~Core, NHibernate};
\node[hl, text width=62mm] (jav) at (3.7,-3.2) {wrappery Javy\\Hibernate, EclipseLink, MyBatis};
\node[box, text width=44mm] (cs) at (-5.0,-5.0) {\texttt{CSharpEntityParsing}\\\texttt{LinqParsing}, \texttt{LinqBuilding}};
\node[box, text width=30mm] (sql) at (0,-5.0) {\texttt{TransactSql}\\T-SQL nad ScriptDom};
\node[box, text width=44mm] (jpa) at (5.0,-5.0) {\texttt{JavaEntityParsing}\\\texttt{JakartaPersistence}};
\node[box, text width=120mm] (aw) at (0,-6.7) {\texttt{AbstractWrappers}\\rozhraní parserů, abstraktní buildery, deskriptor cíle, typy záznamů};
\node[box, text width=120mm] (mod) at (0,-8.2) {\texttt{Common}, \texttt{Model}\\sdílené převodníky, čtení a kompilace C\#, mezireprezentace};
\draw[arr] (api) -- (orch);
\draw[arr] (orch) -- (cat);
\draw[arr] (orch.south) -- ++(0,-3mm) -| (net.north);
\draw[arr] (orch.south) -- ++(0,-3mm) -| (jav.north);
\draw[arr] (net.south) -- ++(0,-3mm) -| (cs.north);
\draw[arr] (net.south) -- ++(0,-3mm) -| ([xshift=-6mm]sql.north);
\draw[arr] (jav.south) -- ++(0,-3mm) -| (jpa.north);
\draw[arr] (jav.south) -- ++(0,-3mm) -| ([xshift=6mm]sql.north);
\draw[arr] (cs.south) -- (cs.south |- aw.north);
\draw[arr] (sql.south) -- (aw.north);
\draw[arr] (jpa.south) -- (jpa.south |- aw.north);
\draw[arr] (cat.east) -- ++(4mm,0) |- (aw.east);
\draw[arr] (aw) -- (mod);
\end{tikzpicture}
\caption{Závislosti projektů solution. Šipka vede od projektu k~projektu, který používá. Projekty Advisoru, vzorků a testů schéma neuvádí.}
\label{fig:vrstvy}
\end{figure}

\textbf{Vrstvy.} Orchestrace vybere parser, builder a deskriptor cíle přes factory a wrappery používá jen přes rozhraní z~\path{AbstractWrappers}. Wrapper nereferencuje jiný wrapper ani \path{DatabaseCatalog}. Co mají frameworky společné, je ve sdílených vrstvách jazyka: čtení entit v~C\#, čtení a zápis LINQ, čtení a zápis T-SQL, lexer a parser podmnožiny Javy a vrstva Jakarta Persistence (anotace, \texttt{orm.xml}, JPQL). Hibernate a EclipseLink jsou profily nad vrstvou JPA a liší se výchozími hodnotami implementace: jmény sloupců, generátorem za strategií \texttt{AUTO} a zápisem typu Unicode. MyBatis vrstvu frameworku nemá a stojí na vrstvách jazyka Java a T-SQL. Vrstvu T-SQL používají i~ostatní cíle pro fallback na nativní SQL.

\textbf{Mezireprezentace} (projekt \texttt{Model}) má dvě části, které uvádí tab.~\ref{tab:model}. Na frameworku nezávisí: typ sloupce je jedna z~dvaceti typových rodin, jazykový typ je skalár, reference, kolekce nebo neznámý typ a strategie klíče jmenuje mechanismus, ne generátor frameworku. Fakt, který zdroj neuvedl, je prázdný. Podle toho fáze doplnění pozná, co má vzít z~katalogu, a kontrola úplnosti pozná, co cíli chybí. Model nevaliduje (tab.~\ref{tab:zajimavosti}).

\begin{table}[H]
\centering\footnotesize
\begin{tabularx}{\textwidth}{L{3.1cm} Y}
\toprule
\textbf{Typ} & \textbf{Obsah} \\
\midrule
\multicolumn{2}{l}{\emph{Entity a mapování}} \\
\texttt{EntityMap} & entita s~vlastnostmi a jejich jazykovými typy, tabulka, schéma a fakta z~dalších řádků \\
\texttt{PropertyMap} & sloupec vlastnosti: jméno, typová rodina, Unicode, délka, přesnost, měřítko, zda smí být \texttt{NULL}, verze řádku, nepersistovaná vlastnost \\
\texttt{PrimaryKey} & části klíče v~pořadí, každá se strategií generování a jejími parametry \\
\texttt{Relation} & kardinalita, role (\emph{Owning} má cizí klíč, \emph{Inverse} je navigace bez sloupce), jména obou entit a uspořádané dvojice sloupců. Vztah N:M je spojovací entita se dvěma vztahy N:1 \\
\texttt{UniqueConstraint} & jméno a členy unikátního omezení \\
\midrule
\multicolumn{2}{l}{\emph{Dotazy}} \\
instrukce & seznam instrukcí dotazu: zdroj, join, filtr, seskupení, \texttt{HAVING}, řazení, projekce, stránkování, \texttt{DISTINCT}, množinová operace a poddotaz. CTE a odvozené tabulky jsou pojmenované definice vedle seznamu \\
\texttt{ConditionNode} & strom podmínek: porovnání, \texttt{AND} nebo \texttt{OR} nad seznamem uzlů a negace \\
\texttt{QueryOperand} & šest tvarů operandu: sloupec, konstanta, poddotaz, výčet hodnot, parametr a výraz (binární operace, funkce, \texttt{CASE}, okenní funkce) \\
\bottomrule
\end{tabularx}%
\caption{Datový model mezireprezentace.}
\label{tab:model}
\end{table}

\textbf{Čtení.} Parser do modelu nezapisuje přímo. Volá veřejné metody abstraktního builderu, které fakt zapíšou jen na prázdné místo (priorita zdrojů, rozh.~017) a odlišný údaj o~klíči zapíšou jako záznam \emph{Conflict}. Jednotku orchestrace nabídne čtení entit i~čtení dotazů, takže soubor smí obsahovat obojí.

\textbf{Zápis} má pevné pořadí kroků (návrhový vzor template method). U~entity mu předchází kontrola úplnosti: entita, které chybí fakt vyžadovaný cílem, artefakt nedostane a vznikne záznam \emph{Failure}. Sedm kroků pak zapíše importy, tabulku a schéma, klíč, vlastnosti, cizí klíče, členy vynucené frameworkem a výsledné soubory. Krok, pro který framework nic nezapisuje, je prázdný. Builder dotazu má osm kroků v~pořadí relačního vyhodnocení (zdroj, joiny, filtr, seskupení, \texttt{HAVING}, řazení, projekce, stránkování) a text skládá v~pořadí cíle. Konstrukci, kterou jazyk cíle nevyjádří, zapíše builder nativního SQL, má-li cíl API pro nativní SQL (rozh.~113). Jinak vznikne záznam \emph{Failure}.

\textbf{Deskriptor cíle} je jediné místo, kde wrapper uvádí vlastnosti frameworku: připnutou verzi, dialekt databáze, u~každé kategorie mapovacích faktů a každé konstrukce dotazu jednu ze tří hodnot (vyžaduje, vyjádří, nevyjádří), seznam funkcí a API pro nativní SQL. Z~deskriptoru vychází kontrola úplnosti, záznamy o~ztrátě a verze ve výsledku převodu (rozh.~009).

\textbf{Přidání frameworku} (S1) je nový projekt wrapperu:

\begin{enumerate}[leftmargin=*,itemsep=1pt]
\item projekt wrapperu v~solution a v~seznamu projektů v~Dockerfile,
\item deskriptor cíle,
\item parser entit nad sdílenou vrstvou jazyka nebo vlastní, který zapisuje jen přes veřejné metody abstraktního builderu,
\item builder entit se sedmi kroky,
\item volitelně parser a builder dotazů. Bez nich orchestrace jednotku s~dotazem odmítne záznamem,
\item registrace ve čtyřech factory orchestrace, hodnota výčtu frameworků a vzorky v~\texttt{SampleData},
\item testy ve složce frameworku a řádek ve vstupech matice směrů, u~javovského frameworku i~javovská sada.
\end{enumerate}

Kontrolu úplnosti, záznamy o~ztrátě, tabulku faktů, doplnění z~katalogu a přiřazení záznamů k~jednotkám dostane nový wrapper bez vlastního kódu. Potřeba, kterou kroky a deskriptor nepokrývají, mění sdílená rozhraní (\path{AbstractWrappers}) a vyžaduje číslované rozhodnutí (rozh.~153).

\textbf{Technologie} uvádí tab.~\ref{tab:technologie}. Verze jsou připnuté na jednom místě pro každý ekosystém (\path{Directory.Packages.props}, \texttt{pom.xml} javovské sady) a test hlídá, že se shodují s~verzí v~deskriptoru.

\begin{table}[H]
\centering\footnotesize
\begin{tabularx}{\textwidth}{L{4.4cm} L{4.2cm} Y}
\toprule
\textbf{Technologie} & \textbf{K~čemu} & \textbf{Důvod volby} \\
\midrule
.NET~10, C\# & celý nástroj včetně javovských wrapperů & platforma prototypu. Javovské wrappery v~C\# drží nasazení v~jednom procesu bez JVM (rozh.~076) \\
Roslyn 5.6.0 & čtení zdrojů v~C\#, kompilace ve 2.~stupni ověření & jedno čtení a jedna kompilace C\# pro všechny wrappery \\
ScriptDom 180.78.1 (gramatika T-SQL 160) & čtení a zápis T-SQL & jeden sdílený projekt pro Dapper, MyBatis, NHibernate a fallback (rozh.~082) \\
vlastní lexer Javy a parsery HQL a JPQL & čtení javovských zdrojů a dotazů & bez JVM a bez generátoru parserů. Neznámý tvar je záznam \emph{Failure} s~řádkem a sloupcem \\
ASP.NET Core, statické stránky, Pico CSS 2.1.1, highlight.js 11.12.0 & REST API a frontend & frontend bez frameworku a bez buildu: prohlížeč spouští soubory z~repozitáře (rozh.~032) \\
SQL Server 2022 & katalog, databáze testů a validátorů & jediný dialekt, pro který verze zapisuje \\
xUnit v3, JUnit~6.1.3, Maven 3.9.11, JDK~25 & .NET sada a javovská sada & javovská sada sestaví a spustí artefakty v~připnutém kontejneru (rozh.~076, 078) \\
Docker Compose & systém, databáze a obě sady & jeden soubor se dvěma profily, sada se spustí jedním příkazem (S5, rozh.~039) \\
GLPK přes P/Invoke & Advisor & převzato z~prototypu \\
\bottomrule
\end{tabularx}%
\caption{Použité technologie. Připnuté verze šesti frameworků: Dapper 2.1.79, EF~Core 10.0.10, NHibernate 5.7.0, Hibernate 7.4.5, EclipseLink 5.0.0 a MyBatis 3.5.19.}
\label{tab:technologie}
\end{table}











\subsection{Rozsah a postup práce}\label{sec:postup}

Rozsah práce uvádí tab.~\ref{tab:cisla}. Proti prototypu se architektura změnila v~pěti oblastech:

\begin{itemize}[leftmargin=*,itemsep=1pt]
\item \textbf{.NET část} (release 1.0 až 1.2.0): složené klíče, vztahy, doplnění faktů z~katalogu databáze, diagnostické záznamy, překlad dotazů mezi třemi frameworky a čtvrtý stupeň ověření.
\item \textbf{Javovský ekosystém} (2.0.0): Hibernate, EclipseLink a MyBatis jako wrappery nad sdílenými vrstvami, javovská testovací sada a překlad mezi ekosystémy oběma směry.
\item \textbf{Rozsah dotazů} (2.1.0): výrazy, CTE, rekurze a fallback na nativní SQL.
\item \textbf{Ověření cizími katalogy} (2.1.0 až 2.4.0): katalogy dotazů LDBC, FIBEN a SEDE, validační sada LDBC a tabulka faktů.
\item \textbf{Infrastruktura experimentů} (2.3.0): experimentální běh, běh mezí, případové studie a srovnání s~jazykovými modely.
\end{itemize}

\begin{table}[H]
\centering\footnotesize
\begin{tabular}{l l}
\toprule
\textbf{Údaj} & \textbf{Hodnota} \\
\midrule
začátek práce na forku & 14.~7.~2026, nad prototypem se dvěma desítkami testovacích metod \\
commity & asi 400 \\
kód mimo testovací projekt & asi 94~tisíc řádků C\#, z~toho asi 18~tisíc vzorky a katalogy v~\texttt{SampleData} \\
rozhodnutí, audity, release & 218, 9, 10 \\
testy ve verzi 1.2.0 (specifikace) & 615 \\
testy v~release \Vydani{} (23.~11.~2026) & 26~108 (.NET) a 16~962 (Java), obě sady bez selhání \\
z~toho mimo svůj běh přeskakuje & 13~737 a 13~661, převážně případy běhu SEDE \\
\bottomrule
\end{tabular}
\caption{Rozsah projektu.}
\label{tab:cisla}
\end{table}

Práce se řídila čtyřmi pravidly:

\begin{itemize}[leftmargin=*,itemsep=1pt]
\item \textbf{Rozhodnutí před implementací.} Změna, která volí mezi variantami, začíná souborem v~\texttt{docs/decisions/} s~jednou doporučenou variantou a zdůvodněním vůči požadavkům a článku~\cite{ormorpher-jss}. Kód následuje po něm.
\item \textbf{Aktualizace dokumentace.} Každá změna chování aktualizuje dokumentaci současného stavu.
\item \textbf{Audity.} Z~každého nálezu auditu plyne oprava, rozhodnutí nebo otevřená položka.
\item \textbf{Dva stroje.} Kontejnery a zaznamenané běhy sad, katalogů a experimentů běží na notebooku s~Dockerem (měřicí stroj). Běh na stolním počítači bez Dockeru je kontrola. Výjimky uvádí kap.~\ref{sec:exp}.
\end{itemize}

\subsection{Vybraná rozhodnutí návrhu}\label{sec:zajimavosti}

Projekt má 218 číslovaných rozhodnutí, jejich seznam je v~\path{docs/decisions/README.md}. Tab.~\ref{tab:zajimavosti} z~nich uvádí ta s~největším vlivem na návrh: o~mezireprezentaci, o~hranici odmítnutí dotazu, o~javovské části bez JVM a o~způsobu ověření. Čtyři další případy, u~nichž vadu nebo mez ukázalo až měření, jsou popsány jednotlivě.

{\footnotesize
\begin{longtable}{P{0.20} P{0.66} P{0.14}}
\toprule
\textbf{Téma} & \textbf{Rozhodnutí} & \textbf{Rozh.} \\
\midrule
\endfirsthead
\toprule
\textbf{Téma} & \textbf{Rozhodnutí} & \textbf{Rozh.} \\
\midrule
\endhead
\rule{0pt}{26pt}\\
\endfoot
\bottomrule
\caption{Vybraná rozhodnutí s~největším vlivem na návrh.}\label{tab:zajimavosti}\\
\endlastfoot
Mezireprezentace & vztah je uložen na entitě, entity se odkazují jménem, složený klíč je seřazený seznam částí, N:M je spojovací entita a podmínka dotazu je strom. Model nevaliduje, validace patří builderu cíle & 001, 002, 005, 006 \\
Hranice odmítnutí dotazu & hranicí je množina řádků: dotaz, který by ji změnil, se nevydá. Změna údajů o~řádku je ztráta se záznamem & 053, 065, 070 \\
Javovská část bez JVM & wrappery jsou v~C\#: lexer Javy, parser podmnožiny tříd, parser a builder JPQL. Nasazení je jeden proces. Přijetí frameworkem ověřují jen testy & 076, 077, 080, 084 \\
Javovská sada jako klient instance & sada bere artefakty přes HTTP, takže ověřuje výstup, který dostane uživatel API. Diferenční ověření porovnává řádky zdrojové a přeložené varianty s~jedním kanonickým výsledkem. Mutovaný artefakt musí skončit rozdílem & 078, 087, 089 \\
Fallback na nativní SQL & bez něj šlo do všech pěti cílů 22 ze 41 dotazů katalogu LDBC. S~ním vyjde dotaz i~v~cíli, jehož jazyk konstrukci nevyjádří & 113 \\
Měřicí prostředí & časy se měří v~zafixovaném prostředí (tab.~\ref{tab:rezimy}), krátký dotaz se opakuje a zaznamenává se medián. Validátor FIBEN vznikl z~běhu nástroje a měří regresi mezi verzemi, ne správnost dotazu vůči otázce & 145, 149, 165, 193 \\
Cizí artefakt & výstup modelu vstupuje jako adresář souborů s~manifestem a prochází týmiž stupni ověření. Model a jeho verzi určuje kanál, ne odpověď: Qwen se v~odpovědích označil za GPT-4 i~Claude & 172, 175, 204, 205 \\
Připnuté prostředí & lock file u~každého projektu, base image digestem, akce CI commitem. Generovaný zdroj ani záznam běhu neobsahuje přihlašovací údaje & 160, 186 \\
\end{longtable}}

\subsubsection*{Druh joinu u~navigace přes referenci}

\textbf{Problém.} Navigace přes referenci v~operandu LINQ, například \texttt{o.Customer.CustomerName}, je join, který v~textu dotazu není zapsán. Jeho druh určuje provider podle mapování. Čtení zdroje musí druh joinu zvolit tak, aby cíl vrátil tytéž řádky.

\textbf{Alternativy.} Rozh.~148 zamítlo \emph{vždy vnitřní join} (u~volitelné reference ve zdroji EF~Core by vynechal řádky) a \emph{vždy levý join}. Zvolilo hook provideru: EF~Core levý join pro cizí klíč, který smí být \texttt{NULL}, NHibernate vždy vnitřní. Tvrzení o~NHibernatu nebylo změřeno.

\textbf{Vada.} Měření z~18.~10. proti NHibernate 5.7.0 tvrzení vyvrátilo: provider vytvoří levý vnější join v~projekci, řazení, klíči seskupení i~v~testu na \texttt{NULL} a vnitřní jen ve filtru, který \texttt{NULL} splnit nemůže (výpis~\ref{lst:join-148}). Čtení podle rozh.~148 proto u~zdroje NHibernate vynechávalo řádky s~chybějící referencí. První tři stupně ověření tuto vadu neodhalí: dotaz se přeloží, zkompiluje a framework ho přijme. Liší se až počet řádků, a jen nad daty, v~nichž reference chybí. Totéž měření ukázalo druhou vadu u~obou providerů: v~řetězu \texttt{o.Customer.Region.RegionName} nad volitelným zákazníkem a povinným regionem následoval po levém joinu vnitřní, který vypustil řádky zachované levým.

\noindent\begin{minipage}{\linewidth}
\begin{lstlisting}[language={[Sharp]C},caption={Navigace přes volitelnou referenci: co vrací zdroj a co z~něj dělalo čtení podle rozh.~148.},label={lst:join-148}]
session.Query<Order>().Select(o => o.Customer.CustomerName)

// NHibernate 5.7.0:        ... left outer join Customers ...  objednávka bez zákazníka má NULL
// čtení podle rozh. 148:   vnitřní join                       objednávka bez zákazníka chybí
\end{lstlisting}
\end{minipage}

\textbf{Řešení.} Rozh.~151 zavádí jedno pravidlo pro EF~Core i~NHibernate: join je levý, smí-li řádek reference chybět, jinak vnitřní, a volitelnost se přenáší po celém řetězu. Zamítnuta byla varianta \emph{u~NHibernatu vždy levý join}, protože mezireprezentace téhož dotazu nad povinnou referencí by se mezi providery lišila bez rozdílu v~řádcích.

\textbf{Ověření.} Plán dotazu NHibernatu 5.7.0 čte test \texttt{NHibernateLinqProviderTest}. SQL, které vytvoří EF~Core 10.0.10, čte test \texttt{EFCoreNavigationJoinProviderTest}. Chování překladu ověřuje test \texttt{LinqReferencePathTest}. Z~případu plyne pracovní pravidlo: předpoklad o~chování frameworku se před rozhodnutím ověří testem proti připnuté verzi a test zůstává v~sadě, takže upgrade, který chování změní, v~něm selže. Na straně .NET jsou to testy \path{*ProviderTest} a \path{EFCoreQueryAcceptance} (SQL EF~Core přes \path{ToQueryString()}), na straně Javy testy \path{*ClaimsTest}. Připnuté verze jsou Dapper 2.1.79, EF~Core 10.0.10, NHibernate 5.7.0, Hibernate 7.4.5, EclipseLink 5.0.0 a MyBatis 3.5.19. Rozhodnutí: 148, 151.

\subsubsection*{\texttt{SUM} nad žádnou hodnotou}

\textbf{Problém.} \texttt{SUM} v~SQL, HQL i~JPQL je nad žádnou hodnotou \texttt{NULL}. \texttt{Enumerable.Sum} je nula a EF~Core tuto sémantiku zapisuje do SQL jako \texttt{COALESCE(SUM(x), 0)}. Sdílený builder LINQ psal \texttt{SUM} mezireprezentace jako \texttt{Sum}. Rozdíl se projeví při porovnání se součtem. Dotaz q207 katalogu FIBEN porovnává v~\texttt{HAVING} součet skupiny s~korelovaným součtem. Skupina, pro kterou vnitřní dotaz nemá řádek, se v~SQL porovnává s~\texttt{NULL} a z~výsledku vypadne. V~EF~Core se porovnává s~nulou a ve výsledku zůstane. Text katalogu vracel 1~řádek, artefakt EF~Core 211.

\textbf{Alternativy.} Rozh.~155 zamítlo uvedenou mez cíle EF~Core (dotaz by dál vracel jiné řádky), nativní SQL pro každý takový \texttt{SUM} (jazyk cíle má pro konstrukci zápis) a podmínku \texttt{Any()} (dva průchody a nepokrývá případ samých \texttt{NULL}).

\textbf{Řešení.} Cíl EF~Core píše \texttt{SUM} jako projekci nad jedinou skupinou seskupenou konstantou. Nad žádným řádkem skupina nevznikne a poddotaz vrátí \texttt{NULL}. Opačným směrem se \texttt{Sum} zdroje EF~Core čte jako \texttt{COALESCE(SUM(x), 0)}. Měření proti připnutým verzím ukázalo, že jde o~vlastnost jednoho agregátu jednoho provideru: \texttt{AVG}, \texttt{MIN} ani \texttt{MAX} EF~Core takto nezapisuje a provider NHibernatu píše \texttt{Sum} jako \texttt{sum} bez \texttt{COALESCE}. Úpravu proto provádí hook provideru EF~Core.

\textbf{Výkonnostní důsledek.} Forma končila voláním \texttt{FirstOrDefault()}, které EF~Core píše jako \texttt{TOP(1)}. Přes \texttt{TOP(1)} SQL Server korelovaný poddotaz nepřevede na join a spouští ho pro každou vnější skupinu. V~běhu FIBEN z~19.~10. mělo šest dotazů tohoto tvaru v~EF~Core medián 7,7--35~s, text katalogu a ostatní cíle 0,07--0,1~s. Podmínky měření: měřicí stroj, prostředí rozh.~145 se zapnutým dynamickým navýšením frekvence procesoru, strom \texttt{9f6082d} s~následujícím commitem, medián nejméně tří měřených spuštění po zahřátí. Příčinu ukázala kontrola z~21.~10. na druhém stroji nad SQL Serverem 2025 Express, který měřicím prostředím není: dotaz q207 četl tabulku \texttt{HOLDING} 214krát s~917~846 logickými čteními, text katalogu dvakrát s~8~578 čteními (\texttt{SET STATISTICS IO}). Plán dotazu záznam neobsahuje. Rozh.~159 formu zakončuje voláním \texttt{Max()}, které vrátí tytéž řádky bez \texttt{TOP}. Zamítlo ponechání formy, nativní SQL, podmínku \texttt{Any()} (v~téže kontrole 13,0~s) a ruční dekorelaci.

\textbf{Ověření.} Chování providerů ověřují \texttt{EFCoreSumProviderTest} a \texttt{NHibernateLinqProviderTest}, shodu řádků na čtvrtém stupni \texttt{SumOverNoValuesTest}. V~běhu FIBEN z~22.~10. na měřicím stroji mělo týchž šest buněk medián 63--122~ms. Že SQL Server formu s~\texttt{MAX} převede na join a formu s~\texttt{TOP(1)} ne, žádný test neověřuje. Ukáže to jen běh katalogu FIBEN. Rozhodnutí: 155, 159.

\subsubsection*{Doba validačního běhu LDBC}

\textbf{Problém.} Validační sada LDBC má 138~747 řádků, čtení a vkládání v~pevném pořadí, a ke každému čtení očekávané řádky. Sadu vydává LDBC a její očekávané řádky jsou na nástroji nezávislé. První dvě přehrání celé sady (12. a 14.~10., server bez omezení) trvala 20 a 25 hodin. Rozh.~145 pak zafixovalo prostředí pro měření času (\texttt{MAXDOP}~1, databáze na dvou procesorech) a validátor v~něm běžel také. Projekce z~18.~10. dala pro celou sadu 2,5--5 dnů.

\textbf{Alternativy.} Rozh.~164 zamítlo validaci v~měřicím prostředí (dny běhu kvůli vlastnosti, kterou validace nepotřebuje), přepínání serveru za běhu (režim by byl v~kódu testů a po pádu by server zůstal ve špatném režimu) a samostatnou databázovou službu (druhá kopie dat a druhý server). Rozh.~166 zamítlo sekvenční přehrání (o~čtvrtinu pomalejší), obě sady současně nad vlastními kopiemi dat (nevejdou se do paměti serveru) a souběžné čtení s~paralelními plány.

\textbf{Řešení.} Validátor neměří čas. Jeho veličinou je shoda řádků a ta na paralelismu serveru nezávisí. Server má proto dva režimy přepínané jedním souborem proměnných (tab.~\ref{tab:rezimy}). Pořadí sady je závazné jen vůči vkládáním: všechna čtení mezi dvěma vkládáními vidí týž stav. Validátor proto úsek mezi dvěma vkládáními čte ve všech frameworcích sady současně, každý framework ve vlastním vlákně.

\textbf{Ověření.} Nad prefixem 250 řádků (233 čtení, 17 vkládání, měřicí stroj, 22.~10.) kleslo přehrání .NET sady ze 7,2~min v~měřicím režimu (18.~10.) na 2,6--2,8~min. S~paralelními plány trvalo 2,9--3,2~min a každý framework četl o~10--15~\% déle, protože paralelní plán a souběžné dotazy sdílejí tatáž jádra. Celá sada pak proběhla za 29,6~h (.NET) a 27,2~h (Java) se shodou v~každém čtení (tab.~\ref{tab:meze}). Rozhodnutí: 110, 117, 164, 166.

\subsubsection*{Limity vstupu a meze cílů}

\textbf{Problém.} Přetečení zásobníku .NET nezachytí: proces skončí uprostřed požadavku a bez záznamu. U~služby, která čte cizí text, je to vada dostupnosti.

\textbf{Alternativy.} Rozh.~092 zamítlo větší zásobník (mez posouvá, neohraničuje) a přepis sestupu na explicitní zásobník (přepis tří parserů, který nepomůže parseru T-SQL ani Roslynu). Rozh.~191 zamítlo limit délky řetězu v~textu, protože by to byla heuristika pro každý jazyk zvlášť.

\textbf{Řešení.} Rozh.~092 před sestupem parseru počítá závorky a hloubku nad 128 odmítne záznamem \emph{Failure} s~řádkem a sloupcem. Testy ukázaly hloubku, která se v~textu neprojevuje závorkami: součet \texttt{a + b + c} je v~mezireprezentaci binární strom hluboký jako řetěz, stejně tak \texttt{CASE} v~\texttt{ELSE}, řada \texttt{NOT} a řetěz \texttt{UNION}. Rozh.~191 proto přidává tři části: hloubka mezireprezentace se měří po čtení, \texttt{CASE} a řada prefixových operátorů se počítají už v~textu a parser, kterému dochází zásobník, jednotku odmítne záznamem. Rozh.~194 omezuje délku textu dotazu na 1~048~576 znaků.

\textbf{Ověření.} Limity ověřují testy \path{NestingDepthCapTest} a \path{ModelDepthCapTest}. Meze mimo nástroj hledal běh mezí (rozh.~182, \path{LimitSearchTest}): syntetické řady rostoucí po jedné z~patnácti os z~každého zdroje, který řadu vyjádří, do všech šesti cílů, jedno měření na bod. Překladač zpracoval řadu po její poslední bod ve všech osách kromě tří (\texttt{joins}, \texttt{ctes}, \texttt{entities}), v~nichž ze zdroje EF~Core nebo Hibernate překročil mez S3. V~ostatních osách selhala dřív jiná vrstva (tab.~\ref{tab:vrstvy}). Na tyto meze nástroj nemá vliv, jen je uvádí (kap.~\ref{sec:limity}). Běh našel také dvě vady překladu do EclipseLinku, řetěz množinových operací a \texttt{CASE} v~\texttt{ELSE}, obě na druhém bodě řady. Řetěz množinových operací opravil commit bez rozhodnutí, \texttt{CASE} v~\texttt{ELSE} rozh.~201. Rozhodnutí: 092, 182, 191, 194, 201.

{\footnotesize
\begin{longtable}{P{0.20} P{0.38} P{0.24} P{0.18}}
\toprule
\textbf{Vrstva} & \textbf{Mez} & \textbf{Typ selhání} & \textbf{Osa běhu mezí} \\
\midrule
\endfirsthead
\toprule
\textbf{Vrstva} & \textbf{Mez} & \textbf{Typ selhání} & \textbf{Osa běhu mezí} \\
\midrule
\endhead
\rule{0pt}{26pt}\\
\endfoot
\bottomrule
\caption{Meze mimo nástroj podle vrstvy, běh mezí 29.~10.--1.~11.~2026.}\label{tab:vrstvy}\\
\endlastfoot
kompilátor C\# (Roslyn), cíl EF~Core & 16 úrovní vnořených poddotazů (4 úrovně 2--3~s, 8 úrovní 31--37~s) & kompilace neskončí do 30~min & \texttt{subqueries} \\
kompilátor C\# (Roslyn), cíl EF~Core & 2~048 vnořených \texttt{CASE}, 4~096 podmínek & chyba \texttt{CS8078} & \texttt{case}, \texttt{conditions} \\
kompilátor Javy (\texttt{javac}), cíle Hibernate a EclipseLink & text dotazu nad 65~535 bajtů: 4~096 podmínek, 2~048 \texttt{CASE}, 1~024 poddotazů & chyba \emph{constant string too long} & \texttt{conditions}, \texttt{case}, \texttt{subqueries} \\
EF~Core 10.0.10 & výraz o~4~096 členech & přetečení zásobníku v~generátoru SQL & \texttt{expression} \\
NHibernate 5.7.0 & 1~024 úrovní poddotazů & přetečení zásobníku v~parseru HQL & \texttt{subqueries} \\
Hibernate 7.4.5 & řetěz negací: 16 článků 5~s, 128 článků 8,5~min, 256 článků nad 30~min & přijetí neskončí do 30~min & \texttt{not} \\
Hibernate 7.4.5 & 2~048 vlastností entity & chyba \emph{Could not instantiate persister} & \texttt{properties} \\
EclipseLink 5.0.0 & 2~048 podmínek, 1~024 členů výrazu, 1~024 \texttt{CASE} & výjimka bez zprávy & \texttt{conditions}, \texttt{expression}, \texttt{case} \\
SQL Server 2022 & 64 úrovní poddotazů, 32 úrovní \texttt{EXISTS}, 256 joinů v~textu .NET cílů (512 joinů v~textu javovských cílů proběhlo za 27~min) & chyba 8623, plán nesestaven & \texttt{subqueries}, \texttt{exists}, \texttt{joins} \\
SQL Server 2022 & \texttt{CASE} vnořený víc než desetkrát & chyba 125 & \texttt{case} \\
běhové prostředí .NET & zásobník vlákna 1,5~MB na Windows, v~kontejneru několikanásobek. Týž vstup proto selže na hostiteli při menší hloubce & ukončení procesu bez záznamu & test \texttt{Nesting\-Depth\-CapTest} \\
\end{longtable}}

\subsection{Vybraná selhání a nápravná opatření}\label{sec:nepovedlo}

Tab.~\ref{tab:nepovedlo} uvádí vybraná selhání s~příčinou a nápravou. Další jsou v~nálezech auditů (\texttt{docs/audits/}) a v~záznamech běhů (\texttt{docs/runs/}). Přerušení validačního běhu, vyčerpání haldy a validita běhu SEDE ovlivňují platnost měření a jsou popsány jednotlivě.

{\footnotesize
\begin{longtable}{P{0.24} P{0.40} P{0.36}}
\toprule
\textbf{Selhání} & \textbf{Příčina} & \textbf{Náprava} \\
\midrule
\endfirsthead
\toprule
\textbf{Selhání} & \textbf{Příčina} & \textbf{Náprava} \\
\midrule
\endhead
\rule{0pt}{26pt}\\
\endfoot
\bottomrule
\caption{Vybraná selhání a nápravná opatření.}\label{tab:nepovedlo}\\
\endlastfoot
vnitřní join místo levého (rozh.~148) & nezměřené tvrzení o~provideru. Vadu našel audit, ne test & rozh.~151 a testy chování obou providerů (kap.~\ref{sec:zajimavosti}) \\
validační běh LDBC přerušený po 29~h (24.~10.) & navazující skript znovu vytvořil databázi pod běžícím validátorem & skript pracuje s~databází, až sám přečte stav kontejneru běhu jako \texttt{exited} \\
procesy ukončené pro nedostatek paměti (22.~10., 5.~11.) & systém a profil \texttt{test} běžely současně & před profilem \texttt{test} se systém zastaví. Součet mezí paměti všeho, co běží současně, se musí vejít do stroje (rozh.~193) \\
javovská část běhu mezí vyčerpala haldu (31.~10.) & \texttt{OutOfMemoryError} ukončil vlákno klienta a další body se chybně zapsaly jako \emph{crashed} & sada se při vyčerpané haldě ukončí a běh pokračuje od místa přerušení (rozh.~195) \\
první běh SEDE bez výsledku (8.~11.) & běh držel řádky všech buněk v~paměti do konce a pro 1~714 dotazů paměť nestačila & soubory běhu se zapisují po dotazech \\
validátor SEDE nevznikl & obě sady braly za zdroj jiný text & rozh.~209: zdrojem obou sad je text katalogu. Běh podle něj zaznamenán není \\
výňatky projektů se nezkompilovaly (24.~10.) & z~63 řádků první kontroly (21 příkladů × 3 cíle) došlo do třetího stupně 44. Výňatky odkazovaly na typy mimo výňatek, výčty a mapy & vstupy upraveny. Kontrola, která změnila vstupy experimentu, je zaznamenána jako běh v~\texttt{docs/runs/} \\
výňatek Cofoundry pod třetím stupněm: v~pěti cílech se nezkompiloval, v~NHibernatu ho framework nepřijal (5.~11.) & typ mimo výňatek a čtyři vady překladu: člen zachyceného jména čtený jako sloupec, vyhrazená jména v~\texttt{hbm.xml}, dvojí navigace v~EF~Core a tabulka \texttt{User} v~SQL Hibernatu bez oddělovačů & výňatek upraven, vady opraveny s~testem, poslední v~release \Vydani. Experimentální běh nad opraveným stromem zaznamenán není \\
\end{longtable}}

\subsubsection*{Přerušení validačního běhu LDBC}

Validační běh LDBC .NET sady byl 24.~10. přerušen po 29 hodinách u~řádku 128~816 ze 138~747, protože navazující skript znovu vytvořil databázi. Selhal skript, který řadí běhy za sebe, ne nástroj ani validátor.

\textbf{Příčina.} Příkaz \texttt{docker compose up -d} nad službou, která na databázi závisí, databázi znovu vytvoří, nedostane-li tytéž proměnné, s~nimiž byla spuštěna (rozh.~164). Pravidlo \path{--env-file judge.env} na každém příkazu bylo v~návodu. Navazující skript měl počkat na konec běhu. Naplánovaná úloha běhu byla předtím ukončena, ale její proces \texttt{docker wait} běžel dál a držel otevřený soubor, do něhož přesměrovával výstup. Příkaz \texttt{cmd} s~přesměrováním do souboru, který nejde otevřít, se neprovede a dávka pokračuje dalším řádkem. Dalším řádkem bylo \texttt{up -d mssql\_db} bez souboru proměnných. Databáze vznikla znovu v~měřicím režimu a běh skončil bez výsledku.

\textbf{Opatření.} Návrat čekacího příkazu není doklad, že běh skončil. Skript proto pracuje s~databází, až sám přečte stav kontejneru běhu jako \texttt{exited}, zapisuje do nových souborů a nejdřív ukončí proces \texttt{docker wait}, který po běhu zůstal.

\textbf{Ověření.} Samostatný test opatření neexistuje. Další validační běh nad celou sadou (2.--5.~11.) proběhl v~obou sadách bez přerušení (tab.~\ref{tab:meze}).

\subsubsection*{Vyčerpání haldy v~běhu mezí}

Javovská část běhu mezí překládá body nasazenou instancí přes HTTP. Selhala dvakrát: poprvé instance, podruhé sada.

\textbf{Selhání instance (29.~10.).} Instance několik hodin hlásila vyčerpání poolu vláken a pak ji jádro ukončilo pro nedostatek paměti. Klient ukončil čekání na překlad po dvou minutách, ale instance překládala dál, protože obsluha požadavku je synchronní a zrušení se k~ní nepřenese. Mez paměti instance neměla. Rozh.~195 proto dalo každé části vlastní mez: překlad 35~s, bod třicet minut, instance mez paměti a restart. Javovská část běhu pokračuje od místa přerušení.

\textbf{Selhání sady (31.~10.).} JVM sady měl výchozí haldu 1,5~GB a Hibernate ji vyčerpal při startu nad řetězem 1~024 negací. \texttt{OutOfMemoryError} ukončil jen vlákno, kterým klient HTTP obsluhuje spojení. Každý další bod se zapsal jako \emph{crashed} s~důvodem \emph{instance neodpověděla}, ačkoli instance odpovídala za šest milisekund. Za 5,6~hodiny tak vzniklo 91 chybných řádků. Záznam přisoudil selhání měřené instanci, chyba byla v~sadě. Řádky nebyly nápadné, protože \emph{crashed} je v~běhu mezí očekávaný výsledek.

\textbf{Opatření.} JVM sady se při vyčerpané haldě ukončí a skript běhu spustí řadu v~novém kontejneru od místa přerušení. Bod, na kterém JVM skončil, je \emph{crashed} oprávněně. Takových restartů proběhlo sedmnáct, každý tam, kde Hibernate nad dlouhým řetězem negací vyčerpá 4,5~GB haldy. Je to mez cíle (tab.~\ref{tab:vrstvy}), ne selhání běhu.

\subsubsection*{Validita běhu SEDE}

\textbf{Chyba metodiky.} Rozhodnutí o~katalozích uváděla, že zdrojem je Dapper, a každá sada to vyložila jinak. Javovská sada spouštěla text katalogu přes JDBC bez překladu. .NET sada spouštěla artefakt, který nástroj vyrobil překladem z~Dapperu do Dapperu. U~FIBEN, jehož dotazy do T-SQL převedl import projektu, dalo obojí tytéž řádky ve všech 237 dotazech.

\textbf{Dopad na ověřené dotazy} (druhý běh, 8.--9.~11., strom \texttt{1f7430b}). Katalog má 1~714 dotazů a 1~201 z~nich nárokuje jako přeložené ve znění autora. Zdroj .NET sady neproběhl u~279 z~1~201: 205 odmítl parser, 73 selhalo nad daty a 1 se nezkompiloval. Text přes JDBC spustil 1~200 z~1~201. Každý z~279 dotazů pak v~.NET sadě selhal ve všech třech cílech jako \emph{zdroj neproběhl}, ačkoli selhal překlad. Kandidát na validátor z~.NET sady proto pokrývá jen 922 nárokovaných dotazů, tedy dotazy, které nástroj přeložil. Dapper se navíc srovnával sám se sebou, takže ztrátu v~jeho překladu žádné srovnání neodhalí. Kandidáti obou sad se na 922 společných dotazech shodují v~907 a liší v~15. Rozbor těchto 15 dotazů hotový není. Shoda se zdrojem v~tab.~\ref{tab:sede} se proto u~cílů EF~Core a NHibernate vztahuje k~překladu nástroje do Dapperu, ne k~textu katalogu, a referenční validátor SEDE neexistuje.

\textbf{Opatření.} Podle rozh.~209 dává řádky, se kterými se cíle srovnávají, v~obou sadách text katalogu spuštěný beze změny a Dapper je v~.NET sadě cílem jako ostatní. Běh podle rozh.~209 zaznamenán není a validátor vznikne až z~něj.

\textbf{První běh (8.~11.)} skončil bez výsledku. Obě části držely řádky každé buňky v~paměti do konce a soubory zapisovaly po posledním dotazu. .NET část skončila po 8~h 16~min chybou alokace při zápisu výsledkového souboru. Od té doby se buňka porovná hned po přečtení a soubory se zapisují po dotazech. Rozhodnutí: 208, 209.

\subsection{Odchylky od specifikace}\label{sec:odchylky}

Tab.~\ref{tab:odchylky} uvádí, kde se výsledek liší od specifikace, a důvod každé odchylky.

{\footnotesize
\begin{longtable}{P{0.30} P{0.30} P{0.40}}
\toprule
\textbf{Specifikace} & \textbf{Výsledek} & \textbf{Důvod} \\
\midrule
\endfirsthead
\toprule
\textbf{Specifikace} & \textbf{Výsledek} & \textbf{Důvod} \\
\midrule
\endhead
\rule{0pt}{26pt}\\
\endfoot
\bottomrule
\caption{Odchylky od specifikace.}\label{tab:odchylky}\\
\endlastfoot
zobrazení mezireprezentace (F14) & nezobrazuje se & tabulka faktů a záznamy uvádějí totéž v~pojmech vstupu a výstupu (rozh.~154) \\
obrazovka s~pěti očíslovanými kroky & oddíly bez čísel, vstup vedle výstupu (rozh.~137) & číslo v~nadpisu neuvádí nic, co neuvádí pořadí oddílů shora dolů. Scénář S7 má čtyři kroky, mez je pět \\
rozšíření Advisoru podle časových možností (F15, T7) & nerealizováno: měří dál jen Dapper a EF~Core & byl zařazen až za experimenty (rozh.~171) a v~čase projektu se nerealizoval \\
T1 & splněn je vstup experimentu, ne výsledek nad ním & případová studie je nad prázdnými tabulkami jen \emph{spustitelná} \\
testovací sady & navíc katalogy LDBC, FIBEN a SEDE & vlastní matice nestačila jako důkaz praktické použitelnosti \\
LLM experimenty mimo podmínky projektu & částečně provedeny (kap.~\ref{sec:llm}) & plné varianty zůstávají diplomové práci \\
\end{longtable}}

\subsection{Využití AI}\label{sec:ai}

Při vývoji řešitel soustavně používal asistenta pro programování (Claude Code od společnosti Anthropic) pod stejným režimem, jaký platí pro každou jinou změnu: nejdřív návrh rozhodnutí, který řešitel přijme, upraví nebo zamítne, potom kód. \textbf{Asistent} psal podle přijatých rozhodnutí kód, testy a dokumentaci, připravoval audity a záznamy běhů a navrhoval znění této zprávy. \textbf{Řešitel} odpovídá za volbu každého rozhodnutí, čte a schvaluje změny, spouští a hodnotí běhy a měření, commituje a dělá release. Ve srovnání s~jazykovými modely (kap.~\ref{sec:llm}) je týž produkt předmětem měření. Obě role jsou oddělené: kanál \texttt{code} volá model bez nástrojů, v~prázdné složce mimo repozitář a s~každým požadavkem zachyceným proxy (rozh.~205).

\section{Experimentální část}\label{sec:exp}

Kapitola odděluje implementovanou infrastrukturu od empiricky ověřených tvrzení. Infrastruktura je vstup experimentu, matice kategorií a metriky (T1--T3). Výsledná čísla předmětem nároku nejsou (rozh.~183, 184): jsou to záznamy jednotlivých běhů v~\texttt{docs/runs/}. Záznam uvádí commit, stroj a prostředí (tab.~\ref{tab:zaznamy-behu}). Záznamy měření S3 ze srpna commit neuvádějí. Běhy proběhly na měřicím stroji (kap.~\ref{sec:behy}). Výjimkou jsou první a třetí měření S3 a pokusy modelu v~běhu z~3.~11., které proběhly na druhém stroji (Intel Core i9-9900, 16~GB).

\textbf{Čísla nejsou naměřena nad odevzdávaným release.} Nad release \texttt{\Vydani} je zaznamenána jen velikost obou testovacích sad. Ostatní čísla pocházejí ze stromů z~21.~8. až 9.~11.~2026 a běh nad opraveným stromem zaznamenán není. Historické jsou proto tři skupiny výsledků, které popisují vady opravené po běhu: řádky výňatku Cofoundry a příkladu \texttt{marketplace-fluent} v~EclipseLinku (tab.~\ref{tab:t2t3}, tab.~\ref{tab:studie}, opraveno rozh.~206 a v~release 2.4.0 a \texttt{\Vydani}), systematická selhání javovských cílů v~běhu SEDE (tab.~\ref{tab:meze}, opraveno rozh.~210--212 a 217) a buňka Cofoundry → Hibernate ve srovnání s~jazykovými modely.

\begin{table}[H]
\centering\footnotesize
\begin{tabular}{l l l}
\toprule
\textbf{Běh} & \textbf{Datum (2026)} & \textbf{Strom} \\
\midrule
výkon překladu (S3) & 21.--26.~8. & záznamy commit neuvádějí \\
podíl F4 & 21.~10. & \texttt{223e791} a následující commit \\
běh mezí & 29.~10.--1.~11. & \texttt{fd27497}, javovská část se čtyřmi soubory z~\texttt{6d5d5a5} \\
celý vstup & 30.~10. & \texttt{6d5d5a5} \\
běh FIBEN & 2.~11. & strom release 2.3.0 (\texttt{63d0b81}) \\
validátor LDBC & 2.--5.~11. & strom release 2.3.0 (\texttt{63d0b81}) \\
experimentální běh & 5.~11. & strom release 2.3.0 (\texttt{a105b3d}) \\
srovnání s~jazykovými modely & 1., 3., 5. a 6.~11. & \texttt{861e228}, \texttt{08b4e4f}, \texttt{a1adad9}, \texttt{5023805} \\
běh SEDE (druhý) & 8.--9.~11. & \texttt{1f7430b} \\
velikost sad & 23.~11. & strom release \Vydani \\
\bottomrule
\end{tabular}
\caption{Záznamy, ze kterých čísla této kapitoly pocházejí. \emph{Strom release} je commit, který se od tagu liší jen dokumentací, číslem verze nebo změnou bez vlivu na výsledek (\texttt{docs/runs/README.md}).}
\label{tab:zaznamy-behu}
\end{table}

\subsection{Metodika}\label{sec:metodika}

Artefakt nástroje i~cizí artefakt se ověřuje ve čtyřech stupních (rozh.~016, 027, 172). Stupně jsou samostatné kontroly a jejich výsledky se uvádějí zvlášť:

\begin{itemize}[leftmargin=*,itemsep=1pt]
\item \textbf{1 přečteno:} parser jazyka cíle přečte text artefaktu bez chyby,
\item \textbf{2 zkompilováno:} artefakt projde překladačem cíle (Roslyn, \texttt{javac}) s~referencemi připnuté verze frameworku,
\item \textbf{3 přijato frameworkem:} framework artefakt přijme, aniž dotaz provede: sestaví z~entit model (EF~Core) nebo factory (NHibernate, Hibernate, EclipseLink, MyBatis) a dotaz připraví,
\item \textbf{4 běží a shoduje se:} dotaz se provede nad databází a vrácené řádky se shodují s~referenčním výsledkem. Tab.~\ref{tab:t2t3} uvádí běh a shodu ve dvou sloupcích.
\end{itemize}

Příklad: metoda dotazu v~LINQ pro EF~Core, která se zkompiluje, kterou EF~Core přijme a která vrátí jiné řádky než referenční výsledek, splňuje stupně 1 až 3 a stupeň 4 ne.

Metodiku dále určují metriky, vstupy, měření času a populace:

\begin{itemize}[leftmargin=*,itemsep=1pt]
\item \textbf{Metriky T3} jsou podíly artefaktů na každém stupni a úplnost přenesených mapovacích faktů: podíl faktů, které parser \emph{téhož} frameworku přečte ze svého artefaktu zpět (rozh.~173, 188).
\item \textbf{Vstupy} jsou pěti druhů (tab.~\ref{tab:vstupy}) a druh určuje referenční výsledek 4.~stupně. .NET sada vstupy vypíše do adresáře běhu. Tentýž výpis dostává jazykový model, takže se srovnává nad týmiž soubory.
\item \textbf{Časy} se měří nad SQL Serverem, který je pro všechny frameworky týž a záměrně omezený (tab.~\ref{tab:rezimy}). Čísla srovnávají artefakty frameworků mezi sebou a mezi verzemi nástroje. Benchmarkem serveru nejsou. Buňka má jedno zahřívací spuštění a pak měřená spuštění, dokud jejich úhrn nedosáhne 1~000~ms, nejméně tři. Zaznamenává se medián (rozh.~165).
\item \textbf{Populace} experimentálního běhu je 624 vstupů, každý do šesti cílů, tedy 3~744 řádků. Matice má 59 kategorií.
\end{itemize}

\begin{table}[H]
\centering\footnotesize
\begin{tabularx}{\textwidth}{L{1.5cm} L{1.2cm} Y L{5.2cm}}
\toprule
\textbf{Druh (klíč)} & \textbf{Vstupů} & \textbf{Co to je} & \textbf{Referenční výsledek 4.~stupně} \\
\midrule
\texttt{matrix:} & 318 & vlastní dotazy matice a kategorie manifestu nad doménou se složenými klíči, z~každého zdroje & kanonický výsledek v~repozitáři \\
\texttt{cross:} & 6 & scénář překladu přes hranici ekosystémů z~každého zdroje, jen s~entitami & -- (3.~stupeň) \\
\texttt{example:} & 22 & příklady včetně šesti výňatků skutečných projektů (Keycloak, Cofoundry, Ambari, Chinook ×~2, JPetStore) & případová studie běží nad prázdnými tabulkami (\emph{spustitelný}), ostatní končí 3.~stupněm \\
\texttt{ldbc:} & 41 & čtecí dotazy LDBC SNB & validační sada LDBC (validátor LDBC, zvlášť) \\
\texttt{fiben:} & 237 & dotazy FIBEN s~otázkou v~přirozeném jazyce & regresní validátor z~doloženého běhu nástroje \\
\bottomrule
\end{tabularx}%
\caption{Vstupy experimentů a jejich počet v~experimentálním běhu za release 2.3.0 (celkem 624). Katalog SEDE (1~714 dotazů) běží jen vlastním během.}
\label{tab:vstupy}
\end{table}

\subsection{Výsledky nástroje}\label{sec:vysledky}

\subsubsection*{Matice T2 × T3}

\textbf{Účel.} Matice odpovídá na otázku požadavků T2 a T3: jaký podíl překladů dosáhne každého stupně ověření v~každé dvojici ekosystémů a jaký podíl mapovacích faktů se přenese. Tab.~\ref{tab:t2t3} je záznam jednoho experimentálního běhu.

\textbf{Sloupce.} \emph{n} je počet překladů. \emph{odm.} je počet překladů bez artefaktu (záznam \emph{Failure}). \emph{nat.\ SQL} je počet dotazů vydaných přes fallback. \emph{přečteno}, \emph{zkomp.} a \emph{přijato} jsou počty artefaktů, které dosáhly stupně 1 až 3, se jmenovatelem \emph{n}. \emph{n4} je počet překladů vstupů se čtvrtým stupněm: matice, FIBEN a případové studie. \emph{běží} má jmenovatel \emph{n4}. \emph{shoda} má jmenovatel jen překlady s~referenčním výsledkem (matice a FIBEN): 1~179, 1~179, 486 a 486. \emph{mapování} je podíl faktů, které parser cíle přečte z~artefaktu zpět, ze všech faktů, které cíl vyjádří, sečtený přes řádky dvojice.

\begin{table}[H]
\centering\footnotesize
\setlength{\tabcolsep}{4pt}
\begin{tabular}{l r r r r r r r r r r}
\toprule
\textbf{Dvojice} & \textbf{n} & \textbf{odm.} & \textbf{nat.\ SQL} & \textbf{přečteno} & \textbf{zkomp.} & \textbf{přijato} & \textbf{n4} & \textbf{běží} & \textbf{shoda} & \textbf{mapování} \\
\midrule
.NET → .NET & 1~341 & 1 & 129 & 1~340 & 1~338 & 1~337 & 1~186 & 1~183 & 1~179 & 99~\% \\
.NET → Java & 1~341 & 2 & 95 & 1~339 & 1~336 & 1~335 & 1~186 & 1~183 & 1~179 & 69~\% \\
Java → .NET & 531 & 4 & 42 & 527 & 527 & 527 & 492 & 489 & 486 & 99~\% \\
Java → Java & 531 & 4 & 29 & 527 & 527 & 527 & 492 & 490 & 486 & 75~\% \\
\bottomrule
\end{tabular}
\caption{Tabulka T2 × T3 po dvojicích ekosystémů, experimentální běh 5.~11.~2026.}
\label{tab:t2t3}
\end{table}

\textbf{Interpretace.} Každý překlad, který má referenční výsledek, se s~ním shoduje: 1~179 z~1~179, 1~179 z~1~179, 486 z~486 a 486 z~486. Tvrzení o~shodě se týká jen těchto 3~330 řádků. Zbylých 414 řádků referenční výsledek nemá:

\begin{itemize}[leftmargin=*,itemsep=1pt]
\item \textbf{388 řádků bez čtvrtého stupně} (\emph{n} $-$ \emph{n4}): dotazy LDBC, jejichž čtvrtým stupněm je validátor LDBC mimo tento běh, příklady a scénáře jen s~entitami, které končí třetím stupněm, a případové studie do Dapperu a MyBatisu kromě výňatku Cofoundry.
\item \textbf{26 řádků případových studií}: 15 běží nad prázdnými tabulkami bez referenčního výsledku (\emph{běží} $-$ \emph{shoda}), 6 patří výňatku Cofoundry, jehož artefakty v~žádném cíli nedosáhly třetího stupně, a 5 je odmítnutých.
\end{itemize}

V~celé populaci je 11 odmítnutí a 7 artefaktů, které nedosáhly třetího stupně:

\begin{itemize}[leftmargin=*,itemsep=1pt]
\item \textbf{11 odmítnutí} je podle rozhodnutí: zdroje Dapper a MyBatis bez katalogu do cílů, které klíč vyžadují (6), případová studie JPetStore do NHibernatu, Hibernatu a EclipseLinku, protože klíč neobsahuje ani její katalog (3), a případové studie JPetStore a Keycloak do EF~Core (2).
\item \textbf{7 artefaktů pod třetím stupněm}: pět se nezkompilovalo a dva odmítl framework. Šest patří výňatku Cofoundry (tab.~\ref{tab:nepovedlo}) a jeden příkladu \texttt{marketplace-fluent} v~EclipseLinku (rozh.~206).
\end{itemize}

Úplnost mapování javovských cílů je nižší (Hibernate 66~\%, EclipseLink 72~\%, MyBatis 80~\%, .NET cíle 99~\%). Způsobují ji fakta, která parser ze svého artefaktu nečte zpět: databázový typ, který JPA určuje z~javovského typu, a \texttt{<id>} MyBatisu. Je to uvedená mez metriky (rozh.~188), ne ztráta v~překladu. Dapper jmenovatele nemá, protože jeho deskriptor nevyjadřuje žádný mapovací fakt.

\subsubsection*{Správnost nad katalogy}

Tab.~\ref{tab:meze} uvádí běhy, které ověřují shodu řádků nad dotazy, jež projekt nepsal. Nezávislý referenční validátor má jen LDBC. Validátor FIBEN vznikl z~běhu nástroje a ověřuje regresi mezi verzemi. SEDE validátor nemá.

\begin{table}[H]
\centering\footnotesize
\begin{tabularx}{\textwidth}{L{4.3cm} Y}
\toprule
\textbf{Běh} & \textbf{Výsledek a omezení} \\
\midrule
validátor LDBC: 138~747 řádků validační sady & každé čtení ve shodě s~očekávanými řádky ve všech šesti frameworcích. Oba běhy také odmítly pět mutací dotazu IC~2 v~každém frameworku: v~negativní kontrole musí pozměněný artefakt skončit rozdílem \\
FIBEN: 237 dotazů & v~obou sadách každá buňka ve shodě se zdrojem i~s~validátorem. Přes fallback jde 61 dotazů do EF~Core, 23 do NHibernatu, 41 do EclipseLinku a 14 do Hibernatu \\
SEDE: 1~201 z~1~714 dotazů přeložených ve znění autora & shoda se zdrojem jen nad buňkami, v~nichž proběhl zdroj i~cíl (tab.~\ref{tab:sede}). V~.NET sadě byl zdrojem překlad nástroje do Dapperu, ne text katalogu, a referenční validátor neexistuje (kap.~\ref{sec:nepovedlo}). Běh našel tři systematické vady javovských cílů, opravené rozh.~210--212 a 217 \\
F4: metadata testovací databáze & 462 ze 462 sledovaných údajů obou schémat testovací databáze přečteno správně (\texttt{CatalogMetadataShareTest}, mez je 95~\%) \\
\bottomrule
\end{tabularx}%
\caption{Správnost nad katalogy a čtení metadat.}
\label{tab:meze}
\end{table}

\begin{table}[H]
\centering\footnotesize
\begin{tabular}{l r r r r r}
\toprule
 & \textbf{EF~Core} & \textbf{NHibernate} & \textbf{Hibernate} & \textbf{EclipseLink} & \textbf{MyBatis} \\
\midrule
buněk, v~nichž proběhl zdroj i~cíl & 293 & 1~031 & 833 & 760 & 866 \\
z~nich ve shodě se zdrojem & 286 & 1~013 & 813 & 746 & 800 \\
\bottomrule
\end{tabular}
\caption{Druhý běh SEDE: shoda se zdrojem po cílech. Počítají se jen buňky, v~nichž proběhl zdroj i~cíl. Zdroj proběhl u~1~060 dotazů v~.NET sadě (922 nárokovaných a 138 dalších) a u~1~200 dotazů v~javovské.}
\label{tab:sede}
\end{table}

\subsubsection*{Případové studie (T1)}

\textbf{Původ a výběr.} Případových studií je šest. Jsou to výňatky pěti open-source projektů pod licencí MIT nebo Apache~2.0 (tab.~\ref{tab:studie-puvod}). Výběr určila čtyři kritéria rozh.~184: čísla požadavku T1 (nejméně 15 entit, vztahy, složený klíč a nejméně 20 čtecích dotazů), permisivní licence, zveřejněné schéma pro SQL Server nebo převoditelné do T-SQL a dotazy ve tvaru, který nástroj čte. Čísla T1 splňují Keycloak a Cofoundry, po jedné aplikaci na každé straně hranice ekosystémů. Ostatní čtyři studie čísla nesplňují a doplňují pokrytí zdrojových frameworků. Studie se zdrojem NHibernate chybí.

\textbf{Meze.} Studie je výňatek, ne celá aplikace. Výňatky jsou upravené: zápisy jsou vynechány a některé konstrukce přepsány, seznam úprav je u~každého výňatku v~repozitáři. Reprezentativnost výběru vůči aplikacím v~praxi nebyla ověřována. Projekty nezveřejňují data, proto čtvrtý stupeň u~studie znamená, že každý dotaz proběhl nad prázdnými tabulkami, bez porovnání řádků. Cíle Dapper a MyBatis končí třetím stupněm, protože jejich SQL neobsahuje schéma.

\begin{table}[H]
\centering\footnotesize
\begin{tabularx}{\textwidth}{L{2.6cm} L{5.4cm} L{2cm} Y}
\toprule
\textbf{Studie} & \textbf{Projekt a licence} & \textbf{Zdroj} & \textbf{Rozsah výňatku} \\
\midrule
Keycloak & \path{keycloak/keycloak}, Apache~2.0 & Hibernate & 22 entit, 11 složených klíčů, 19 vztahů, 47 dotazů \\
Cofoundry & \path{cofoundry-cms/cofoundry}, MIT & EF~Core & 29 entit, 39 jednotek dotazů \\
Ambari & \path{apache/ambari}, Apache~2.0 & EclipseLink & 6 entit, 5 vztahů, 39 dotazů \\
Chinook & \path{cwoodruff/ChinookASPNETWebAPI}, MIT & Dapper & 11 entit, 38 dotazů \\
Chinook & tentýž projekt & EF~Core & 11 entit, 1 složený klíč, 11 vztahů, 35 dotazů \\
JPetStore & \path{mybatis/jpetstore-6}, Apache~2.0 & MyBatis & 7 entit, 14 dotazů \\
\bottomrule
\end{tabularx}%
\caption{Případové studie: původ a rozsah výňatků.}
\label{tab:studie-puvod}
\end{table}

Tab.~\ref{tab:studie} uvádí stupeň (kap.~\ref{sec:metodika}), kterého projekt studie dosáhl jako celek v~každém cíli.

\begin{table}[H]
\centering\footnotesize
\begin{tabular}{l c c c c c c}
\toprule
\textbf{Studie (zdroj)} & \textbf{Dapper} & \textbf{EF~Core} & \textbf{NHibernate} & \textbf{Hibernate} & \textbf{EclipseLink} & \textbf{MyBatis} \\
\midrule
Keycloak (Hibernate) & 3 & odmítnuto & 4 & 4 & 4 & 3 \\
Cofoundry (EF~Core) & 1 & 1 & 2 & 1 & 1 & 1 \\
Ambari (EclipseLink) & 3 & 4 & 4 & 4 & 4 & 3 \\
Chinook (Dapper) & 3 & 4 & 4 & 4 & 4 & 3 \\
Chinook (EF~Core) & 3 & 4 & 4 & 4 & 4 & 3 \\
JPetStore (MyBatis) & 3 & odmítnuto & odmítnuto & odmítnuto & odmítnuto & 3 \\
\bottomrule
\end{tabular}
\caption{Případové studie: dosažený stupeň projektu v~každém cíli (5.~11.~2026). 1 přečteno, 2 zkompilováno, 3 přijato frameworkem, 4 běží nad prázdnými tabulkami.}
\label{tab:studie}
\end{table}

\subsubsection*{Rozsah a výkon překladu}

Tab.~\ref{tab:vykon} uvádí měření času překladu. Doby dlouhých běhů jsou srovnatelné jen v~témže režimu serveru. Obě přehrání validační sady LDBC (29,6~h pro .NET a 27,2~h pro Javu) proběhla ve validačním režimu nad týmž stromem. Běhy FIBEN a SEDE proběhly v~měřicím režimu. U~běhu SEDE nebyl zaznamenán stav procesoru, jeho mediány proto nepocházejí z~doloženého měřicího prostředí.

\begin{table}[H]
\centering\footnotesize
\begin{tabularx}{\textwidth}{L{4.3cm} Y}
\toprule
\textbf{Běh} & \textbf{Výsledek} \\
\midrule
běh mezí: 8~963 bodů .NET části v~patnácti osách a 4~776 řádků javovské & překladač zpracoval řadu po její poslední bod ve všech osách kromě tří. Mez S3 (30~s) překročil řetěz 32 CTE a 512 joinů ze zdroje EF~Core a projekt mezi 1~600 a 6~400 entitami. Meze ostatních vrstev uvádí tab.~\ref{tab:vrstvy} \\
celý vstup: 604--607 tehdejších vstupů v~jednom projektu na framework & každý z~36 překladů pod třetinou meze S3 \\
S3: projekt o~100 entitách a 100 dotazech & 0,2~s na stroji s~Intel Core i9-9900 (21.~8.), 0,4--0,7~s na měřicím stroji (24.~8.), s~katalogem 0,4--0,5~s na stroji s~Intel Core i9-9900, z~toho čtení katalogu 210--330~ms (26.~8.) \\
\bottomrule
\end{tabularx}%
\caption{Rozsah a výkon překladu.}
\label{tab:vykon}
\end{table}

\subsection{Srovnání s~jazykovými modely}\label{sec:llm}

Protokol srovnání (rozh.~175, 189, 202--205, 207) staví na témže výpisu vstupů a na týchž stupních ověření jako experimentální běh. Model je \emph{producent cizího artefaktu} (rozh.~172): jeho odpověď se importuje jako adresář souborů a ověří se stejně jako artefakt nástroje. Kroky uvádí tab.~\ref{tab:protokol}, příkazy výpis~\ref{lst:llm-min}. Protokol byl proveden ve čtyřech zaznamenaných bězích (tab.~\ref{tab:llm-behy}), každá buňka ve třech pokusech. Srovnání je pilotní a jeho výsledky platí jen pro uvedené modely, buňky a varianty.

{\footnotesize
\begin{longtable}{P{0.14} P{0.86}}
\toprule
\textbf{Krok} & \textbf{Co vznikne} \\
\midrule
\endfirsthead
\toprule
\textbf{Krok} & \textbf{Co vznikne} \\
\midrule
\endhead
\rule{0pt}{26pt}\\
\endfoot
\bottomrule
\caption{Kroky protokolu srovnání s~jazykovým modelem.}\label{tab:protokol}\\
\endlastfoot
1.~Buňky & \texttt{cells.tsv}: vstup, zdroj, cíl a varianta \texttt{zero-shot}, \texttt{few-shot} (tři ručně přeložené příklady) nebo \texttt{schema} (s~DDL katalogu). Ve výchozím nastavení tři pokusy na buňku \\
2.~Prompty & \texttt{prompt.md} na pokus. Ukládá modelu pravidla: bez webu, bez nástrojů, odpověď jako celé soubory a závěrečná zpráva \\
3.~Pokusy & \texttt{reply.md} jedním ze čtyř kanálů: \texttt{console} nebo \texttt{chat} (obsluha ručně ve webovém rozhraní, nová konverzace na pokus), \texttt{api} (endpoint kompatibilní s~OpenAI) nebo \texttt{code} (\texttt{claude -p} bez nástrojů, každý požadavek zachytí proxy, rozh.~205) \\
4.~Import & adresář artefaktu z~každé odpovědi. Odpověď bez artefaktu je výsledek modelu na stupni 0. U~variant mimo \texttt{zero-shot} a \texttt{schema} import opraví boilerplate (chybějící \texttt{import} nebo \texttt{using}, třída okolo metody) a zapíše to (rozh.~204, 207) \\
5.~Nástroj & experimentální běh nad týmiž vstupy \\
6.~Validace & \texttt{foreign-run.sh}: proces na artefakt s~mezí času, tytéž stupně ověření \\
7.~Souhrn & oddíl \emph{Foreign artifacts beside the tool over the same cells} v~\texttt{summary.md}: podíl nástroje, podíl po pokusech a best-of-n \\
\end{longtable}}

\textbf{Výběr vzorku.} Buňky vybral řešitel, nejsou náhodným výběrem. Prvních třináct buněk tvoří 11 buněk nad osmi kategoriemi matice, 1 scénář jen s~entitami a 1 příklad: filtrace v~každém ze čtyř směrů mezi ekosystémy a kategorie, které se liší jazykem zdroje nebo cíle (LINQ, HQL, JPQL, mapper XML, kolekční parametr, kategorie mimo jazyk cíle). Případové studie a dotazy FIBEN a LDBC v~něm nejsou, protože se nevejdou do kontextu pilotního modelu nebo potřebují databázi validátoru. Čtvrtý běh má 43 buněk ve třech sadách: \emph{základní} (12 buněk, kategorie matice, které první vzorek neměl), \emph{velká} (14 buněk: sedm buněk s~případovou studií jako celým projektem, z~toho Keycloak ve dvou směrech, dlouhý dotaz, čtyři dotazy FIBEN a dva LDBC) a \emph{záludná} (17 buněk: vstupy, na nichž dříve chyboval nástroj a k~nimž po opravě existuje test).

\textbf{Co model dostal.} Ve variantě \texttt{zero-shot} zadání úlohy, cílový framework a jeho připnutou verzi, tvar odpovědi a tytéž vstupní soubory jako nástroj. Varianta \texttt{few-shot} přidává tři ručně přeložené příklady v~témž směru a import u~ní doplní chybějící \texttt{import} nebo \texttt{using}. Model nedostal schéma databáze, které nástroj čte z~katalogu, dokumentaci frameworku ani nástroje. Variantu \texttt{schema} žádný běh nepoužil.

\textbf{Možná zkreslení.} Ve prospěch nástroje působí dvě okolnosti. Sada \emph{záludná} obsahuje vstupy, pro které byl nástroj opraven. Validátor v~těchto bězích znal prostředí artefaktu (schéma tabulek, reference projektu, jméno assembly) lépe, než je popisoval prompt. Rozh.~207 to pro další běhy odstraňuje. Ve prospěch modelu působí, že v~běhu z~1.~11. obsahoval prompt komentář se jmény tabulek, která nástroj bere z~katalogu. Kategorie matice jsou ve veřejném repozitáři od září 2026. Záznamy uvádějí u~modelů Fable 5.1 a Opus 5.5 hranici trénovacích dat v~červnu 2026, u~modelu Qwen3.5-9B ji výrobce neuvádí. Srovnání hodnotí správnost. Efektivitu běhu přeloženého kódu, kterou u~překladu kódu jazykovými modely hodnotí benchmark TRACE~\cite{trace}, neměří.

\textbf{Pilotní model} je Qwen3.5-9B~\cite{qwen} v~kvantizaci Q4\_K\_M v~serveru llama.cpp~\cite{llamacpp}. Větší modely téže řady se do 12~GB virtuálního stroje nevejdou (rozh.~202). Nasazení uvádí výpis~\ref{lst:llm-server}. Není to silný model ve smyslu požadavku T4: je dolní mezí a ověřením kanálu \texttt{api}. Na čtyřech jádrech trval pokus 5 až 62 minut s~mediánem 16 minut. Běh z~5.~11. volá týmž kanálem \texttt{api} model Claude Fable 5.1 na endpointu výrobce. Jeho soubor \texttt{api.json} se od nastavení pilotního modelu liší adresou, jménem modelu, limitem tokenů, cenami a klíčem. Endpoint výrobce teplotu nepřijímá a seed ignoruje, takže tři pokusy na buňku nejsou bajtově shodné (rozh.~203).

\begin{table}[H]
\centering\footnotesize
\begin{tabularx}{\textwidth}{L{1.8cm} L{3.6cm} L{1.5cm} Y}
\toprule
\textbf{Záznam} & \textbf{Model} & \textbf{Kanál} & \textbf{Vzorek, varianty, nastavení} \\
\midrule
1.~11. & Qwen3.5-9B, Q4\_K\_M (otevřené váhy, lokálně v~llama.cpp na CPU) & \texttt{api} & 13 buněk, \texttt{zero-shot}, teplota 0, seed 1, limit 8~192 tokenů, 39 pokusů, 17,1~h strojového času, 0~USD \\
3.~11. & Claude Fable 5.1 & \texttt{code} & 13 buněk, \texttt{zero-shot} a \texttt{few-shot}, 78 pokusů, teplota nenastavitelná, effort \texttt{high}, předplatné, cena podle ceníku 19,61~USD \\
5.~11. & Claude Fable 5.1 & \texttt{api} & tytéž buňky a prompty jako 3.~11., 78 pokusů, 61,5~min, cena podle ceníku 21,68~USD \\
6.~11. & Claude Opus 5.5 & \texttt{code} & 43 buněk ve třech sadách, obě varianty, 258 pokusů, dvě pětihodinová okna předplatného, cena podle ceníku 111,53~USD \\
\bottomrule
\end{tabularx}%
\caption{Zaznamenané běhy srovnání s~jazykovým modelem.}
\label{tab:llm-behy}
\end{table}

\textbf{Metriky.} Tab.~\ref{tab:llm-vysledky} uvádí počty pokusů, ne buněk: každá buňka má tři pokusy. \emph{přečteno}, \emph{zkompilováno} a \emph{přijato} mají jmenovatel všechny pokusy. \emph{shoda} má jmenovatel pokusy buněk s~referenčním výsledkem: 11 buněk ze 13 a 27 ze 43. \emph{best-of-n} počítá buňky, v~nichž se s~referenčním výsledkem shodl aspoň jeden ze tří pokusů. Je to horní odhad pro uživatele, který pokus opakuje a výsledek ověří, ne podíl úloh vyřešených jedním pokusem.

\begin{table}[H]
\centering\footnotesize
\setlength{\tabcolsep}{3pt}
\begin{tabular}{l l r r r r r}
\toprule
\textbf{Model, kanál} & \textbf{Varianta} & \textbf{přečteno} & \textbf{zkompilováno} & \textbf{přijato} & \textbf{shoda} & \textbf{best-of-n} \\
\midrule
Qwen3.5-9B, api & zero-shot & 3/39 (8~\%) & 3/39 (8~\%) & 0/39 (0~\%) & 0/33 (0~\%) & 0/11 (0~\%) \\
Fable 5.1, code & zero-shot & 39/39 (100~\%) & 39/39 (100~\%) & 34/39 (87~\%) & 28/33 (85~\%) & 10/11 (91~\%) \\
Fable 5.1, code & few-shot & 39/39 (100~\%) & 39/39 (100~\%) & 33/39 (85~\%) & 27/33 (82~\%) & 9/11 (82~\%) \\
Fable 5.1, api & zero-shot & 39/39 (100~\%) & 39/39 (100~\%) & 34/39 (87~\%) & 28/33 (85~\%) & 10/11 (91~\%) \\
Fable 5.1, api & few-shot & 39/39 (100~\%) & 39/39 (100~\%) & 34/39 (87~\%) & 28/33 (85~\%) & 10/11 (91~\%) \\
Opus 5.5, code & zero-shot & 126/129 (98~\%) & 115/129 (89~\%) & 94/129 (73~\%) & 62/81 (77~\%) & 24/27 (89~\%) \\
Opus 5.5, code & few-shot & 126/129 (98~\%) & 120/129 (93~\%) & 102/129 (79~\%) & 65/81 (80~\%) & 23/27 (85~\%) \\
\bottomrule
\end{tabular}
\caption{Pokusy jazykových modelů po stupních ověření: počet a podíl. Sloupec \emph{best-of-n} počítá buňky.}
\label{tab:llm-vysledky}
\end{table}

Nástroj dosáhl nad týmiž buňkami nejvyššího stupně, který buňka umožňuje, ve všech 13 buňkách prvních tří běhů (shoda 11 z~11) a ve 40 ze 43 buněk čtvrtého běhu (shoda 27 z~27). Dvě buňky odmítl podle rozhodnutí: JPetStore → Hibernate bez klíče a Keycloak → EF~Core s~kolekčním parametrem v~nativním SQL. V~buňce Cofoundry → Hibernate selhal dotaz, protože Hibernate tehdy psal tabulku \texttt{User} do SQL bez oddělovačů. Vadu opravuje release \texttt{\Vydani}. Ze čtyř běhů plynou čtyři zjištění:

\begin{itemize}[leftmargin=*,itemsep=1pt]
\item \textbf{Qwen3.5-9B na 13 buňkách dodržel formát odpovědi, ne boilerplate.} Všech 39 odpovědí obsahuje bloky šablony, žádný artefakt není přijatý. Ruční oprava mimo záznam dovedla 8 ze 13 buněk na nejvyšší stupeň (tři jen opravou boilerplate, pět až po opravách podle hlášení frameworku nebo databáze), 5 zůstalo na chybách obsahu. Záznam tedy měří dodržení předepsaného tvaru, ne obsah odpovědi. Dvanáct z~39 pokusů (4 buňky) skončilo na limitu 8~192 tokenů: soubory jsou úplné, závěrečná zpráva opakuje jednu větu až do limitu.
\item \textbf{Fable 5.1 na 13 buňkách: neúspěchy mají příčinu v~promptu.} Každý ze 156 artefaktů obou kanálů prošel čtením i~kompilací. Jedinou příčinou neúspěchu je jméno assembly v~mapování NHibernatu. Prompt ho neuváděl, validátor ho dodával sám (rozh.~028) a mapování s~jiným jménem odmítl. Je to selhání promptu, ne překladu. Od rozh.~207 validátor jméno dosazuje a přepis zaznamená. Vzorek neobsahuje buňku, v~níž by tento model chyboval v~obsahu překladu.
\item \textbf{Opus 5.5 na 43 buňkách: neúspěchy mají dvě skupiny příčin.} Selhání protokolu: žádný pokus žádné ze sedmi buněk s~případovou studií nedošel ke čtvrtému stupni. V~pěti z~nich prompt neobsahoval schéma projektu (25 pokusů s~chybou \emph{Invalid object name}), další pokusy skončily na vzorových hodnotách parametrů, které validátor nemá. Dvanáct pokusů se nezkompilovalo kvůli referencím validátoru. Schéma a reference doplňuje do protokolu rozh.~207. Selhání překladu: \texttt{CASE} v~\texttt{ELSE} do EclipseLinku (všech šest pokusů), cesta, kterou entity neobsahují (dotaz FIBEN \texttt{q32} do Hibernatu, šest pokusů), a dotaz v~souboru \texttt{.sql} s~parametrem HQL (\texttt{order-repository} do NHibernatu, šest pokusů).
\item \textbf{Cena.} 258 pokusů Opus 5.5 stojí podle ceníku 111,53~USD a pokus nad Cofoundry trvá 10--12 minut. Nástroj přeloží projekt o~100 entitách a 100 dotazech pod sekundu (tab.~\ref{tab:vykon}).
\end{itemize}

Varianty \texttt{rag}, \texttt{repair} a \texttt{agent} protokol jmenuje, ale driver je neimplementuje. S~variantou \texttt{schema} patří navazující diplomové práci.

\subsection{Reprodukce experimentů}\label{sec:reprodukce}

Reprodukce má tři úrovně: minimální ověření, reprodukci tabulky T2 × T3 a dlouhé běhy (tab.~\ref{tab:reprodukce}). Společné předpoklady: virtuální stroj Dockeru se čtyřmi procesory a 12~GB bez swapu, asi 30~GB místa, síť pro první build a čistý klon repozitáře (kap.~\ref{sec:spusteni}). Příkazy se spouštějí z~\texttt{ORMConvertor/}. Řešení chyb uvádí tab.~\ref{tab:problemy}.

{\footnotesize
\begin{longtable}{P{0.25} P{0.25} P{0.14} P{0.36}}
\toprule
\textbf{Krok} & \textbf{Prostředí a vstupy} & \textbf{Doba} & \textbf{Kritérium úspěchu} \\
\midrule
\endfirsthead
\toprule
\textbf{Krok} & \textbf{Prostředí a vstupy} & \textbf{Doba} & \textbf{Kritérium úspěchu} \\
\midrule
\endhead
\rule{0pt}{26pt}\\
\endfoot
\bottomrule
\caption{Reprodukce ve třech úrovních. Doby jsou z~měřicího stroje (kap.~\ref{sec:behy}).}\label{tab:reprodukce}\\
\endlastfoot
\multicolumn{4}{@{}l}{\textbf{Minimální ověření}} \\
1.~Systém: \texttt{docker compose up -{}-build} & datové sady z~tab.~\ref{tab:sady} & orientačně desítky minut & první převod z~kap.~\ref{sec:tutorial} vrátí stav \emph{catalog read} a 27 záznamů \emph{Supplied} \\
2.~Obě sady: \texttt{build} a \texttt{run} (kap.~\ref{sec:testy}) & zastavený systém, archiv LDBC scale factoru 0,1 & záznam ji neuvádí & žádný selhaný test a počet testů shodný se záznamem pro tag \texttt{\Vydani} (tab.~\ref{tab:cisla}) \\
\midrule
\multicolumn{4}{@{}l}{\textbf{Reprodukce tabulky T2 × T3 a srovnání s~modelem}} \\
3.~Experimentální běh bez druhu \texttt{fiben} (výpis~\ref{lst:exp-beh}) & profil \texttt{test} & záznam ji neuvádí & \texttt{summary.md} ve volume s~tabulkami ve sloupcích tab.~\ref{tab:t2t3} a 2~322 řádky \\
4.~Experimentální běh s~druhem \texttt{fiben} (výpis~\ref{lst:exp-fiben}) & navíc databáze \texttt{Fiben} v~\texttt{mssql\_db} a instance s~katalogem nad ní & 58~min včetně buildu & \texttt{summary.md} s~celou populací 3~744 řádků \\
5.~Srovnání s~LLM: jedna buňka, tři pokusy, kanál \texttt{api} s~lokálním modelem & model zabere asi 6~GB paměti & pokus buňky z~výpisu~\ref{lst:llm-min} asi 12~minut & tři bajtově shodné odpovědi a oddíl o~cizích artefaktech v~\texttt{summary.md} \\
\midrule
\multicolumn{4}{@{}l}{\textbf{Dlouhé běhy}} \\
6.~Běh FIBEN (\emph{The FIBEN run} v~README): .NET část a po ní javovská & měřicí režim serveru & 2,3~h a 3,1~h & \texttt{summary.txt}: shoda se zdrojem a s~validátorem stejná jako v~záznamu. Mediány se liší podle stroje \\
7.~Validátor LDBC (\emph{The LDBC judge} v~README) & validační režim, \texttt{judge.env} na každém příkazu compose & 29,6~h a 27,2~h & souhrn testů bez selhání \\
\end{longtable}}

\textbf{Experimentální běh} (T2 × T3) má tři kroky (obr.~\ref{fig:beh}). .NET část vypíše vstupy, přeloží každý do šesti cílů, tři .NET cíle ověří ve všech stupních a změří úplnost přenesených mapovacích faktů. Javovská část ověří tři javovské cíle. Souhrn spočítá podíly po dvojicích ekosystémů, směrech, kategoriích a druzích vstupu. Odmítnutí i~nesplněný stupeň se zapíší jako řádky a běh neukončí. Běh ukončí jen jeho vlastní vada: vstup bez manifestu nebo vstup, pro který běh nemá referenční výsledek. Druh \texttt{fiben} běh bez nastavené databáze \texttt{Fiben} vynechá bez chyby. Běh se spouští v~kontejnerech (výpis~\ref{lst:exp-beh}), nebo na hostiteli s~oběma toolchainy a testovací databází (výpis~\ref{lst:exp-host}).

\begin{figure}[H]
\centering
\begin{tikzpicture}[node distance=12mm and 8mm]
\node[box, text width=30mm] (inp) {výpis vstupů\\\texttt{inputs/}};
\node[hl, text width=30mm, right=of inp] (net) {.NET část\\\texttt{tool-rows.tsv}\\\texttt{mapping-rows.tsv}};
\node[hl, text width=30mm, right=of net] (jav) {javovská část\\\texttt{java-tool-rows.tsv}};
\node[box, text width=30mm, right=of jav] (sum) {souhrn\\\texttt{summary.md}};
\node[box, text width=30mm, below=of inp] (prm) {buňky a prompty\\\texttt{cells.tsv}, \texttt{prompt.md}};
\node[hl, text width=30mm, right=of prm] (att) {pokusy modelu\\kanál console, chat,\\api nebo code\\\texttt{reply.md}};
\node[box, text width=30mm, right=of att] (imp) {import\\\texttt{artifact.json}};
\node[hl, text width=30mm, right=of imp] (jud) {validace cizích artefaktů\\\texttt{foreign-run.sh}};
\draw[arr] (inp) -- (net);
\draw[arr] (net) -- (jav);
\draw[arr] (jav) -- (sum);
\draw[arr] (inp) -- (prm);
\draw[arr] (prm) -- (att);
\draw[arr] (att) -- (imp);
\draw[arr] (imp) -- (jud);
\draw[arr] (jud) -- (sum);
\node[lbl, anchor=south west] at (inp.north west) {nástroj: experimentální běh};
\node[lbl, anchor=north west] at (prm.south west) {jazykový model: protokol srovnání};
\end{tikzpicture}
\caption{Experimentální běh a srovnání s~jazykovým modelem nad týmž výpisem vstupů a s~týmž souhrnem.}
\label{fig:beh}
\end{figure}

\noindent\begin{minipage}{\linewidth}
\begin{lstlisting}[language=bash,caption={Experimentální běh na hostiteli.},label={lst:exp-host}]
export ORMCONVERTOR_EXPERIMENT_RESULTS=<adresář běhu>
dotnet test Tests/Tests.csproj --configuration Release --filter "FullyQualifiedName~Tests.Experiment.ExperimentRunTest"
(cd JavaTests && mvn -B test -Dtest=ExperimentRunTest)   # ORMCONVERTOR_API_URL: instance nad testovací databází
dotnet test Tests/Tests.csproj --configuration Release --filter "FullyQualifiedName~ExperimentSummaryTest"
\end{lstlisting}
\end{minipage}

Na hostiteli potřebují vstupy druhu \texttt{fiben} klíč \path{ConnectionStrings:FibenDatabase} a pro Javu nastavení běhu FIBEN s~instancí, jejíž katalog je nad databází \texttt{Fiben}. Vstupy druhu \texttt{ldbc} potřebují pro Javu instanci s~katalogem nad \texttt{LdbcSnb} (\path{ORMCONVERTOR_LDBC_API_URL}). Databázi \texttt{ORMConvertorCaseStudies} a pro Javu instanci nad ní (\path{ORMCONVERTOR_CASESTUDY_API_URL}) potřebují případové studie.

\noindent\begin{minipage}{\linewidth}
\begin{lstlisting}[language=bash,caption={Experimentální běh v~kontejnerech bez druhu \texttt{fiben}: 2~322 řádků.},label={lst:exp-beh}]
docker volume create exp
docker compose --profile test run --rm -v exp:/results -e ORMCONVERTOR_EXPERIMENT_RESULTS=/results \
  -e ORMCONVERTOR_EXPERIMENT_KINDS=matrix,cross,example,ldbc \
  tests --filter "FullyQualifiedName~Tests.Experiment.ExperimentRunTest"
docker compose --profile test run --rm -v exp:/results -e ORMCONVERTOR_EXPERIMENT_RESULTS=/results \
  -e ORMCONVERTOR_EXPERIMENT_KINDS=matrix,cross,example,ldbc java_tests -Dtest=ExperimentRunTest
docker compose --profile test run --rm -v exp:/results -e ORMCONVERTOR_EXPERIMENT_RESULTS=/results \
  tests --filter "FullyQualifiedName~ExperimentSummaryTest"
\end{lstlisting}
\end{minipage}

Výpis~\ref{lst:exp-beh} vynechává druh \texttt{fiben}, protože profil \texttt{test} databázi \texttt{Fiben} neobsahuje. Dává proto (624 $-$ 237) × 6 = 2~322 řádků. Publikovanou populaci 3~744 řádků reprodukují tytéž tři příkazy bez proměnné \path{ORMCONVERTOR_EXPERIMENT_KINDS}, doplněné o~nastavení z~výpisu~\ref{lst:exp-fiben} (\path{ORMConvertor/README.md}, \emph{The experiment run} a \emph{The FIBEN run}). Zaznamenaný běh měl oba databázové servery omezené na 2~048~MB paměti (\path{DATABASE_MAX_SERVER_MEMORY_MB=2048} u~příkazů compose), aby se systém a profil \texttt{test} vešly do 12~GB.

\noindent\begin{minipage}{\linewidth}
\begin{lstlisting}[language=bash,caption={Doplnění experimentálního běhu o~druh \texttt{fiben}: 3~744 řádků.},label={lst:exp-fiben}]
# databáze Fiben a instance s katalogem nad ní
docker compose up -d --build mssql_db
ORMCONVERTOR_CATALOG_DATABASE=Fiben docker compose up -d --build ormconvertor
# argument navíc pro .NET část (tests)
-e "ConnectionStrings__FibenDatabase=Server=mssql_db,1433;Database=Fiben;User ID=sa;Password=Testingorms123;TrustServerCertificate=true;"
# argumenty navíc pro javovskou část (java_tests)
-e "ORMCONVERTOR_TEST_FIBEN_JDBC_URL=jdbc:sqlserver://mssql_db:1433;databaseName=Fiben;user=sa;password=Testingorms123;encrypt=true;trustServerCertificate=true"
-e ORMCONVERTOR_FIBEN_API_URL=http://ormconvertor:5072/orm
\end{lstlisting}
\end{minipage}

\textbf{Srovnání s~jazykovým modelem} (krok~5 tab.~\ref{tab:reprodukce}) se spouští na hostiteli s~oběma toolchainy. Lokální model pro kanál \texttt{api} nasadí výpis~\ref{lst:llm-server}. Adresář běhu obsahuje soubor \texttt{cells.tsv} s~buňkami a soubor \texttt{api.json} s~adresou endpointu a nastavením modelu (\path{ORMConvertor/README.md}, \emph{A language model over an API}). Výpis~\ref{lst:llm-min} uvádí příkazy v~pořadí kroků tab.~\ref{tab:protokol}. Očekávané soubory jsou \path{attempts/<id>/prompt.md} a \texttt{reply.md} pro každý pokus, \path{artifacts/<id>/artifact.json} pro každou odpověď s~artefaktem, \texttt{foreign-rows.tsv} a \texttt{summary.md}. Běh je dokončen, když \texttt{summary.md} obsahuje oddíl \emph{Foreign artifacts beside the tool over the same cells}.

\noindent\begin{minipage}{\linewidth}
\begin{lstlisting}[language=bash,caption={Nasazení lokálního modelu pro kanál \texttt{api}.},label={lst:llm-server}]
sha256sum <models>/Qwen3.5-9B-Q4_K_M.gguf
docker pull ghcr.io/ggml-org/llama.cpp:server
docker image inspect -f '{{index .RepoDigests 0}}' ghcr.io/ggml-org/llama.cpp:server
MSYS_NO_PATHCONV=1 docker run -d --name llm_server -p 8080:8080 -v "<models>:/models:ro" \
  ghcr.io/ggml-org/llama.cpp:server -m /models/Qwen3.5-9B-Q4_K_M.gguf --host 0.0.0.0 --port 8080 \
  -c 16384 -t 4 --load-mode none --temp 0 --seed 1 --reasoning off
curl -s http://localhost:8080/health          # {"status":"ok"} po načtení, asi minuta
\end{lstlisting}
\end{minipage}

\noindent\begin{minipage}{\linewidth}
\begin{lstlisting}[language=bash,caption={Srovnání s~jazykovým modelem kanálem \texttt{api}: řádek souboru \texttt{cells.tsv} a příkazy.},label={lst:llm-min}]
# cells.tsv (v souboru sloupce odděluje tabulátor)
answers                               source   target   variant     set
matrix:AggregationGroupingAndHaving   Dapper   EFCore   zero-shot   basic

export ORMCONVERTOR_LLM_RUN=<adresář běhu>
export ORMCONVERTOR_EXPERIMENT_RESULTS=<adresář běhu>
dotnet test Tests/Tests.csproj --filter "FullyQualifiedName~WebAttemptTest.ThePromptsAreWrittenWhereAsked"     # krok 2
dotnet test Tests/Tests.csproj --filter "FullyQualifiedName~ApiAttemptTest.TheAttemptsAreMadeWhereAsked"       # krok 3
dotnet test Tests/Tests.csproj --filter "FullyQualifiedName~WebAttemptTest.TheRepliesAreImportedWhereAsked"    # krok 4
# krok 5: příkazy ExperimentRunTest obou sad z výpisu experimentálního běhu na hostiteli
./foreign-run.sh <adresář běhu>/artifacts <adresář běhu>                                                       # krok 6
dotnet test Tests/Tests.csproj --filter "FullyQualifiedName~ExperimentSummaryTest"                             # krok 7
\end{lstlisting}
\end{minipage}

Po skončení se kontejnery běhů odstraní jménem (\path{docker rm -f llm_server fiben_run_dotnet fiben_run_java}) a systém se spustí příkazem \path{docker compose start ormconvertor mssql_db}. \texttt{docker compose down -v} odstraní i~volume s~databázemi a další start načítá vše znovu.

\section{Závěr}\label{sec:zaver}

Projekt splnil, co specifikace označila za hlavní část rozšíření: společná mezireprezentace a překladový přístup fungují i~pro druhý programovací jazyk a druhý technologický ekosystém. Tři javovské frameworky byly přidány jako samostatné wrappery nad sdílenými vrstvami (S1), překlad jde ve všech 36 směrech šesti frameworků. Překládaná podmnožina je vymezená číslovanými rozhodnutími a kategorie matice jsou z~každého zdroje, který je vyjádří, do všech šesti cílů ověřené na čtyřech stupních. Stav požadavků shrnuje tab.~\ref{tab:prehled-pozadavku}, navazující práci tab.~\ref{tab:dalsi}. Specifikace jmenuje tři přínosy a baseline pro srovnání s~jazykovými modely. Praktickou použitelnost hodnotí výsledky nad šesti případovými studiemi a třemi katalogy dotazů (kap.~\ref{sec:vysledky}). Reprodukovatelné experimentální prostředí tvoří kontejnery požadavku S5 a postup v~kap.~\ref{sec:reprodukce}. Zbylé dva body:

\begin{itemize}[leftmargin=*,itemsep=1pt]
\item \textbf{Cross-language použitelnost přístupu.} Táž mezireprezentace a tytéž překladové kroky dávají artefakty, které javovské frameworky přijmou a které vracejí tytéž řádky jako zdroj. V~experimentálním běhu to platí pro každý překlad s~referenčním výsledkem: pro kategorie matice a pro 237 dotazů benchmarku FIBEN, jehož dotazy projekt nepsal (tab.~\ref{tab:t2t3}). Každé čtení validační sady LDBC se ve všech šesti frameworcích shoduje s~jejími očekávanými řádky. Validační sada LDBC je jediný validátor nezávislý na projektu.
\item \textbf{Deterministická baseline a společné metriky.} Artefakt nástroje i~jazykového modelu prochází týmiž stupni ověření. Ve čtyřech zaznamenaných bězích měly neúspěšné pokusy modelů Fable 5.1 a Opus 5.5 dvě skupiny příčin. První je kontext, který prompt neobsahoval: schéma projektu a jméno assembly. Rozh.~207 ho do protokolu doplnilo a jeho vliv změří až další běh. Druhá jsou chyby obsahu překladu, mezi nimi chování připnuté verze frameworku, které má nástroj ověřené testy. Srovnání je pilotní: tři modely, vzorky 13 a 43 buněk vybrané řešitelem, tři pokusy na buňku a varianty \texttt{zero-shot} a \texttt{few-shot} bez schématu databáze. Na jiné modely, vstupy a varianty protokolu se výsledky nezobecňují. Schopnost modelu překládat od úplnosti promptu oddělí až běh s~variantou \texttt{schema}.
\end{itemize}

\begin{table}[H]
\centering\footnotesize
\begin{tabularx}{\textwidth}{L{2.6cm} Y}
\toprule
\textbf{Oblast} & \textbf{Na čem lze pokračovat} \\
\midrule
experimenty & plné experimenty T4--T6 v~navazující diplomové práci (varianta \texttt{schema}, RAG, opravný cyklus, agentní varianta, ablace), běh SEDE s~validátorem (rozh.~209) \\
Advisor & měření všech šesti frameworků (rozh.~146, 147), srovnání s~heuristikou (T7), izolace spouštění cizího kódu (S4) \\
\bottomrule
\end{tabularx}%
\caption{Navazující práce. Úplný seznam je v~\texttt{docs/open-items.md}.}
\label{tab:dalsi}
\end{table}

Projekt mi přinesl zkušenosti z~několika oblastí. Největší přínos má softwarověinženýrská stránka vývoje: vedoucí se v~roli stakeholdera pravidelně dotazoval na stav a během vývoje jsem si vytvořil postup, kterým jsem úkoly zadával nástroji Claude Code a jeho výstupy kontroloval testy a revizí kódu. Díky tomu jsem se mohl více soustředit na rozhodování o~implementačních detailech, a tím nepřímo poznat detaily ORM frameworků. Nesporně k~tomu přispěla i~nepřetržitá podpora vedoucího, který mě držel na správné cestě a ve správném tempu, za což mu tímto děkuji.

\begin{thebibliography}{99}
\addcontentsline{toc}{section}{\refname}

\bibitem{spec}
ŠTOCHEL, Miroslav. \emph{Rozšíření nástroje ORMorpher: Specification of team software project}. Schválená specifikace výzkumného projektu, MFF UK, 2026. V~repozitáři projektu jako \path{docs/specifikace.pdf}.

\bibitem{zamer}
ŠTOCHEL, Miroslav. \emph{Rozšíření nástroje ORMorpher: záměr projektu}. Schválený záměr, MFF UK, 2026. V~repozitáři projektu jako \texttt{docs/zamer.tex}.

\bibitem{fork-repo}
ŠTOCHEL, Miroslav. \emph{orm-convertor} [online]. Repozitář projektu ORMorpher, release \Vydani, 2026 [cit.~\mbox{2026-11-24}]. Dostupné z: \url{https://github.com/MiraStochel/orm-convertor}

\bibitem{abraham-repo}
ABRAHÁM, Milan. \emph{Framework-Agnostic Query Adaptation: Ensuring SQL Compatibility Across .NET Database Frameworks}. Diplomová práce, vedoucí Koupil, Pavel. Praha: Univerzita Karlova, Matematicko-fyzikální fakulta, Katedra softwarového inženýrství, 2025. Repozitář prototypu [online] [cit.~\mbox{2026-11-24}]. Dostupné z: \url{https://github.com/milan252525/orm-convertor}

\bibitem{ormorpher-ase}
ABRAHÁM, Milan; KOUPIL, Pavel. ORMorpher: An Interactive Framework for ORM Translation and Optimization. In: \emph{2025 40th IEEE/ACM International Conference on Automated Software Engineering (ASE)}. IEEE, 2025, s.~4037--4040. Dostupné z: \url{https://doi.org/10.1109/ASE63991.2025.00369}

\bibitem{ormorpher-jss}
ABRAHÁM, Milan; KOUPIL, Pavel. \emph{A Unified Framework for Object-Relational Mapping Translation and Performance-Aware Selection}. Rozšířená časopisecká verze, preprint (under review), 2026.

\bibitem{adam-web}
\emph{Adaptive Data Management Research Group} [online]. CNAIP [cit.~\mbox{2026-11-24}]. Dostupné z: \url{https://www.cnaip.cz/en/entity/adaptive-data-management-research-group}

\bibitem{mmcat-intro}
MM-cat Docs. \emph{Motivation} [online] [cit.~\mbox{2026-11-09}]. Dostupné z: \url{https://mmcatdb.com/getting-started/introduction/}

\bibitem{mmcat-aimm}
MM-cat Docs. \emph{Database Mapping Optimization} [online] [cit.~\mbox{2026-11-09}]. Dostupné z: \url{https://mmcatdb.com/aimm/}

\bibitem{koupil-holubova-2022}
KOUPIL, Pavel; HOLUBOVÁ, Irena. A unified representation and transformation of multi-model data using category theory. \emph{Journal of Big Data}, 2022, roč.~9, č.~1, čl.~61. Dostupné z: \url{https://doi.org/10.1186/s40537-022-00613-3}

\bibitem{koupil-crha-holubova-2025}
KOUPIL, Pavel; CRHA, Daniel; HOLUBOVÁ, Irena. A universal approach for simplified redundancy-aware cross-model querying. \emph{Information Systems}, 2025, roč.~127, čl.~102456. Dostupné z: \url{https://doi.org/10.1016/j.is.2024.102456}

\bibitem{mm-quecat}
KOUPIL, Pavel; CRHA, Daniel; HOLUBOVÁ, Irena. MM-quecat: A Tool for Unified Querying of Multi-Model Data. In: \emph{Proceedings of the 26th International Conference on Extending Database Technology (EDBT 2023)}. OpenProceedings.org, 2023, s.~831--834. Dostupné z: \url{https://doi.org/10.48786/edbt.2023.76}

\bibitem{mm-evoque}
KOUPIL, Pavel; BÁRTÍK, Jáchym; HOLUBOVÁ, Irena. MM-evoque: Query Synchronisation in Multi-Model Databases. In: \emph{Proceedings of the 27th International Conference on Extending Database Technology (EDBT 2024)}. OpenProceedings.org, 2024, s.~818--821. Dostupné z: \url{https://doi.org/10.48786/edbt.2024.78}

\bibitem{mm-evocat}
KOUPIL, Pavel; BÁRTÍK, Jáchym; HOLUBOVÁ, Irena. MM-evocat: A Tool for Modelling and Evolution Management of Multi-Model Data. In: \emph{Proceedings of the 31st ACM International Conference on Information \& Knowledge Management (CIKM~'22)}. ACM, 2022, s.~4892--4896. Dostupné z: \url{https://doi.org/10.1145/3511808.3557180}

\bibitem{mm-mapsearch}
KOUPIL, Pavel et al. \emph{mm-mapsearch} [online]. Stránka nástroje [cit.~\mbox{2026-11-24}]. Dostupné z: \url{https://www.ksi.mff.cuni.cz/~koupil/software/mm-mapsearch/}

\bibitem{corovcak-dp}
ČOROVČÁK, Martin. \emph{Universal Object Mapping Assistant: LLM-Assisted Migration of Application Persistence Layers}. Diplomová práce, vedoucí Koupil, Pavel. Praha: Univerzita Karlova, Matematicko-fyzikální fakulta, Katedra softwarového inženýrství, 2026.

\bibitem{obonova-dp}
OBOŇOVÁ, Ivona. \emph{Query Rewriting Framework for Translating SQL into Graph Database Query Languages}. Diplomová práce, vedoucí Koupil, Pavel. Praha: Univerzita Karlova, Matematicko-fyzikální fakulta, Katedra softwarového inženýrství, 2026.

\bibitem{ldbc}
LDBC. \emph{LDBC Social Network Benchmark (LDBC SNB)}, Interactive workload v1 [online]. Graph Data Council [cit.~\mbox{2026-11-24}]. Dostupné z: \url{https://ldbcouncil.org/benchmarks/snb/}. Data pod licencí LDBC k~datovým sadám. Nic v~projektu není LDBC Benchmark a žádné zde uvedené číslo není jeho výsledkem.

\bibitem{fiben}
IBM Research. \emph{FIBEN: A Finance Dataset Benchmark for Natural Language Queries} [online]. Repozitář \path{IBM/fiben-benchmark}, v~listopadu 2026 nedostupný. Archivovaná kopie z~1.~3.~2026 [cit.~\mbox{2026-11-24}]. Dostupné z: \url{http://web.archive.org/web/20260301045548/https://github.com/IBM/fiben-benchmark}

\bibitem{sede}
HAZOOM, Moshe; MALIK, Vibhor; BOGIN, Ben. Text-to-SQL in the Wild: A Naturally-Occurring Dataset Based on Stack Exchange Data. In: \emph{Proceedings of the 1st Workshop on Natural Language Processing for Programming (NLP4Prog 2021)}. Association for Computational Linguistics, 2021, s.~77--87. Dostupné z: \url{https://doi.org/10.18653/v1/2021.nlp4prog-1.9}. Repozitář sady [online] [cit.~\mbox{2026-11-24}]. Dostupné z: \url{https://github.com/hirupert/sede}. Texty dotazů a data (dump dba.stackexchange.com z~2.~4.~2024) jsou obsah uživatelů Stack Exchange pod CC~\mbox{BY-SA}.

\bibitem{qwen}
Qwen Team. \emph{Qwen3.5-9B} [online]. Hugging Face [cit.~\mbox{2026-11-24}]. Dostupné z: \url{https://huggingface.co/Qwen/Qwen3.5-9B}. Kvantizace Q4\_K\_M z~repozitáře \path{unsloth/Qwen3.5-9B-GGUF}.

\bibitem{llamacpp}
GERGANOV, Georgi et al. \emph{llama.cpp} [online]. Serverový image \path{ghcr.io/ggml-org/llama.cpp:server} [cit.~\mbox{2026-11-24}]. Dostupné z: \url{https://github.com/ggml-org/llama.cpp}

\bibitem{hibernate-docs}
Hibernate Project. \emph{Hibernate ORM Documentation} [online] [cit.~\mbox{2026-11-24}]. Dostupné z: \url{https://hibernate.org/orm/documentation/}

\bibitem{eclipselink-docs}
Eclipse Foundation. \emph{EclipseLink Documentation} [online] [cit.~\mbox{2026-11-24}]. Dostupné z: \url{https://eclipse.dev/eclipselink/documentation/}

\bibitem{mybatis-docs}
MyBatis Project. \emph{MyBatis 3 Documentation} [online] [cit.~\mbox{2026-11-24}]. Dostupné z: \url{https://mybatis.org/mybatis-3/}

\bibitem{trace}
GONG, Zhihao et al. TRACE: Evaluating Execution Efficiency of LLM-Based Code Translation. In: \emph{Proceedings of the 64th Annual Meeting of the Association for Computational Linguistics (Volume 1: Long Papers)}. Association for Computational Linguistics, 2026, s.~3089--3117. Dostupné z: \url{https://doi.org/10.18653/v1/2026.acl-long.140}

\end{thebibliography}

\end{document}
